\chapter{状态估计与最小二乘}
\label{ch:slam_est}

% ════════════════════════════════════════════════════════════════
%  P1 首章：全吸收重建 —— Barfoot SER2 ch1(概率论基元 2.x)+ ch3(线性高斯估计 3.x)
%  + 视觉SLAM十四讲 ch6(§6.1) + SLAM Handbook ch1(Factor Graphs)。
%  自包含独立记录：每步推导/每例/每定理+证明/每表全录，读者无需翻原书。
%  范围：概率/Bayes/高斯/马氏/信息形式/Schur/SMW + MLE/MAP + 高斯→最小二乘
%        + 线性高斯批量正规方程 + 批量vs递归(等价) + 因子图与 MAP 等价。
%  边界：递归滤波 KF/ESKF 的完整递推留 ch:eskf（本章只证它是批量 MAP 的递归特例）；
%        GN/LM/Schur 求解器内核留 ch:nlopt（本章只写"要解什么"）。
%  教学层遵循 v5.0 理论教学规范。右扰动为主（本章状态为向量空间，含位姿处用右扰动）。
% ════════════════════════════════════════════════════════════════

\begin{practice}[前置自测]
开始之前，先试答下列问题。若已能流畅作答，可只速读本章表格与速查卡；若卡住，正是本章要为你解决的。
\begin{enumerate}
  \item 运动方程与观测方程因噪声不可能\emph{精确}成立，那么“求解 SLAM”到底在求什么？
  \item 最大似然估计（MLE）与最大后验估计（MAP）有何区别？什么情形下二者给出同一答案？
  \item 为什么在\emph{高斯噪声}假设下，“最大化概率”会变成“最小化误差平方和”？式子里的\emph{信息矩阵}扮演什么角色？
  \item 协方差形式与信息（逆协方差）形式各擅长表达什么？为什么 SLAM 后端偏爱信息形式？给定联合高斯 $p(\mathbf{x},\mathbf{y})$，\emph{条件化}得 $p(\mathbf{x}\mid\mathbf{y})$、\emph{边缘化}得 $p(\mathbf{x})$，分别要做什么矩阵运算？
  \item “因子图”是什么？它与最小二乘、与你在 GTSAM/g2o 里搭的图，是同一回事吗？“批量平滑”与“递归滤波”谁的信息矩阵稀疏？
\end{enumerate}
\noindent\emph{答案线索}：(1) 概率推断——求最可能状态及其不确定度，非解方程，见 \cref{sec:est-why}；(2) MAP $=$ MLE $+$ 先验，均匀先验下同解，见 \cref{sec:est-mle-map}；(3) 取负对数把高斯指数变成马氏二次型，信息矩阵 $=$ 权重，见 \cref{sec:est-ls}；(4) 见 \cref{tab:cov-info} 与 \cref{sec:est-gauss}；(5) 见 \cref{sec:est-fg}。答错 $\ge 2$ 题，建议精读全章。
\end{practice}

\section*{本章目标}
学完本章，你应当能够：
\begin{enumerate}
  \item \textbf{说清}为何状态估计是\emph{概率推断}而非“解方程”，并写出 SLAM 的运动/观测模型与后验 $p(\mathbf{x}\mid\mathbf{z})$；
  \item \textbf{推导}本章所需的全部概率论与高斯分布工具（贝叶斯、期望/协方差、多元高斯、马氏距离、卡方、\emph{信息形式}、舒尔补条件化、舒尔补边缘化、归一化乘积$=$信息相加、SMW 矩阵求逆引理、CRLB/Fisher 信息），且每条都能独立复现；
  \item \textbf{推导} MLE 与 MAP，解释先验的作用以及二者何时等价；
  \item \textbf{独立推导}高斯噪声下批量 MAP $\Longrightarrow$ 加权最小二乘，理解信息矩阵即权重；
  \item \textbf{写出}线性高斯批量问题的\emph{正规方程} $\mathbf{H}^\top\mathbf{W}^{-1}\mathbf{H}\,\hat{\mathbf{x}}=\mathbf{H}^\top\mathbf{W}^{-1}\mathbf{z}$，并证明“解的协方差 $=$ 逆信息矩阵 $=$ 逆 Hessian $=$ 克拉默–拉奥下界”；
  \item \textbf{证明}线性高斯下 MAP $=$ 贝叶斯后验均值（三条等价路径），并据可观测性判定解的存在唯一性；
  \item \textbf{解释}批量与递归的等价：信息矩阵的\emph{块三对角}稀疏（源于马尔可夫性）使批量可 $O(K)$ 求解，并把卡尔曼滤波识别为批量 MAP 的\emph{前向递归特例}；
  \item \textbf{用因子图}统一表示估计问题，列出 SLAM 各类因子（先验/里程计/观测/回环）的残差，区分平滑（信息矩阵\emph{稀疏}）与滤波（\emph{稠密}），并理解变量消元 $=$ 稀疏矩阵分解；
  \item \textbf{定位}后续：递归滤波（\cref{ch:eskf}）是本章 MAP 的特例，GN/LM 与稀疏求解（\cref{ch:nlopt}）是解正规方程的迭代法与数值内核。
\end{enumerate}

\subsection*{本章知识导航}
本章主线：\emph{把“带噪声的 SLAM 模型”翻译成一个可优化的最小二乘 / 因子图问题，并把它解到底}。结构如下：

\begin{center}
\small
\begin{tikzpicture}[
  node distance=5mm and 7mm, >=Stealth,
  blk/.style={draw, rounded corners, align=center, font=\small, inner sep=3pt, fill=main!6, draw=main!60!black},
  grp/.style={draw, rounded corners, align=center, font=\small, inner sep=3pt, fill=second!6, draw=second!60!black},
  fwd/.style={draw, rounded corners, align=center, font=\small, inner sep=3pt, fill=third!8, draw=third!60!black},
  ln/.style={->, thick, gray!60!black},
  fl/.style={->, thick, dashed, third!60!black}]
\node[blk] (prob) {概率 · Bayes};
\node[blk, right=of prob] (gauss) {高斯 · 马氏 · Schur};
\node[blk, right=of gauss] (mle) {MLE · MAP};
\node[blk, right=of mle] (ls) {最小二乘 $J(\mathbf{x})$};
\node[grp, below=9mm of prob] (lin) {线性高斯\\正规方程};
\node[grp, right=of lin] (info) {信息形式\\稀疏 · 可观性};
\node[grp, right=of info] (br) {批量 vs 递归\\（等价）};
\node[grp, right=of br] (fg) {因子图\\（统一框架）};
\node[fwd, below=9mm of br] (filt) {滤波 KF/ESKF\\（递归特例）};
\node[fwd, below=9mm of fg] (opt) {GN/LM 求解器\\（解正规方程）};
\draw[ln] (prob)--(gauss); \draw[ln] (gauss)--(mle); \draw[ln] (mle)--(ls);
\draw[ln] (ls.south) -- ++(0,-4.5mm) -| (lin.north);
\draw[ln] (lin)--(info); \draw[ln] (info)--(br); \draw[ln] (br)--(fg);
\draw[fl] (br.south)--(filt.north);
\draw[fl] (fg.south)--(opt.north);
\end{tikzpicture}
\end{center}

\noindent\textbf{推荐路径（脉络层）}：\emph{骨干}是第 1--4 节（从概率到最小二乘，建立“为什么 SLAM 是最小二乘”这一全书地基）与第 5 节的核心结论（正规方程、信息矩阵、CRLB）。\emph{次优}是第 2 节的若干高阶恒等式（SMW、Isserlis、香农信息、KL）、第 6 节的批量$\leftrightarrow$递归等价证明、第 7 节因子图的变量消元。\emph{进阶}是连续时间高斯过程回归与 STEAM（\cref{sec:est-steam} 完整展开），以及 iSAM/Bayes 树增量平滑（点到为止，给出去向）。虚线指向的“滤波”“GN/LM 求解器”是本章\emph{刻意推迟}的两支：分别见 \cref{ch:eskf} 与 \cref{ch:nlopt}。

\subsection*{前置知识桥接}
\cref{ch:rigid_body} 与 \cref{ch:lie} 解决了“如何\emph{表示}并\emph{微分}位姿”——当残差 $\mathbf{e}(\mathbf{T})$ 含位姿时，求导用右扰动 $\mathbf{T}\,\mathrm{Exp}(\delta\boldsymbol{\xi})$、右雅可比 $\mathbf{J}_r$。本章接着回答：当 $\mathbf{T}$ 是\emph{待估计}量时，从带噪观测把它\emph{估}出来，目标函数从何而来、怎么解、解得多准。所需概率论只需高中概率 $+$ 线性代数；高斯分布的全部必要性质本章 \cref{sec:est-gauss} \emph{自给自足、逐步推导}，无需另查。读到含 $\SE(3)$ 状态时，把欧氏增量 $\delta\mathbf{x}$ 替换为流形右扰动 $\delta\boldsymbol{\xi}$（即 \cref{ch:lie} 的 $\boxplus/\boxminus$）即可，本章先在向量空间把估计骨架讲透。

\subsection*{如果跳过本章会怎样}
\begin{itemize}
  \item \cref{ch:eskf} 的卡尔曼滤波，本质是本章\emph{批量 MAP} 在马尔可夫假设下的\emph{前向递归特例}（本章 \cref{sec:est-batch-recursive} 给出完整证明）；不懂 MAP，就只能把 KF 的“预测–更新”当公式背；
  \item \cref{ch:nlopt} 的高斯牛顿、列文伯格–马夸尔特，是\emph{求解}本章正规方程的迭代法；不先写出最小二乘目标与它的稀疏结构，就无从谈“高效解它”；
  \item \cref{ch:vo}、\cref{ch:vio}、\cref{ch:lidar_slam} 的后端，几乎都是“在因子图上做 MAP”；第 7 节的因子图语言是读懂它们的入场券。
\end{itemize}

\begin{table}[H]\centering\small
\caption{本章阅读方式建议}
\begin{tabular}{lll}
\toprule
阅读方式 & 时间 & 适合谁 \\
\midrule
精读（含全部推导与练习） & 6--8 小时 & 首次系统学习、要做后端/滤波推导 \\
速读（跳过 \cref{sec:est-gauss} 高阶恒等式与 \cref{sec:est-batch-recursive} 细节） & 2 小时 & 已懂 MLE/MAP、想对齐本书记号与因子图视角 \\
速查（只看小结的符号表 + 三张速查表） & 15 分钟 & 写代码时回来查残差/信息矩阵形式 \\
\bottomrule
\end{tabular}
\end{table}

\begin{note}
\textbf{本章记号与约定.}\;
\emph{状态} $\mathbf{x}$（在 SLAM 中含全部待估位姿与路标；需要时单列路标 $\mathbf{y}$ 或 $\boldsymbol{\ell}$）；\emph{观测} $\mathbf{z}$（递归小节按习惯写 $\mathbf{y}_k$）；\emph{控制/输入} $\mathbf{u}$（确定输入也记 $\mathbf{v}_k$）。
\emph{过程噪声} $\mathbf{w}\sim\mathcal{N}(\mathbf{0},\boldsymbol{\Sigma}_{\mathbf{w}})$、\emph{观测噪声} $\mathbf{v}\sim\mathcal{N}(\mathbf{0},\boldsymbol{\Sigma}_{\mathbf{v}})$。
\emph{单位矩阵}记 $\mathbf{I}$。\emph{信息矩阵} $\boldsymbol{\Lambda}:=\boldsymbol{\Sigma}^{-1}$（精度矩阵），\emph{信息向量} $\boldsymbol{\eta}:=\boldsymbol{\Sigma}^{-1}\boldsymbol{\mu}$。
\emph{先验/预测}估计戴 check $\check{\mathbf{x}},\check{\boldsymbol{\Sigma}}$，\emph{后验/校正}估计戴 hat $\hat{\mathbf{x}},\hat{\boldsymbol{\Sigma}}$（沿用 Barfoot\cite{barfoot2024state}）。\emph{连续时间/STEAM 一节}（\cref{sec:est-steam}）按滤波惯例把状态/轨迹协方差记 $\check{\mathbf{P}}/\hat{\mathbf{P}}$（即 $\check{\boldsymbol{\Sigma}}/\hat{\boldsymbol{\Sigma}}$ 的滤波惯用写法，与 \cref{ch:eskf}、\cref{ch:nlopt} 的 Laplace 协方差 $\hat{\mathbf{P}}=\boldsymbol{\Lambda}^{-1}$ 一致），本章其余处协方差一律用 $\boldsymbol{\Sigma}$。
\emph{平方马氏距离} $\|\mathbf{e}\|_{\boldsymbol{\Sigma}}^2:=\mathbf{e}^\top\boldsymbol{\Sigma}^{-1}\mathbf{e}$。

\textbf{跨书差异（重要，避免混淆）}：(1) 噪声协方差的字母三书不一致——Barfoot\cite{barfoot2024state} 记过程 $\mathbf{Q}$、观测 $\mathbf{R}$；十四讲\cite{gaoxiang2019slam14} 恰好\emph{相反}（过程 $\mathbf{R}$、观测 $\mathbf{Q}$）；Handbook\cite{carlone2026handbook} 记观测 $\boldsymbol{\Sigma}_R$、运动 $\boldsymbol{\Sigma}_Q$、里程计 $\boldsymbol{\Sigma}_S$。本书\textbf{一律}用带下标的 $\boldsymbol{\Sigma}_{\mathbf{w}},\boldsymbol{\Sigma}_{\mathbf{v}}$，既消歧义，又把字母 $\mathbf{R}$ 留给旋转矩阵（\cref{ch:lie}）。(2) 单位阵 Barfoot 用粗体 $\mathbf{1}$、本书用 $\mathbf{I}$；Barfoot 的“信息矩阵 $\mathbf{I}_k$”在本书写 $\boldsymbol{\Lambda}_k$。(3) Handbook 的平方根因子取\emph{上三角} $\mathbf{R}$（$\boldsymbol{\Lambda}=\mathbf{R}^\top\mathbf{R}$），与下三角 Cholesky $\mathbf{L}=\mathbf{R}^\top$ 互为转置；本书随上下文标明。
\end{note}

% ════════════════════════════════════════════════════════════════
\section{状态估计是一个概率问题}
\label{sec:est-why}

\paragraph{动机：与其假设数据符合方程，不如在噪声中估计。}
SLAM 的两条基本方程——\emph{运动方程}（机器人怎样动）与\emph{观测方程}（传感器看到什么）——由于噪声，等式\emph{必定}不精确成立\cite{gaoxiang2019slam14}。即便相机几乎完美符合针孔模型，得到的数据仍受各种未知噪声污染。于是范式转变：
\begin{quote}
\emph{与其假设数据必须符合方程，不如讨论如何在有噪声的数据中进行准确的\textbf{状态估计}。}
\end{quote}
这一句是本书后端的总纲：我们不“解”方程，而是\emph{估计}最可能的状态及其不确定度。

\paragraph{SLAM 的概率模型。}
设 $k=1,\dots,N$ 为时刻，路标 $j=1,\dots,M$。经典离散 SLAM 模型为\cite{gaoxiang2019slam14}
\begin{equation}
  \begin{cases}
    \mathbf{x}_k=f(\mathbf{x}_{k-1},\mathbf{u}_k)+\mathbf{w}_k, & \text{运动方程},\\
    \mathbf{z}_{k,j}=h(\mathbf{x}_k,\mathbf{y}_j)+\mathbf{v}_{k,j}, & \text{观测方程},
  \end{cases}
  \qquad
  \mathbf{w}_k\sim\mathcal{N}(\mathbf{0},\boldsymbol{\Sigma}_{\mathbf{w}}),\;
  \mathbf{v}_{k,j}\sim\mathcal{N}(\mathbf{0},\boldsymbol{\Sigma}_{\mathbf{v}}).
  \label{eq:slam-model}
\end{equation}
对视觉 SLAM，位姿 $\mathbf{x}_k$ 由 $\mathbf{T}_k\in\SE(3)$ 表达，观测方程是针孔投影 $s\,\mathbf{z}_{k,j}=\mathbf{K}(\mathbf{R}_k\mathbf{y}_j+\mathbf{t}_k)$（$\mathbf{K}$ 为内参，$s$ 为深度）。在噪声影响下，我们要从带噪的 $\mathbf{z},\mathbf{u}$ 推断位姿与地图\emph{及其概率分布}——这正是一个状态估计问题。

\paragraph{把所有未知量一起估：批量后验。}
把全部待估量堆为 $\mathbf{x}=\{\mathbf{x}_{1:N},\mathbf{y}_{1:M}\}$。从概率视角，对状态的估计就是求\emph{后验分布}\cite{gaoxiang2019slam14,carlone2026handbook}
\begin{equation}
  p(\mathbf{x}\mid\mathbf{z},\mathbf{u}).
  \label{eq:posterior}
\end{equation}
特别地，若没有控制输入、只有图像（只用观测方程），则退化为 $p(\mathbf{x}\mid\mathbf{z})$，此即\emph{运动恢复结构}（SfM）。得到这样一个分布的过程，统称\textbf{概率推断}（probabilistic inference）\cite{carlone2026handbook}。要做推断，须先指定变量的概率模型，以及它们如何产生（不确定的）测量——这正是下文概率工具箱要给的。

\paragraph{两种范式：批量 vs 递归。}
处理 \cref{eq:posterior} 有两条路\cite{gaoxiang2019slam14}：
\begin{itemize}
  \item \textbf{递归/滤波}（incremental）：数据随时间到来，持有当前估计、用新数据更新它（卡尔曼滤波及其衍生）。仅关心当前时刻 $\mathbf{x}_k$。
  \item \textbf{批量}（batch）：把数据“攒”起来一并处理。可在更大范围达到最优，被认为优于传统滤波\cite{dellaert2017factor}，是当代视觉 SLAM 的主流。
\end{itemize}
理论上批量更易讲清，且\emph{理解了批量也更容易理解递归}（\cref{sec:est-batch-recursive} 会证明递归是批量的高效特例）。故本章主讲批量；递归滤波的完整递推留 \cref{ch:eskf}。

\begin{pitfall}[概念误区：把 SLAM 当成“解一个方程组”]
\cref{eq:slam-model} 看起来像一组方程，初学者容易想“凑一组 $\mathbf{x}$ 让等式成立”。\textbf{现象}：试图用高斯消元或求逆“精确解”，发现方程数（观测）通常\emph{远多于}未知数，且每条都带噪——线性系统超定，精确解一般\textbf{不存在}；勉强取最小范数解则完全忽略各观测的可信度差异。\textbf{根本原因}：观测是随机变量的实现，模型等式只在期望意义下成立。\textbf{正确做法}：在所有候选 $\mathbf{x}$ 中，找\emph{最可能}产生这批观测的那个，并给出它的不确定度。这一步“解方程 $\to$ 概率估计”的转身，是本章其余内容的全部动机。\textbf{自检}：写出目标后看它是不是 $\arg\min$ 一个加权残差平方和，而非 $\arg$ 让残差为零。
\end{pitfall}

\begin{practice}
\begin{enumerate}
  \item （概念）对一个只有 2 个未知位姿、却有 5 条带噪相对观测的位姿图，说明为何“解方程组”无解，而“最小化加权残差”有唯一解。答案线索：观测数 $>$ 未知数 $\Rightarrow$ 超定 $\Rightarrow$ 转最小二乘。
  \item 写出纯 SfM（无运动方程）的后验 $p(\mathbf{x}\mid\mathbf{z})$，指出与 \cref{eq:posterior} 的区别仅在于缺少运动因子。答案线索：去掉 $\mathbf{u}$ 与运动项。
\end{enumerate}
\end{practice}

% ════════════════════════════════════════════════════════════════
\section{概率与高斯工具箱}
\label{sec:est-gauss}

本节\emph{从零自给}本章所需的全部概率与高斯结论，每条都给完整推导，使读者无需另查即可掌握。组织顺序：概率密度基元（贝叶斯、边缘化、期望、独立/不相关）$\to$ 多元高斯定义 $\to$ 联合高斯的条件化与边缘化（舒尔补）$\to$ 信息形式与协方差形式的对偶 $\to$ 线性变换与过非线性（EKF 来源）$\to$ 归一化乘积（信息相加，正规方程的根）$\to$ 马氏距离与卡方（MAP$\leftrightarrow$最小二乘的桥）$\to$ SMW 矩阵求逆引理 $\to$ CRLB 与 Fisher 信息。所有结论摘自 Barfoot\cite{barfoot2024state} 第 1 章（其内部式号为 2.x），本节将其完整复现并合入十四讲\cite{gaoxiang2019slam14} 的直觉。

\subsection{概率密度、贝叶斯、边缘化}
\label{sec:est-prob-basics}

\paragraph{概率密度函数。}
随机变量 $\mathbf{x}$ 按某密度 $p(\mathbf{x})$ 分布——一个非负函数，满足\emph{总概率公理} $\int p(\mathbf{x})\,\mathrm{d}\mathbf{x}=1$。注意这是\emph{密度}而非概率：概率 $=$ 密度下的面积，$\Pr(\mathbf{x}\in\Omega)=\int_\Omega p(\mathbf{x})\,\mathrm{d}\mathbf{x}$；累积分布 $P(x)=\int_{-\infty}^{x}p(x')\,\mathrm{d}x'$。机器人学的贝叶斯框架里\emph{始终直接用密度}（连续变量），不走 Kolmogorov 三公理 $+$ CDF 的经典路线\cite{barfoot2024state}。

\paragraph{贝叶斯公式：先验、似然、后验。}
联合密度恒可分解为“条件 $\times$ 无条件” $p(\mathbf{x},\mathbf{z})=p(\mathbf{z}\mid\mathbf{x})p(\mathbf{x})=p(\mathbf{x}\mid\mathbf{z})p(\mathbf{z})$，由此得贝叶斯公式\cite{barfoot2024state}
\begin{equation}
  \underbrace{p(\mathbf{x}\mid\mathbf{z})}_{\text{后验}}
  =\frac{\overbrace{p(\mathbf{z}\mid\mathbf{x})}^{\text{似然}}\,\overbrace{p(\mathbf{x})}^{\text{先验}}}{p(\mathbf{z})},
  \qquad p(\mathbf{z})=\int p(\mathbf{z}\mid\mathbf{x})p(\mathbf{x})\,\mathrm{d}\mathbf{x}.
  \label{eq:bayes}
\end{equation}
术语：$p(\mathbf{x})$ 是\emph{先验}（封装全部先验信息），$p(\mathbf{x}\mid\mathbf{z})$ 是\emph{后验}（含观测后的信息）。分母 $p(\mathbf{z})$（归一化常数、证据）与 $\mathbf{x}$ 无关，由\emph{边缘化}算得，一般情形计算昂贵；幸而下文的\emph{点估计}并不需要它。

\paragraph{边缘化：从联合密度积掉变量。}
对联合密度 $p(\mathbf{x},\mathbf{y})$ 积掉一个变量得另一个的\emph{边缘}\cite{barfoot2024state}：
\begin{equation}
  \int p(\mathbf{x},\mathbf{y})\,\mathrm{d}\mathbf{x}
  =\int p(\mathbf{y}\mid\mathbf{x})p(\mathbf{x})\,\mathrm{d}\mathbf{x}
  \quad\text{或}\quad
  \int p(\mathbf{x}\mid\mathbf{y})p(\mathbf{y})\,\mathrm{d}\mathbf{x}=\underbrace{\Big(\textstyle\int p(\mathbf{x}\mid\mathbf{y})\,\mathrm{d}\mathbf{x}\Big)}_{=1}p(\mathbf{y})=p(\mathbf{y}).
  \label{eq:marginalize}
\end{equation}
边缘化 $=$ “我不再关心 $\mathbf{x}$，把它的所有可能性加权积掉”。在 SLAM 滑动窗口里，把离开窗口的旧位姿\emph{边缘化}掉，正是这一操作（其高斯特化见 \cref{eq:info-marginal}）。

\paragraph{期望与矩。}
期望算子 $\mathbb{E}[f(\mathbf{x})]=\int f(\mathbf{x})p(\mathbf{x})\,\mathrm{d}\mathbf{x}$（矩阵函数逐元素解释）。零阶矩恒为 $1$（总概率公理）；一阶矩是\emph{均值} $\boldsymbol{\mu}=\mathbb{E}[\mathbf{x}]$；二阶中心矩是\emph{协方差} $\boldsymbol{\Sigma}=\mathbb{E}[(\mathbf{x}-\boldsymbol{\mu})(\mathbf{x}-\boldsymbol{\mu})^\top]$（对称正定）；两变量的\emph{交叉协方差}
\begin{equation}
  \mathrm{cov}(\mathbf{x},\mathbf{y})=\mathbb{E}[(\mathbf{x}-\mathbb{E}[\mathbf{x}])(\mathbf{y}-\mathbb{E}[\mathbf{y}])^\top]=\mathbb{E}[\mathbf{x}\mathbf{y}^\top]-\mathbb{E}[\mathbf{x}]\mathbb{E}[\mathbf{y}]^\top.
  \label{eq:crosscov}
\end{equation}
（更高阶矩为偏度、峰度，多元需张量，本章不需。）

\paragraph{独立与不相关。}
两变量\emph{统计独立}指 $p(\mathbf{x},\mathbf{y})=p(\mathbf{x})p(\mathbf{y})$；\emph{不相关}指 $\mathbb{E}[\mathbf{x}\mathbf{y}^\top]=\mathbb{E}[\mathbf{x}]\mathbb{E}[\mathbf{y}]^\top$（即 $\mathrm{cov}(\mathbf{x},\mathbf{y})=\mathbf{0}$）。一般地\emph{独立 $\Rightarrow$ 不相关}，反之不成立——但\textbf{高斯是例外}（\cref{sec:est-gauss-cond} 证明二者等价）。SLAM 推导常假设各观测、各时刻噪声相互独立以简化计算，这把联合似然因式分解，是后文“批量目标 $=$ 各因子之和”的前提。

\paragraph{归一化乘积：融合两个独立估计。}
同一变量上两密度 $p_1,p_2$ 的\emph{归一化乘积} $p(\mathbf{x})=\eta\,p_1(\mathbf{x})p_2(\mathbf{x})$，$\eta=(\int p_1 p_2\,\mathrm{d}\mathbf{x})^{-1}$。其贝叶斯用途：在均匀先验下融合两条独立观测 $\mathbf{y}_1,\mathbf{y}_2$，
\begin{equation}
  p(\mathbf{x}\mid\mathbf{y}_1,\mathbf{y}_2)=\eta\,p(\mathbf{x}\mid\mathbf{y}_1)\,p(\mathbf{x}\mid\mathbf{y}_2).
  \label{eq:norm-prod}
\end{equation}
\emph{推导}：贝叶斯写 $p(\mathbf{x}\mid\mathbf{y}_1,\mathbf{y}_2)=p(\mathbf{y}_1,\mathbf{y}_2\mid\mathbf{x})p(\mathbf{x})/p(\mathbf{y}_1,\mathbf{y}_2)$；设 $\mathbf{y}_1,\mathbf{y}_2$ 在给定 $\mathbf{x}$ 下条件独立，$p(\mathbf{y}_1,\mathbf{y}_2\mid\mathbf{x})=p(\mathbf{y}_1\mid\mathbf{x})p(\mathbf{y}_2\mid\mathbf{x})=\tfrac{p(\mathbf{x}\mid\mathbf{y}_1)p(\mathbf{y}_1)}{p(\mathbf{x})}\tfrac{p(\mathbf{x}\mid\mathbf{y}_2)p(\mathbf{y}_2)}{p(\mathbf{x})}$（各因子再用贝叶斯），代回即得 \cref{eq:norm-prod}，其中 $\eta=\tfrac{p(\mathbf{y}_1)p(\mathbf{y}_2)}{p(\mathbf{y}_1,\mathbf{y}_2)p(\mathbf{x})}$；先验均匀时 $\eta$ 为常数。这就是“多测量融合 $=$ 后验之乘积”的概率原型，因子图中“多个因子相乘”即此（其高斯特化“信息相加”见 \cref{prop:info-add}）。

\subsection{多元高斯分布}
\label{sec:est-gauss-def}

\begin{definition}[多元高斯分布]\label{def:gaussian}
$\mathbf{x}\sim\mathcal{N}(\boldsymbol{\mu},\boldsymbol{\Sigma})$，$\boldsymbol{\Sigma}$ 对称正定（SPD），密度为
\begin{equation}
  p(\mathbf{x})=\frac{1}{\sqrt{(2\pi)^{N}\det\boldsymbol{\Sigma}}}
  \exp\!\Big(-\tfrac12(\mathbf{x}-\boldsymbol{\mu})^\top\boldsymbol{\Sigma}^{-1}(\mathbf{x}-\boldsymbol{\mu})\Big).
  \label{eq:gaussian}
\end{equation}
指数里的二次型 $\|\mathbf{x}-\boldsymbol{\mu}\|_{\boldsymbol{\Sigma}}^2=(\mathbf{x}-\boldsymbol{\mu})^\top\boldsymbol{\Sigma}^{-1}(\mathbf{x}-\boldsymbol{\mu})$ 即\emph{平方马氏距离}。一维特例 $p(x)=\tfrac{1}{\sqrt{2\pi\sigma^2}}\exp(-\tfrac12\tfrac{(x-\mu)^2}{\sigma^2})$；标准正态 $\mathcal{N}(\mathbf{0},\mathbf{I})$。
\end{definition}

\paragraph{一维高斯的累积分布与误差函数.}
\cref{def:gaussian} 给的是\emph{密度}；其\emph{累积分布函数}（CDF，即 \cref{sec:est-prob-basics} 的 $P(x)=\int_{-\infty}^{x}p(x')\,\mathrm{d}x'$）在一维有闭式
\begin{equation}
  P(x\mid\mu,\sigma^2)=\frac12\Big(1+\mathrm{erf}\big(\tfrac{x-\mu}{\sqrt{2}\,\sigma}\big)\Big),\qquad
  \mathrm{erf}(z)=\frac{2}{\sqrt{\pi}}\int_{0}^{z}\mathrm{e}^{-t^2}\,\mathrm{d}t,
  \label{eq:gauss-cdf}
\end{equation}
其中\emph{误差函数} $\mathrm{erf}(\cdot)$ 定义域 $(-\infty,\infty)$、值域 $(-1,1)$、呈 S 形（sigmoid）。它与 \cref{sec:est-sampling} 采样所用的\emph{分位数函数} $Q=P^{-1}$ 互为反函数：解 $P(x)=\alpha$ 得 $x=\mu+\sqrt{2}\,\sigma\,\mathrm{erf}^{-1}(2\alpha-1)$（前向 CDF 把分位映成概率、分位数函数把概率映回分位），\cref{eq:gauss-sample} 的逆变换采样正用后者。

\noindent 高斯的几个性质是后文支柱（均独立成立、逐条给依据）：
\begin{itemize}
  \item \textbf{均值即众数}：\cref{eq:gaussian} 在 $\mathbf{x}=\boldsymbol{\mu}$ 处取最大（指数为零）。于是“最可能点（mode）”与“均值（mean）”重合——这是稍后“MAP 与贝叶斯均值在高斯下一致”（\cref{thm:map-bayes}）的根。
  \item \textbf{仿射变换仍高斯}：$\mathbf{y}=\mathbf{G}\mathbf{x}+\mathbf{b}$（$\mathbf{G}$ 常阵）$\Rightarrow\mathbf{y}\sim\mathcal{N}(\mathbf{G}\boldsymbol{\mu}+\mathbf{b},\,\mathbf{G}\boldsymbol{\Sigma}\mathbf{G}^\top)$（\cref{sec:est-lin-change} 证）。
  \item \textbf{几何}：$\sqrt{\det\boldsymbol{\Sigma}}$ 正比于\emph{不确定椭球}体积；等概率面 $\|\mathbf{x}-\boldsymbol{\mu}\|_{\boldsymbol{\Sigma}}^2=M^2$（$M=1,2,3,\dots$）是一族同心椭球，$M$ 越大覆盖概率越高。
\end{itemize}

\begin{insight}
高斯之所以在估计里无处不在，有三重理由：(1) 中心极限定理使大量独立微扰之和近高斯；(2) 它\emph{仅靠均值与协方差}两个矩就完全确定，便于传播；(3) 高斯在线性运算（仿射变换、条件化、边缘化、相乘）下\emph{闭合}——结果仍是高斯。正是第三点让“线性高斯估计”有处处可写的闭式，本章一切闭式结论都源于此。
\end{insight}

\paragraph{协方差形式 vs 信息形式（对偶）。}
高斯既可用 $(\boldsymbol{\mu},\boldsymbol{\Sigma})$ 表示，也可用\emph{信息形式} $(\boldsymbol{\eta},\boldsymbol{\Lambda})$ 表示，其中
\begin{equation}
  \boldsymbol{\Lambda}=\boldsymbol{\Sigma}^{-1}\ (\text{信息/精度矩阵}),\qquad
  \boldsymbol{\eta}=\boldsymbol{\Sigma}^{-1}\boldsymbol{\mu}\ (\text{信息向量}).
  \label{eq:info-form}
\end{equation}
（由 $\mathbf{x}\sim\mathcal{N}(\boldsymbol{\mu},\boldsymbol{\Sigma})$ 作线性替换 $\mathbf{u}=\boldsymbol{\Sigma}^{-1}\mathbf{x}$ 得 $\mathbf{u}\sim\mathcal{N}(\boldsymbol{\Sigma}^{-1}\boldsymbol{\mu},\boldsymbol{\Sigma}^{-1})$，是 \cref{sec:est-lin-change} 的特例。）\emph{为什么要两种形式}？因为二者“好算”的操作恰好互补（\cref{tab:cov-info}）：协方差形式下\emph{边缘化}容易（直接取子块）、\emph{条件化}难；信息形式下恰好相反。SLAM 的观测融合是反复\emph{条件化/相乘}，故后端偏爱信息形式。还有一个工程理由：当对某变量“一无所知”时 $\boldsymbol{\Sigma}\to\infty$ 无法表示，而 $\boldsymbol{\Lambda}\to\mathbf{0}$ 一表即得。

\begin{table}[H]\centering\small
\caption{协方差形式与信息形式的对偶（贯穿全章）}
\label{tab:cov-info}
\begin{tabular}{lll}
\toprule
操作 & 协方差形式 $(\boldsymbol{\mu},\boldsymbol{\Sigma})$ & 信息形式 $(\boldsymbol{\eta},\boldsymbol{\Lambda})$ \\
\midrule
条件化 $p(\mathbf{x}_a\mid\mathbf{x}_b)$ & 舒尔补（难，\cref{eq:schur-cond}） & 取子块（易，\cref{eq:info-cond}） \\
边缘化 $p(\mathbf{x}_a)$ & 取子块（易，\cref{eq:cov-marginal}） & 舒尔补（难，\cref{eq:info-marginal}） \\
独立信息融合（相乘） & 逐对 SMW（繁） & \textbf{信息相加}（易，\cref{prop:info-add}） \\
KF 预测步 & $\check{\boldsymbol{\Sigma}}=\mathbf{A}\hat{\boldsymbol{\Sigma}}\mathbf{A}^\top+\boldsymbol{\Sigma}_{\mathbf{w}}$（易） & 需 SMW \\
KF 校正步 & 需 SMW / 增益 & $\hat{\boldsymbol{\Lambda}}=\check{\boldsymbol{\Lambda}}+\mathbf{C}^\top\boldsymbol{\Sigma}_{\mathbf{v}}^{-1}\mathbf{C}$（易） \\
“一无所知” & $\boldsymbol{\Sigma}\to\infty$（不可表示） & $\boldsymbol{\Lambda}\to\mathbf{0}$（可表示） \\
\bottomrule
\end{tabular}
\end{table}

\subsection{联合高斯的条件化与边缘化（舒尔补）}
\label{sec:est-gauss-cond}

这是\emph{高斯推断的基石}，直接服务 MAP、卡尔曼校正、因子图变量消元。设联合高斯
\begin{equation}
  \begin{bmatrix}\mathbf{x}\\\mathbf{y}\end{bmatrix}\sim
  \mathcal{N}\!\left(\begin{bmatrix}\boldsymbol{\mu}_x\\\boldsymbol{\mu}_y\end{bmatrix},
  \begin{bmatrix}\boldsymbol{\Sigma}_{xx}&\boldsymbol{\Sigma}_{xy}\\\boldsymbol{\Sigma}_{yx}&\boldsymbol{\Sigma}_{yy}\end{bmatrix}\right),
  \qquad \boldsymbol{\Sigma}_{yx}=\boldsymbol{\Sigma}_{xy}^\top.
  \label{eq:joint-gauss}
\end{equation}

\begin{theorem}[高斯条件化（舒尔补）]\label{thm:schur-cond}
给定 $\mathbf{y}$ 后，$\mathbf{x}$ 的条件分布仍高斯：
\begin{equation}
  p(\mathbf{x}\mid\mathbf{y})=\mathcal{N}\Big(\boldsymbol{\mu}_x+\boldsymbol{\Sigma}_{xy}\boldsymbol{\Sigma}_{yy}^{-1}(\mathbf{y}-\boldsymbol{\mu}_y),\ \ \boldsymbol{\Sigma}_{xx}-\boldsymbol{\Sigma}_{xy}\boldsymbol{\Sigma}_{yy}^{-1}\boldsymbol{\Sigma}_{yx}\Big),
  \label{eq:schur-cond}
\end{equation}
条件协方差 $\mathbf{S}=\boldsymbol{\Sigma}_{xx}-\boldsymbol{\Sigma}_{xy}\boldsymbol{\Sigma}_{yy}^{-1}\boldsymbol{\Sigma}_{yx}$ 正是 $\boldsymbol{\Sigma}_{yy}$ 的\emph{舒尔补}。
\end{theorem}
\begin{proof}
对联合协方差作块 $\mathbf{LDU}$ 分解\cite{barfoot2024state}
\[
  \begin{bmatrix}\boldsymbol{\Sigma}_{xx}&\boldsymbol{\Sigma}_{xy}\\\boldsymbol{\Sigma}_{yx}&\boldsymbol{\Sigma}_{yy}\end{bmatrix}
  =\begin{bmatrix}\mathbf{I}&\boldsymbol{\Sigma}_{xy}\boldsymbol{\Sigma}_{yy}^{-1}\\\mathbf{0}&\mathbf{I}\end{bmatrix}
   \begin{bmatrix}\mathbf{S}&\mathbf{0}\\\mathbf{0}&\boldsymbol{\Sigma}_{yy}\end{bmatrix}
   \begin{bmatrix}\mathbf{I}&\mathbf{0}\\\boldsymbol{\Sigma}_{yy}^{-1}\boldsymbol{\Sigma}_{yx}&\mathbf{I}\end{bmatrix},
  \quad \mathbf{S}=\boldsymbol{\Sigma}_{xx}-\boldsymbol{\Sigma}_{xy}\boldsymbol{\Sigma}_{yy}^{-1}\boldsymbol{\Sigma}_{yx}.
\]
两边求逆得联合\emph{信息}矩阵
\[
  \begin{bmatrix}\boldsymbol{\Sigma}_{xx}&\boldsymbol{\Sigma}_{xy}\\\boldsymbol{\Sigma}_{yx}&\boldsymbol{\Sigma}_{yy}\end{bmatrix}^{-1}
  =\begin{bmatrix}\mathbf{I}&\mathbf{0}\\-\boldsymbol{\Sigma}_{yy}^{-1}\boldsymbol{\Sigma}_{yx}&\mathbf{I}\end{bmatrix}
   \begin{bmatrix}\mathbf{S}^{-1}&\mathbf{0}\\\mathbf{0}&\boldsymbol{\Sigma}_{yy}^{-1}\end{bmatrix}
   \begin{bmatrix}\mathbf{I}&-\boldsymbol{\Sigma}_{xy}\boldsymbol{\Sigma}_{yy}^{-1}\\\mathbf{0}&\mathbf{I}\end{bmatrix}.
\]
把它代入 \cref{eq:gaussian} 的指数二次型并逐项展开，发现二次型恰好\emph{分裂成两个互不耦合的二次项之和}：
\[
  \Big(\!\begin{bmatrix}\mathbf{x}\\\mathbf{y}\end{bmatrix}\!-\!\begin{bmatrix}\boldsymbol{\mu}_x\\\boldsymbol{\mu}_y\end{bmatrix}\!\Big)^{\!\top}\!\boldsymbol{\Sigma}^{-1}\!\Big(\cdots\Big)
  =\underbrace{\big(\mathbf{x}-\boldsymbol{\mu}_x-\boldsymbol{\Sigma}_{xy}\boldsymbol{\Sigma}_{yy}^{-1}(\mathbf{y}-\boldsymbol{\mu}_y)\big)^\top\mathbf{S}^{-1}\big(\cdots\big)}_{\text{含 }\mathbf{x}\text{，即 }p(\mathbf{x}\mid\mathbf{y})}
  +\underbrace{(\mathbf{y}-\boldsymbol{\mu}_y)^\top\boldsymbol{\Sigma}_{yy}^{-1}(\mathbf{y}-\boldsymbol{\mu}_y)}_{\text{仅含 }\mathbf{y}\text{，即 }p(\mathbf{y})}.
\]
由“和的指数 $=$ 指数之积”，$p(\mathbf{x},\mathbf{y})=p(\mathbf{x}\mid\mathbf{y})p(\mathbf{y})$，且第一项就是均值 $\boldsymbol{\mu}_x+\boldsymbol{\Sigma}_{xy}\boldsymbol{\Sigma}_{yy}^{-1}(\mathbf{y}-\boldsymbol{\mu}_y)$、协方差 $\mathbf{S}$ 的高斯，即 \cref{eq:schur-cond}；第二项是 $\mathcal{N}(\boldsymbol{\mu}_y,\boldsymbol{\Sigma}_{yy})$，即边缘 $p(\mathbf{y})$。
\end{proof}

\noindent\textbf{直觉}：测得 $\mathbf{y}$ 后，$\mathbf{x}$ 的均值朝观测方向\emph{平移} $\boldsymbol{\Sigma}_{xy}\boldsymbol{\Sigma}_{yy}^{-1}(\mathbf{y}-\boldsymbol{\mu}_y)$，协方差从 $\boldsymbol{\Sigma}_{xx}$ \emph{减去}一个 SPD 项 $\boldsymbol{\Sigma}_{xy}\boldsymbol{\Sigma}_{yy}^{-1}\boldsymbol{\Sigma}_{yx}$——\emph{收缩}（信息增多）。把 $\boldsymbol{\Sigma}_{xy}\boldsymbol{\Sigma}_{yy}^{-1}$ 记作“卡尔曼增益型”修正系数，这条恒等式既是卡尔曼滤波更新步的来源（\cref{ch:eskf}），也是因子图变量消元的代数核心（\cref{sec:est-fg}）。

\begin{corollary}[协方差形式边缘化：取子块]\label{cor:cov-marginal}
对 \cref{eq:joint-gauss} 直接积掉 $\mathbf{y}$，边缘仍是高斯，且\emph{只需取相应子块}\cite{barfoot2024state}：
\begin{equation}
  p(\mathbf{x})=\int p(\mathbf{x},\mathbf{y})\,\mathrm{d}\mathbf{y}=\mathcal{N}(\boldsymbol{\mu}_x,\boldsymbol{\Sigma}_{xx}),\qquad
  p(\mathbf{y})=\mathcal{N}(\boldsymbol{\mu}_y,\boldsymbol{\Sigma}_{yy}).
  \label{eq:cov-marginal}
\end{equation}
（上面证明的副产物已给出 $p(\mathbf{y})=\mathcal{N}(\boldsymbol{\mu}_y,\boldsymbol{\Sigma}_{yy})$。）
\end{corollary}

\begin{proposition}[信息形式的条件化与边缘化（对偶）]\label{prop:info-marg-cond}
设联合高斯的\emph{信息}矩阵与信息向量分块为
\begin{equation*}
  \boldsymbol{\Lambda}=\boldsymbol{\Sigma}^{-1}=\begin{bmatrix}\boldsymbol{\Lambda}_{xx}&\boldsymbol{\Lambda}_{xy}\\\boldsymbol{\Lambda}_{yx}&\boldsymbol{\Lambda}_{yy}\end{bmatrix},\qquad
  \boldsymbol{\eta}=\boldsymbol{\Sigma}^{-1}\boldsymbol{\mu}=\begin{bmatrix}\boldsymbol{\eta}_x\\\boldsymbol{\eta}_y\end{bmatrix}.
\end{equation*}
则\emph{条件化}（给定 $\mathbf{y}$ 求 $\mathbf{x}$）在信息形式下\emph{直接取子块}：
\begin{equation}
  \boldsymbol{\Lambda}_{x\mid y}=\boldsymbol{\Lambda}_{xx},\qquad \boldsymbol{\eta}_{x\mid y}=\boldsymbol{\eta}_x-\boldsymbol{\Lambda}_{xy}\mathbf{y};
  \label{eq:info-cond}
\end{equation}
而\emph{边缘化}（积掉 $\mathbf{y}$ 求 $p(\mathbf{x})$）在信息形式下需对 $\boldsymbol{\Lambda}_{yy}$ 做\emph{舒尔补}\cite{barfoot2024state}：
\begin{equation}
  \boxed{\ \boldsymbol{\Sigma}_{xx}^{-1}=\boldsymbol{\Lambda}_{xx}-\boldsymbol{\Lambda}_{xy}\boldsymbol{\Lambda}_{yy}^{-1}\boldsymbol{\Lambda}_{yx},\qquad
  \boldsymbol{\Sigma}_{xx}^{-1}\boldsymbol{\mu}_x=\boldsymbol{\eta}_x-\boldsymbol{\Lambda}_{xy}\boldsymbol{\Lambda}_{yy}^{-1}\boldsymbol{\eta}_y.\ }
  \label{eq:info-marginal}
\end{equation}
\end{proposition}
\noindent\cref{eq:info-marginal} 正是“在信息/Hessian 矩阵上对要边缘化的变量做舒尔补”——滑窗 SLAM、地标边缘化、因子图消元的核心运算。注意信息向量也要同步舒尔补更新（$\boldsymbol{\eta}_x-\boldsymbol{\Lambda}_{xy}\boldsymbol{\Lambda}_{yy}^{-1}\boldsymbol{\eta}_y$），实现时常被忽略，务必保留。\cref{tab:cov-info} 的对偶到此完整：协方差形式边缘易、条件难（舒尔补）；信息形式条件易、边缘难（舒尔补）。

\begin{theorem}[高斯下：独立 $\Longleftrightarrow$ 不相关]\label{thm:gauss-indep}
对高斯 PDF，统计独立与不相关\emph{等价}（一般 PDF 仅“独立 $\Rightarrow$ 不相关”单向）。
\end{theorem}
\begin{proof}
设独立 $p(\mathbf{x}\mid\mathbf{y})=p(\mathbf{x})=\mathcal{N}(\boldsymbol{\mu}_x,\boldsymbol{\Sigma}_{xx})$。与 \cref{eq:schur-cond} 比对，均值修正项与协方差减项须对任意 $\mathbf{y}$ 为零：$\boldsymbol{\Sigma}_{xy}\boldsymbol{\Sigma}_{yy}^{-1}(\mathbf{y}-\boldsymbol{\mu}_y)=\mathbf{0}$ 且 $\boldsymbol{\Sigma}_{xy}\boldsymbol{\Sigma}_{yy}^{-1}\boldsymbol{\Sigma}_{yx}=\mathbf{0}$，由 $\boldsymbol{\Sigma}_{yy}\succ\mathbf{0}$ 推出 $\boldsymbol{\Sigma}_{xy}=\mathbf{0}$。由 \cref{eq:crosscov}，$\boldsymbol{\Sigma}_{xy}=\mathbb{E}[\mathbf{x}\mathbf{y}^\top]-\mathbb{E}[\mathbf{x}]\mathbb{E}[\mathbf{y}]^\top=\mathbf{0}$ 即不相关。反向：设不相关 $\boldsymbol{\Sigma}_{xy}=\mathbf{0}$，代入 \cref{eq:schur-cond} 得 $p(\mathbf{x}\mid\mathbf{y})=\mathcal{N}(\boldsymbol{\mu}_x,\boldsymbol{\Sigma}_{xx})=p(\mathbf{x})$ 即独立。
\end{proof}

\subsection{线性变换与高斯过非线性（EKF 之源）}
\label{sec:est-lin-change}

\begin{proposition}[线性变量替换]\label{prop:lin-change}
设 $\mathbf{x}\sim\mathcal{N}(\boldsymbol{\mu}_x,\boldsymbol{\Sigma}_{xx})$，$\mathbf{y}=\mathbf{G}\mathbf{x}$（$\mathbf{G}\in\mathbb{R}^{M\times N}$ 常阵），则 $\mathbf{y}\sim\mathcal{N}(\mathbf{G}\boldsymbol{\mu}_x,\mathbf{G}\boldsymbol{\Sigma}_{xx}\mathbf{G}^\top)$。
\end{proposition}
\begin{proof}
直接用期望线性性：$\boldsymbol{\mu}_y=\mathbb{E}[\mathbf{G}\mathbf{x}]=\mathbf{G}\boldsymbol{\mu}_x$，$\boldsymbol{\Sigma}_{yy}=\mathbb{E}[\mathbf{G}(\mathbf{x}-\boldsymbol{\mu}_x)(\mathbf{x}-\boldsymbol{\mu}_x)^\top\mathbf{G}^\top]=\mathbf{G}\boldsymbol{\Sigma}_{xx}\mathbf{G}^\top$。当 $\mathbf{G}$ 可逆（$M=N$），亦可用变量替换法验证密度：由 $\mathrm{d}\mathbf{y}=|\det\mathbf{G}|\,\mathrm{d}\mathbf{x}$，代入 $\int p(\mathbf{x})\,\mathrm{d}\mathbf{x}=1$ 并令 $\mathbf{x}=\mathbf{G}^{-1}\mathbf{y}$，归一化常数 $\det\boldsymbol{\Sigma}_{xx}$ 变为 $\det(\mathbf{G}\boldsymbol{\Sigma}_{xx}\mathbf{G}^\top)$、指数二次型变为 $(\mathbf{y}-\mathbf{G}\boldsymbol{\mu}_x)^\top(\mathbf{G}\boldsymbol{\Sigma}_{xx}\mathbf{G}^\top)^{-1}(\mathbf{y}-\mathbf{G}\boldsymbol{\mu}_x)$，正是 $\mathcal{N}(\mathbf{G}\boldsymbol{\mu}_x,\mathbf{G}\boldsymbol{\Sigma}_{xx}\mathbf{G}^\top)$。
\end{proof}
\noindent 反方向（$M<N$，欲从 $\mathbf{y}$ 的统计量恢复 $\mathbf{x}$）协方差会“膨胀”到大空间而失败，须改用\emph{信息形式}：选 $\boldsymbol{\Sigma}_{xx}^{-1}\mathbf{x}=\mathbf{G}^\top\boldsymbol{\Sigma}_{yy}^{-1}\mathbf{y}$，取期望得\emph{测量信息往状态空间投影}的母式
\begin{equation}
  \boldsymbol{\Sigma}_{xx}^{-1}=\mathbf{G}^\top\boldsymbol{\Sigma}_{yy}^{-1}\mathbf{G},\qquad
  \boldsymbol{\Sigma}_{xx}^{-1}\boldsymbol{\mu}_x=\mathbf{G}^\top\boldsymbol{\Sigma}_{yy}^{-1}\boldsymbol{\mu}_y.
  \label{eq:info-project}
\end{equation}
线性高斯 MAP 中每个测量贡献的信息项 $\mathbf{H}^\top\boldsymbol{\Sigma}_{\mathbf{v}}^{-1}\mathbf{H}$（$\mathbf{G}\leftrightarrow\mathbf{H}$、$\boldsymbol{\Sigma}_{yy}\leftrightarrow\boldsymbol{\Sigma}_{\mathbf{v}}$）正是它的特例；配合下文的信息相加，得批量正规方程。

\begin{proposition}[高斯过非线性（线性化）：EKF 两步之源]\label{prop:gauss-nonlinear}
设 $p(\mathbf{y}\mid\mathbf{x})=\mathcal{N}(g(\mathbf{x}),\boldsymbol{\Sigma}_{\mathbf{v}})$、$p(\mathbf{x})=\mathcal{N}(\boldsymbol{\mu}_x,\boldsymbol{\Sigma}_{xx})$，$g$ 非线性。把 $g$ 在 $\boldsymbol{\mu}_x$ 处一阶线性化 $g(\mathbf{x})\approx g(\boldsymbol{\mu}_x)+\mathbf{G}(\mathbf{x}-\boldsymbol{\mu}_x)$（$\mathbf{G}=\partial g/\partial\mathbf{x}|_{\boldsymbol{\mu}_x}$），则
\begin{equation}
  p(\mathbf{y})=\int p(\mathbf{y}\mid\mathbf{x})p(\mathbf{x})\,\mathrm{d}\mathbf{x}
  \;\approx\;\mathcal{N}\big(g(\boldsymbol{\mu}_x),\ \boldsymbol{\Sigma}_{\mathbf{v}}+\mathbf{G}\boldsymbol{\Sigma}_{xx}\mathbf{G}^\top\big).
  \label{eq:gauss-nonlinear}
\end{equation}
\end{proposition}
\begin{proof}[证明骨架]
代入两个高斯，指数里出现 $\mathbf{x}$ 的二次型 $(\mathbf{x}-\boldsymbol{\mu}_x)^\top(\boldsymbol{\Sigma}_{xx}^{-1}+\mathbf{G}^\top\boldsymbol{\Sigma}_{\mathbf{v}}^{-1}\mathbf{G})(\mathbf{x}-\boldsymbol{\mu}_x)$ 与交叉项。定义 $\mathbf{F}^\top(\mathbf{G}^\top\boldsymbol{\Sigma}_{\mathbf{v}}^{-1}\mathbf{G}+\boldsymbol{\Sigma}_{xx}^{-1})=\boldsymbol{\Sigma}_{\mathbf{v}}^{-1}\mathbf{G}$，对 $\mathbf{x}$ \emph{配方（complete the square）}：把含 $\mathbf{x}$ 的部分凑成 $\big((\mathbf{x}-\boldsymbol{\mu}_x)-\mathbf{F}(\mathbf{y}-\boldsymbol{\mu}_y)\big)$ 的高斯，对 $\mathbf{x}$ 积分得常数（吸入归一化）；剩下仅含 $\mathbf{y}$ 的二次型，其系数矩阵化简（末步用 SMW \cref{eq:smw}）为 $(\boldsymbol{\Sigma}_{\mathbf{v}}+\mathbf{G}\boldsymbol{\Sigma}_{xx}\mathbf{G}^\top)^{-1}$，即 \cref{eq:gauss-nonlinear}（配方、积分吸常数、SMW 三步的逐行矩阵代数见 \cref{der:gauss-nonlinear-full}）。
\end{proof}
\noindent \cref{eq:gauss-nonlinear} 的 $\boldsymbol{\Sigma}_{yy}=\boldsymbol{\Sigma}_{\mathbf{v}}+\mathbf{G}\boldsymbol{\Sigma}_{xx}\mathbf{G}^\top$ 就是 EKF \emph{预测步}的协方差传播（确定部分 $\mathbf{G}\boldsymbol{\Sigma}\mathbf{G}^\top$ 加噪声 $\boldsymbol{\Sigma}_{\mathbf{v}}$）；配合 \cref{eq:schur-cond} 的条件化即得 EKF \emph{观测步}。两步可理解为滤波器中信息的\emph{创造与销毁}\cite{barfoot2024state}（完整 EKF 递推见 \cref{ch:eskf}）。

\paragraph{为何要线性化：先看一个能精确算的反例。}
\cref{prop:gauss-nonlinear} 的“仿射近似”绕过了一个本质事实——\textbf{高斯一旦过了非线性，就不再是高斯}。在动手线性化之前，值得用一个标量、\emph{确定性}（无观测噪声，$\boldsymbol{\Sigma}_{\mathbf{v}}=\mathbf{0}$）、且\emph{可逆}的非线性，把变换后的精确密度算到底；这是后文 EKF 一阶线性化、以及立体相机视差–深度反演产生“偏斜”分布（\cref{ch:vo}）的最小直觉来源。

\begin{example}[$y=\exp(x)$：高斯被确定性非线性掰成非高斯]\label{ex:gauss-through-exp}
设 $x\sim\mathcal{N}(0,\sigma^2)$，密度 $p(x)=\tfrac{1}{\sqrt{2\pi\sigma^2}}\exp\!\big(-\tfrac{x^2}{2\sigma^2}\big)$，经\emph{确定性}映射
\begin{equation}
  y=\exp(x)\quad(\text{可逆},\ x=\ln y,\ y>0),\qquad \mathrm{d}y=\exp(x)\,\mathrm{d}x=y\,\mathrm{d}x.
  \label{eq:exp-map}
\end{equation}
变换后的密度\emph{不能}只把 $x=\ln y$ 代进 $p(x)$ 了事——必须带上\emph{雅可比因子} $|\mathrm{d}x/\mathrm{d}y|=1/y$ 才保持归一。由总概率公理换元
\begin{equation}
  1=\int_{-\infty}^{\infty}\!p(x)\,\mathrm{d}x
   =\int_{0}^{\infty}\underbrace{\frac{1}{\sqrt{2\pi\sigma^2}}\exp\!\Big(-\frac{(\ln y)^2}{2\sigma^2}\Big)\frac{1}{y}}_{=\,p(y)}\,\mathrm{d}y,
  \label{eq:lognormal}
\end{equation}
故 $p(y)=\dfrac{1}{y\sqrt{2\pi\sigma^2}}\exp\!\big(-\tfrac{(\ln y)^2}{2\sigma^2}\big)$（即\emph{对数正态分布}）。这\textbf{不是}高斯：它只在 $y>0$ 上有支撑、右偏（有长尾）、众数 $\mathrm{e}^{-\sigma^2}$、中位数 $1$、均值 $\mathrm{e}^{\sigma^2/2}$ 三者\emph{互不重合}（高斯则三者皆为 $\mu$，见 \cref{def:gaussian} 下“均值即众数”）。
\end{example}

\begin{insight}[确定性非线性破坏高斯性，且“均值 $\ne$ 众数”是 EKF 的系统偏差之源]
\cref{ex:gauss-through-exp} 给出三点贯穿后文的直觉。\textbf{(1) 必带雅可比因子}：漏掉 \cref{eq:exp-map} 的 $1/y$ 会得到一个面积 $\ne 1$ 的非法密度——这正是李群上做协方差传播必须带 $\mathbf{J}$、像素反投影到深度必须带视差雅可比的同一回事。\textbf{(2) 非高斯不可避免}：只要 $g$ 非线性，精确 $p(y)$ 一般\emph{无高斯闭式}（多维、且映射含噪不可逆时连换元都失效，\cref{prop:gauss-nonlinear} 的线性化正是为此而生）。\textbf{(3) 众数偏离均值}：EKF 把 $\hat{y}=g(\hat{x})=\mathrm{e}^{0}=1$ 当估计，而真实均值是 $\mathrm{e}^{\sigma^2/2}>1$——线性化丢掉的二阶曲率\emph{系统性}地把估计偏向众数一侧，$\sigma$ 越大偏得越多。这就是 \cref{ch:eskf} 中“EKF 在强非线性下过自信且有偏”、\cref{ch:vo} 立体匹配“近大远小使深度分布偏斜”的共同根：线性化只搬运了一阶信息。
\end{insight}

\subsection{归一化乘积 \texorpdfstring{$=$}{=} 信息相加（正规方程的代数核）}
\label{sec:est-info-add}

\begin{proposition}[高斯归一化乘积 $=$ 独立观测融合 $=$ 信息相加]\label{prop:info-add}
$K$ 个高斯之归一化乘积仍是高斯，且在信息形式下\emph{相加}：
\begin{equation}
  \exp\!\Big(-\tfrac12(\mathbf{x}-\boldsymbol{\mu})^\top\boldsymbol{\Sigma}^{-1}(\mathbf{x}-\boldsymbol{\mu})\Big)
  \equiv\eta\prod_{k=1}^{K}\exp\!\Big(-\tfrac12(\mathbf{x}-\boldsymbol{\mu}_k)^\top\boldsymbol{\Sigma}_k^{-1}(\mathbf{x}-\boldsymbol{\mu}_k)\Big),
  \quad
  \boldsymbol{\Lambda}=\sum_{k}\boldsymbol{\Sigma}_k^{-1},\ \ \boldsymbol{\eta}=\sum_{k}\boldsymbol{\Sigma}_k^{-1}\boldsymbol{\mu}_k.
  \label{eq:info-sum}
\end{equation}
更一般，带线性映射 $\mathbf{G}_k\in\mathbb{R}^{M_k\times N}$（$M_k\le N$）的版本：
\begin{equation}
  \eta\prod_{k=1}^{K}\exp\!\Big(-\tfrac12(\mathbf{G}_k\mathbf{x}-\boldsymbol{\mu}_k)^\top\boldsymbol{\Sigma}_k^{-1}(\mathbf{G}_k\mathbf{x}-\boldsymbol{\mu}_k)\Big)
  \ \Rightarrow\
  \boxed{\ \boldsymbol{\Lambda}=\sum_{k}\mathbf{G}_k^\top\boldsymbol{\Sigma}_k^{-1}\mathbf{G}_k,\quad
  \boldsymbol{\eta}=\sum_{k}\mathbf{G}_k^\top\boldsymbol{\Sigma}_k^{-1}\boldsymbol{\mu}_k.\ }
  \label{eq:info-sum-G}
\end{equation}
\end{proposition}
\begin{proof}
取乘积的负指数 $\tfrac12\sum_k(\mathbf{G}_k\mathbf{x}-\boldsymbol{\mu}_k)^\top\boldsymbol{\Sigma}_k^{-1}(\mathbf{G}_k\mathbf{x}-\boldsymbol{\mu}_k)$，按 $\mathbf{x}$ 收集：二次项系数 $\sum_k\mathbf{G}_k^\top\boldsymbol{\Sigma}_k^{-1}\mathbf{G}_k=:\boldsymbol{\Lambda}$、一次项系数 $-\sum_k\mathbf{G}_k^\top\boldsymbol{\Sigma}_k^{-1}\boldsymbol{\mu}_k=:-\boldsymbol{\eta}$。配方为 $\tfrac12(\mathbf{x}-\boldsymbol{\Lambda}^{-1}\boldsymbol{\eta})^\top\boldsymbol{\Lambda}(\mathbf{x}-\boldsymbol{\Lambda}^{-1}\boldsymbol{\eta})+\text{常数}$，故乘积仍是高斯 $\mathcal{N}(\boldsymbol{\Lambda}^{-1}\boldsymbol{\eta},\boldsymbol{\Lambda}^{-1})$，信息相加。令所有 $\mathbf{G}_k=\mathbf{I}$ 即 \cref{eq:info-sum}。
\end{proof}
\noindent\cref{eq:info-sum-G} 是\textbf{本章最关键的一式}：它就是\emph{线性高斯批量 MAP / 加权最小二乘的正规方程}的雏形——信息矩阵 $\boldsymbol{\Lambda}=\sum_k\mathbf{G}_k^\top\boldsymbol{\Sigma}_k^{-1}\mathbf{G}_k$（各测量信息相加，$\mathbf{G}_k$ 是测量雅可比）、信息向量 $\boldsymbol{\eta}=\sum_k\mathbf{G}_k^\top\boldsymbol{\Sigma}_k^{-1}\boldsymbol{\mu}_k$，解 $\boldsymbol{\mu}=\boldsymbol{\Lambda}^{-1}\boldsymbol{\eta}$。\cref{sec:est-linear} 把它写成 $\mathbf{H}^\top\mathbf{W}^{-1}\mathbf{H}$；\cref{sec:est-fg} 把每个因子识别为一个信息加项——这正是\emph{因子图 $\leftrightarrow$ MAP 等价}的代数根源。$M_k\le N$（行少于列）意味着\emph{单个因子不能定姿，须多因子联合}。

\subsection{马氏距离、卡方与信息度量}
\label{sec:est-mahalanobis}

\begin{proposition}[平方马氏距离 $\sim$ 卡方分布]\label{prop:chi2}
$\mathbf{x}\sim\mathcal{N}(\boldsymbol{\mu},\boldsymbol{\Sigma})$（$N$ 维），则平方马氏距离 $u=(\mathbf{x}-\boldsymbol{\mu})^\top\boldsymbol{\Sigma}^{-1}(\mathbf{x}-\boldsymbol{\mu})\sim\chi^2(N)$，均值 $N$、方差 $2N$。
\end{proposition}
\begin{proof}
标准化 $\mathbf{e}=\boldsymbol{\Sigma}^{-1/2}(\mathbf{x}-\boldsymbol{\mu})\sim\mathcal{N}(\mathbf{0},\mathbf{I})$，则 $u=\mathbf{e}^\top\mathbf{e}=\sum_i e_i^2$ 是 $N$ 个独立标准正态平方之和，按定义服从 $\chi^2(N)$。求均值用迹技巧：$u=\mathrm{tr}(\boldsymbol{\Sigma}^{-1}(\mathbf{x}-\boldsymbol{\mu})(\mathbf{x}-\boldsymbol{\mu})^\top)$，迹与期望均线性可交换，$\mathbb{E}[u]=\mathrm{tr}(\boldsymbol{\Sigma}^{-1}\mathbb{E}[(\mathbf{x}-\boldsymbol{\mu})(\mathbf{x}-\boldsymbol{\mu})^\top])=\mathrm{tr}(\boldsymbol{\Sigma}^{-1}\boldsymbol{\Sigma})=\mathrm{tr}\,\mathbf{I}=N$。方差 $2N$（用四阶矩 Isserlis 定理 \cref{eq:isserlis}）。
\end{proof}
\noindent\textbf{用途}：因 $\mathbb{E}[u]=N$，平方马氏距离均值常用作“数据是否符合某高斯”的卡方一致性检验——估计残差的归一化平方和若显著偏离自由度，说明噪声协方差设定不当或有外点。更要紧的是：高斯的\emph{负对数似然}恰是半个平方马氏距离 $+$ 常数，
\begin{equation}
  -\ln\mathcal{N}(\mathbf{x};\boldsymbol{\mu},\boldsymbol{\Sigma})=\tfrac12\|\mathbf{x}-\boldsymbol{\mu}\|_{\boldsymbol{\Sigma}}^2+\tfrac12\ln\big((2\pi)^N\det\boldsymbol{\Sigma}\big),
  \label{eq:neglog-gauss}
\end{equation}
所以\emph{最大化高斯似然} $\Longleftrightarrow$ \emph{最小化平方马氏距离之和}，即 MAP $\Longleftrightarrow$ 加权最小二乘；马氏距离里的 $\boldsymbol{\Sigma}^{-1}$ 正是加权阵。这是 \cref{sec:est-ls} 全部内容的代数预告。

\paragraph{香农信息（高斯）与 KL 散度。}
高斯的熵 $H(\mathbf{x})=-\mathbb{E}[\ln p(\mathbf{x})]$ 可\emph{逐步}算出。把 \cref{eq:neglog-gauss} 的 $-\ln p(\mathbf{x})=\tfrac12(\mathbf{x}-\boldsymbol{\mu})^\top\boldsymbol{\Sigma}^{-1}(\mathbf{x}-\boldsymbol{\mu})+\tfrac12\ln((2\pi)^N\det\boldsymbol{\Sigma})$ 代入 $H=\mathbb{E}[-\ln p]$，按期望线性性拆成两项：
\begin{equation}
  H(\mathbf{x})=\underbrace{\tfrac12\ln\big((2\pi)^N\det\boldsymbol{\Sigma}\big)}_{\text{常数项（不含 }\mathbf{x}\text{）}}
  +\tfrac12\underbrace{\mathbb{E}\big[(\mathbf{x}-\boldsymbol{\mu})^\top\boldsymbol{\Sigma}^{-1}(\mathbf{x}-\boldsymbol{\mu})\big]}_{=\,\mathbb{E}[u]=N\ (\cref{prop:chi2})},
  \label{eq:gauss-entropy-expand}
\end{equation}
第二项正是 \cref{prop:chi2} 的平方马氏距离期望 $\mathbb{E}[u]=N$。代入并用 $N=N\ln e$ 并项，得\emph{纯协方差的函数}
\begin{equation}
  H(\mathbf{x})=\tfrac12\ln\big((2\pi)^N\det\boldsymbol{\Sigma}\big)+\tfrac12 N=\tfrac12\ln\big((2\pi e)^N\det\boldsymbol{\Sigma}\big),
  \label{eq:gauss-entropy}
\end{equation}
故 $\sqrt{\det\boldsymbol{\Sigma}}$ 正比于不确定椭球体积、不确定性越大熵越高。沿\emph{第 $M$ 层等概率面} $\|\mathbf{x}-\boldsymbol{\mu}\|_{\boldsymbol{\Sigma}}^2=M^2$（$M=1,2,3,\dots$，\cref{def:gaussian} 的几何项）密度\emph{恒为}常值
\begin{equation}
  p(\mathbf{x})\big|_{\|\mathbf{x}-\boldsymbol{\mu}\|_{\boldsymbol{\Sigma}}^2=M^2}=\frac{1}{\sqrt{(2\pi)^N\,\mathrm{e}^{M^2}\det\boldsymbol{\Sigma}}}
  \label{eq:gauss-equiprob}
\end{equation}
（把 $M^2$ 代回 \cref{eq:gaussian} 指数即得）；二维时椭圆内几何面积 $A=M^2\pi\sqrt{\det\boldsymbol{\Sigma}}$ 随 $M$ 增大，而面上密度按 $\mathrm{e}^{-M^2/2}$ 衰减——这把“等概率椭球”从定性图像落成了定量闭式。两高斯之 KL 散度
\begin{equation}
  \mathrm{KL}(p_2\|p_1)=\tfrac12\Big((\boldsymbol{\mu}_2-\boldsymbol{\mu}_1)^\top\boldsymbol{\Sigma}_1^{-1}(\boldsymbol{\mu}_2-\boldsymbol{\mu}_1)+\ln\tfrac{\det\boldsymbol{\Sigma}_1}{\det\boldsymbol{\Sigma}_2}+\mathrm{tr}(\boldsymbol{\Sigma}_1^{-1}\boldsymbol{\Sigma}_2)-N\Big)\ge 0,
  \label{eq:gauss-kl}
\end{equation}
仅当两高斯相同时为零（不对称，非真距离）。它在比较“近似后验 vs 真后验”（如变分推断、滑窗一致性）时用作度量。

\paragraph{互信息：观测能消去多少不确定度。}
熵之外还有一个直接服务“主动 SLAM / 下一最佳观测”的度量：两变量的\emph{互信息}（mutual information）
\begin{equation}
  I(\mathbf{x},\mathbf{y})=\mathbb{E}\!\Big[\ln\tfrac{p(\mathbf{x},\mathbf{y})}{p(\mathbf{x})p(\mathbf{y})}\Big]
  =H(\mathbf{x})+H(\mathbf{y})-H(\mathbf{x},\mathbf{y})\ \ge 0,
  \label{eq:mutual-info}
\end{equation}
度量“知道一个变量能减少另一个变量多少不确定度”\cite{barfoot2024state}。当 $\mathbf{x},\mathbf{y}$ 统计独立时 $p(\mathbf{x},\mathbf{y})=p(\mathbf{x})p(\mathbf{y})$，对数内为 $1$，$I=0$（知道一个对另一个毫无帮助）。对\emph{联合高斯} \cref{eq:joint-gauss}，把 \cref{eq:gauss-entropy} 代入 \cref{eq:mutual-info}，并用分块行列式 $\det\boldsymbol{\Sigma}=\det\boldsymbol{\Sigma}_{yy}\det(\boldsymbol{\Sigma}_{xx}-\boldsymbol{\Sigma}_{xy}\boldsymbol{\Sigma}_{yy}^{-1}\boldsymbol{\Sigma}_{yx})$（条件协方差 $\mathbf{S}$ 的行列式），化简得纯协方差的闭式
\begin{equation}
  I(\mathbf{x},\mathbf{y})=-\tfrac12\ln\det\big(\mathbf{I}-\boldsymbol{\Sigma}_{xx}^{-1}\boldsymbol{\Sigma}_{xy}\boldsymbol{\Sigma}_{yy}^{-1}\boldsymbol{\Sigma}_{yx}\big)
  =\tfrac12\ln\tfrac{\det\boldsymbol{\Sigma}_{xx}}{\det\mathbf{S}},
  \label{eq:gauss-mutual-info}
\end{equation}
（两写法经 Sylvester 行列式定理 $\det(\mathbf{I}-\mathbf{A}\mathbf{B})=\det(\mathbf{I}-\mathbf{B}\mathbf{A})$ 互换）。直觉：$\det\boldsymbol{\Sigma}_{xx}/\det\mathbf{S}$ 是“观测 $\mathbf{y}$ 前后 $\mathbf{x}$ 的不确定椭球体积之比”，其对数（除以 $2$）就是这次观测带来的信息量。主动感知/最优实验设计正是\emph{最大化候选观测的互信息}（等价于最大化信息矩阵的某标量，如 $\det$ 准则 D-optimality）。

\paragraph{Isserlis（Wick）定理。}
零均值高斯 $\mathbf{x}\sim\mathcal{N}(\mathbf{0},\boldsymbol{\Sigma})$ 的高阶矩按“两两配对乘积之和”计算，四变量情形
\begin{equation}
  \mathbb{E}[x_i x_j x_k x_\ell]=\mathbb{E}[x_i x_j]\mathbb{E}[x_k x_\ell]+\mathbb{E}[x_i x_k]\mathbb{E}[x_j x_\ell]+\mathbb{E}[x_i x_\ell]\mathbb{E}[x_j x_k].
  \label{eq:isserlis}
\end{equation}
标量特例 $\mathbb{E}[x^4]=3\sigma^4$。两个常用矩阵结论由此得：$\mathbb{E}[\mathbf{x}\mathbf{x}^\top\mathbf{x}\mathbf{x}^\top]=\boldsymbol{\Sigma}(\mathrm{tr}(\boldsymbol{\Sigma})\mathbf{I}+2\boldsymbol{\Sigma})$；带中间矩阵 $\mathbf{A}$ 的更一般式（逐对配对 + 收集得）
\begin{equation}
  \mathbb{E}[\mathbf{x}\mathbf{x}^\top\mathbf{A}\mathbf{x}\mathbf{x}^\top]=\boldsymbol{\Sigma}\big(\mathrm{tr}(\mathbf{A}\boldsymbol{\Sigma})\mathbf{I}+\mathbf{A}\boldsymbol{\Sigma}+\mathbf{A}^\top\boldsymbol{\Sigma}\big).
  \label{eq:isserlis-A}
\end{equation}
两者用于不确定度的高阶传播（如李群上协方差传播的四阶项，见 \cref{ch:lie}，以及无迹/sigma-point 滤波的矩匹配，见 \cref{ch:eskf}）。

\emph{分块情形（服务李群四阶协方差传播）。}\;把 $\mathbf{x}=[\mathbf{x}_1;\mathbf{x}_2]$ 按 $\boldsymbol{\Sigma}=\big[\begin{smallmatrix}\boldsymbol{\Sigma}_{11}&\boldsymbol{\Sigma}_{12}\\\boldsymbol{\Sigma}_{12}^\top&\boldsymbol{\Sigma}_{22}\end{smallmatrix}\big]$ 分块（$\dim\mathbf{x}_1=N_1,\dim\mathbf{x}_2=N_2$），逐元素配对再收集，得只用\emph{某一子块}作中间“夹心”的四阶矩
\begin{subequations}\label{eq:isserlis-block}
\begin{align}
  \mathbb{E}\big[\mathbf{x}\,\mathbf{x}_1^\top\mathbf{x}_1\,\mathbf{x}^\top\big]
  &=\boldsymbol{\Sigma}\Big(\mathrm{tr}(\boldsymbol{\Sigma}_{11})\mathbf{I}+2\begin{bmatrix}\boldsymbol{\Sigma}_{11}&\boldsymbol{\Sigma}_{12}\\\mathbf{0}&\mathbf{0}\end{bmatrix}\Big),\label{eq:isserlis-block-1}\\
  \mathbb{E}\big[\mathbf{x}\,\mathbf{x}_2^\top\mathbf{x}_2\,\mathbf{x}^\top\big]
  &=\boldsymbol{\Sigma}\Big(\mathrm{tr}(\boldsymbol{\Sigma}_{22})\mathbf{I}+2\begin{bmatrix}\mathbf{0}&\mathbf{0}\\\boldsymbol{\Sigma}_{12}^\top&\boldsymbol{\Sigma}_{22}\end{bmatrix}\Big).\label{eq:isserlis-block-2}
\end{align}
\end{subequations}
\emph{一致性校验}：因 $\mathbf{x}^\top\mathbf{x}=\mathbf{x}_1^\top\mathbf{x}_1+\mathbf{x}_2^\top\mathbf{x}_2$，两式相加应还原无分块的 $\mathbb{E}[\mathbf{x}\mathbf{x}^\top\mathbf{x}\mathbf{x}^\top]$；确有 $\mathrm{tr}(\boldsymbol{\Sigma}_{11})+\mathrm{tr}(\boldsymbol{\Sigma}_{22})=\mathrm{tr}(\boldsymbol{\Sigma})$、两个分块矩阵相加为 $\boldsymbol{\Sigma}$，故和 $=\boldsymbol{\Sigma}(\mathrm{tr}(\boldsymbol{\Sigma})\mathbf{I}+2\boldsymbol{\Sigma})$，与上面 $\mathbb{E}[\mathbf{x}\mathbf{x}^\top\mathbf{x}\mathbf{x}^\top]$ 吻合。\cref{eq:isserlis-block} 在 \cref{ch:lie} 旋转矢量协方差传播的二/四阶矩里被直接调用。

\paragraph{Stein 引理（导数–协方差恒等式）。}
高斯还有一条把“函数的导数期望”与“协方差”联系起来的恒等式：设 $\mathbf{x}\sim\mathcal{N}(\boldsymbol{\mu},\boldsymbol{\Sigma})$、$f(\mathbf{x})$ 一次可微标量函数，则\cite{barfoot2024state}
\begin{equation}
  \mathbb{E}\big[(\mathbf{x}-\boldsymbol{\mu})\,f(\mathbf{x})\big]=\boldsymbol{\Sigma}\,\mathbb{E}\!\Big[\tfrac{\partial f(\mathbf{x})}{\partial\mathbf{x}^\top}\Big],
  \label{eq:stein}
\end{equation}
\emph{推导骨架（分部积分）}：以标量 $x\sim\mathcal{N}(\mu,\sigma^2)$ 为例，关键观察是高斯密度满足 $p'(x)=-\tfrac{x-\mu}{\sigma^2}p(x)$，故 $(x-\mu)p(x)=-\sigma^2 p'(x)$；于是 $\mathbb{E}[(x-\mu)f(x)]=\int(x-\mu)f(x)p\,\mathrm{d}x=-\sigma^2\int f(x)p'(x)\,\mathrm{d}x=\sigma^2\int f'(x)p(x)\,\mathrm{d}x=\sigma^2\mathbb{E}[f']$（分部积分，边界项因高斯尾部衰减为零），即 \cref{eq:stein} 的标量版；多元逐分量同理（矩阵微积分背景见 \cref{ch:matderiv}）。\emph{二阶（两次应用）}还给出一条把\emph{二阶导期望}与协方差相连的恒等式
\begin{equation}
  \mathbb{E}\big[(\mathbf{x}-\boldsymbol{\mu})(\mathbf{x}-\boldsymbol{\mu})^\top f(\mathbf{x})\big]
  =\boldsymbol{\Sigma}\,\mathbb{E}\!\Big[\tfrac{\partial^2 f(\mathbf{x})}{\partial\mathbf{x}^\top\partial\mathbf{x}}\Big]\boldsymbol{\Sigma}+\boldsymbol{\Sigma}\,\mathbb{E}[f(\mathbf{x})],
  \label{eq:stein-2nd}
\end{equation}
它在 sigma-point/无迹滤波的协方差矩匹配里有用（把二阶矩用非线性映射的 Hessian 期望表达）。更一般地对联合高斯 $\mathbf{x},\mathbf{y}$ 有 $\mathrm{cov}(\mathbf{y},f(\mathbf{x}))=\mathrm{cov}(\mathbf{y},\mathbf{x})\,\mathbb{E}[\partial f/\partial\mathbf{x}^\top]$。它的用途：在无须显式线性化的前提下，把“高斯过非线性后的交叉协方差”表达成“非线性映射雅可比的期望”——这正是 sigma-point / 无迹卡尔曼“用样本点统计量代替解析雅可比”之所以成立的代数根据（\cref{ch:eskf}）。

\subsection{随机采样、样本统计量与白化}
\label{sec:est-sampling}

估计之外，从高斯\emph{采样}是仿真、粒子滤波、蒙特卡洛核验的基础；它也把“白化”这一工程操作的来由讲清。

\paragraph{从高斯采样：协方差平方根。}
对协方差作平方根分解 $\boldsymbol{\Sigma}=\mathbf{V}\mathbf{V}^\top$（$\mathbf{V}$ 可取 Cholesky 下三角因子，见 \cref{sec:est-fg-linear}）。先取 $N$ 个独立标准正态样本堆成 $\mathbf{x}_{\text{std}}\sim\mathcal{N}(\mathbf{0},\mathbf{I})$（一维标准正态由 $x=\sqrt{2}\,\mathrm{erf}^{-1}(2y-1)$、$y\sim\mathcal{U}[0,1]$ 经分位数函数得），则
\begin{equation}
  \mathbf{z}_{\text{meas}}=\boldsymbol{\mu}+\mathbf{V}\mathbf{x}_{\text{std}}\sim\mathcal{N}(\boldsymbol{\mu},\boldsymbol{\Sigma}).
  \label{eq:gauss-sample}
\end{equation}
\emph{验证}：均值 $\mathbb{E}[\mathbf{z}]=\boldsymbol{\mu}+\mathbf{V}\mathbb{E}[\mathbf{x}_{\text{std}}]=\boldsymbol{\mu}$；协方差 $\mathbb{E}[(\mathbf{z}-\boldsymbol{\mu})(\mathbf{z}-\boldsymbol{\mu})^\top]=\mathbf{V}\mathbb{E}[\mathbf{x}_{\text{std}}\mathbf{x}_{\text{std}}^\top]\mathbf{V}^\top=\mathbf{V}\mathbf{V}^\top=\boldsymbol{\Sigma}$（\cref{prop:lin-change} 的特例）。

\paragraph{白化是采样的逆操作。}
\cref{sec:est-ls} 的“白化” $\boldsymbol{\Sigma}^{-1/2}\mathbf{e}$ 恰是 \cref{eq:gauss-sample} 的\emph{逆}：采样是“标准正态 $\xrightarrow{\mathbf{V}}$ 带协方差的样本”，白化是“带协方差的残差 $\xrightarrow{\mathbf{V}^{-1}}$ 标准化残差”。这解释了为什么白化后的残差服从 $\mathcal{N}(\mathbf{0},\mathbf{I})$、其平方和服从 $\chi^2$（\cref{prop:chi2}）——一致性检验正建立在此。

\paragraph{Q\,-\,Q（分位–分位）图：直观判一批样本是否服从某分布。}
采样 \cref{eq:gauss-sample} 用到了\emph{分位数函数}（quantile function）$Q(\cdot)$——它是累积分布 $P(x)=\Pr(X\le x)$ 的反函数 $Q=P^{-1}$：给定概率 $\alpha\in[0,1]$，$x=Q(\alpha)$ 是“恰有 $\alpha$ 比例质量落在其左侧”的那个值（如标准正态的 $x=\sqrt{2}\,\mathrm{erf}^{-1}(2\alpha-1)$，\cref{eq:gauss-sample} 下方已用）。分位数函数除了\emph{逆变换采样}（取 $\alpha\sim\mathcal{U}[0,1]$，则 $Q(\alpha)$ 即服从目标分布），还给出一件估计里极常用的诊断工具：\textbf{Q\,-\,Q 图}\cite{barfoot2024state}。

\emph{怎么画}：设有 $N$ 个待检样本 $\{r_i\}$（如白化后的标量残差），排序为顺序统计量 $r_{(1)}\le\cdots\le r_{(N)}$；第 $i$ 个样本的\emph{经验分位水平}取 $\alpha_i=\tfrac{i-1/2}{N}$。以\emph{理论分位} $Q(\alpha_i)$（待检参考分布，如 $\mathcal{N}(0,1)$）为横轴、\emph{样本分位} $r_{(i)}$ 为纵轴，逐点描 $(\,Q(\alpha_i),\,r_{(i)}\,)$。若样本确实来自该参考分布，这些点应\emph{紧贴对角线 $y=x$}；系统性偏离对角线则按形状指认\emph{偏离的种类}：整体平移 $=$ 均值偏（残差有系统偏差，如 \cref{ex:gauss-through-exp} 的 mode$\ne$mean）、整体斜率$\ne1$ $=$ 方差被高估/低估（协方差 $\boldsymbol{\Sigma}_{\mathbf{v}}$ 标定过大/过小）、两端上翘下垂成 “S/反 S” $=$ 重尾或轻尾（有外点或过度自信）。

\begin{figure}[ht]
\centering
\begin{tikzpicture}[>=Stealth, scale=1.0]
  % axes
  \draw[->, gray!60!black] (-2.7,0)--(2.9,0) node[right, font=\scriptsize, black]{理论分位 $Q(\alpha_i)$};
  \draw[->, gray!60!black] (0,-2.7)--(0,2.9) node[above, font=\scriptsize, black]{样本分位 $r_{(i)}$};
  % reference y=x line
  \draw[densely dashed, second!60!black, thick] (-2.5,-2.5)--(2.5,2.5) node[pos=0.96, above left, font=\scriptsize]{$y=x$};
  % (a) good fit: points hug the line (small deterministic scatter around y=x)
  \foreach \q/\dy in {-2.2/0.08,-1.6/-0.06,-1.1/0.05,-0.7/-0.04,-0.35/0.06,-0.1/-0.05,0.15/0.05,0.4/-0.06,0.75/0.04,1.15/-0.05,1.7/0.07,2.25/-0.06}{
    \fill[main!75!black] (\q,{\q+\dy}) circle (1.3pt);
  }
  % (b) heavy-tailed sample: S-shape (hollow), cubic deviation from the line
  \foreach \q in {-2.2,-1.6,-1.1,-0.7,-0.35,-0.1,0.15,0.4,0.75,1.15,1.7,2.25}{
    \pgfmathsetmacro\yy{1.0*\q + 0.06*\q*\q*\q}
    \draw[third!70!black] (\q,\yy) circle (1.4pt);
  }
  \node[font=\scriptsize, main!60!black, anchor=west] at (0.15,-2.0) {实心：样本$\approx$参考（贴线）};
  \node[font=\scriptsize, third!70!black, anchor=east] at (-0.15,2.1) {空心：重尾（两端外偏，S 形）};
\end{tikzpicture}
\caption{Q\,-\,Q 图示意。横轴为参考分布（此处 $\mathcal{N}(0,1)$）的理论分位、纵轴为样本分位。实心点贴对角线 $y=x$，说明样本与参考分布一致；空心点呈 “S 形”两端外偏，是\emph{重尾}（外点多于高斯预期）的典型签名。斜率偏离 $1$ 指示方差标定有误，整体上下平移指示均值偏。（仿 \cite{barfoot2024state} 重绘）}
\label{fig:qq-plot}
\end{figure}

\begin{insight}[Q\,-\,Q 图与卡方/马氏诊断互补]
Q\,-\,Q 图与 \cref{prop:chi2} 的卡方一致性检验测的是\emph{同一件事的两个侧面}：估计是否自洽。区别在\emph{粒度}。卡方/马氏检验把整条（白化）残差压成\emph{一个标量} $\|\mathbf{e}\|_{\boldsymbol{\Sigma}}^2$，看它是否落在 $\chi^2(N)$ 的均值 $N$ 附近（\cref{ch:eskf} 的 NEES/NIS 即对它取多次平均）——给“通过/不通过”的判决，却\emph{不告诉你怎么不对}。Q\,-\,Q 图保留\emph{整个分布形状}：把白化残差的每个标量分量（或多次试验得到的马氏距离 $u_i$ 对 $\chi^2(N)$ 参考）逐点对照理论分位，于是\emph{偏离的形状直接指认病因}——平移$\to$偏置、斜率$\to$协方差缩放、S 形$\to$重尾/外点。实务上二者配合：先用卡方检验快速筛“是否一致”，再用 Q\,-\,Q 图定位“何处、如何不一致”，据此回去调 $\boldsymbol{\Sigma}_{\mathbf{w}}/\boldsymbol{\Sigma}_{\mathbf{v}}$ 或换鲁棒核（\cref{ch:nlopt}）。这也是“比较近似后验 vs 真后验”（如蒙特卡洛核验滤波器、变分推断）最常用的可视化判据之一。
\end{insight}

\paragraph{样本均值与样本协方差（Bessel 修正）。}
反过来，从 $N$ 个样本估计分布的矩：样本均值 $\boldsymbol{\mu}_{\text{meas}}=\tfrac1N\sum_i\mathbf{x}_{i,\text{meas}}$、样本协方差
\begin{equation}
  \boldsymbol{\Sigma}_{\text{meas}}=\frac{1}{N-1}\sum_{i=1}^{N}(\mathbf{x}_{i,\text{meas}}-\boldsymbol{\mu}_{\text{meas}})(\mathbf{x}_{i,\text{meas}}-\boldsymbol{\mu}_{\text{meas}})^\top.
  \label{eq:sample-cov}
\end{equation}
分母取 $N-1$ 而非 $N$ 是\emph{Bessel 修正}：样本协方差用到的是各样本与\emph{样本}均值之差，而样本均值本身又由同一批样本算出，二者轻微相关；可证除以 $N-1$ 才给出真协方差的\emph{无偏}估计。$N$ 大时 $N-1\approx N$，差异可忽略。标定传感器噪声协方差 $\boldsymbol{\Sigma}_{\mathbf{v}}$ 时（用于设权重）正用此式。

\subsection{Sherman–Morrison–Woodbury 矩阵求逆引理}
\label{sec:est-smw}

\begin{lemma}[SMW 矩阵求逆引理]\label{lem:smw}
对相容分块，恒有
\begin{equation}
  (\mathbf{A}^{-1}+\mathbf{B}\mathbf{D}^{-1}\mathbf{C})^{-1}\equiv\mathbf{A}-\mathbf{A}\mathbf{B}(\mathbf{D}+\mathbf{C}\mathbf{A}\mathbf{B})^{-1}\mathbf{C}\mathbf{A},
  \label{eq:smw}
\end{equation}
以及配套三式 $(\mathbf{D}+\mathbf{C}\mathbf{A}\mathbf{B})^{-1}\equiv\mathbf{D}^{-1}-\mathbf{D}^{-1}\mathbf{C}(\mathbf{A}^{-1}+\mathbf{B}\mathbf{D}^{-1}\mathbf{C})^{-1}\mathbf{B}\mathbf{D}^{-1}$ 等。
\end{lemma}
\begin{proof}
把同一分块矩阵 $\big[\begin{smallmatrix}\mathbf{A}^{-1}&-\mathbf{B}\\\mathbf{C}&\mathbf{D}\end{smallmatrix}\big]$ 分别作 LDU 与 UDL 分解并各自求逆，比较两结果的左上块即得 \cref{eq:smw}，比较右下块得第二式（完整代数见 \cref{sec:est-appendix}）。
\end{proof}
\noindent SMW 是“信息形式 $\leftrightarrow$ 协方差形式 $+$ 卡尔曼增益”互转的母恒等式：左端 $(\boldsymbol{\Sigma}^{-1}+\mathbf{C}^\top\mathbf{R}^{-1}\mathbf{C})^{-1}$ 是\emph{信息形式}的后验协方差，右端化为含\emph{增益} $\boldsymbol{\Sigma}\mathbf{C}^\top(\mathbf{C}\boldsymbol{\Sigma}\mathbf{C}^\top+\mathbf{R})^{-1}$ 的\emph{协方差形式}更新。KF 协方差更新、批量与递归的互证、\cref{eq:gauss-nonlinear} 的末步，都靠它。

\subsection{Cramér–Rao 下界与 Fisher 信息}
\label{sec:est-crlb}

\begin{theorem}[Cramér–Rao 下界（CRLB）]\label{thm:crlb}
设参数 $\boldsymbol{\theta}$ 影响 $p(\mathbf{x}\mid\boldsymbol{\theta})$，则基于测量 $\mathbf{x}_{\text{meas}}$ 的任意\emph{无偏}估计 $\hat{\boldsymbol{\theta}}$ 的协方差被 Fisher 信息矩阵之逆下界：
\begin{equation}
  \mathrm{cov}(\hat{\boldsymbol{\theta}}\mid\mathbf{x}_{\text{meas}})\succeq\boldsymbol{\mathcal{I}}_{\boldsymbol{\theta}}^{-1},\qquad
  \boldsymbol{\mathcal{I}}_{\boldsymbol{\theta}}=\mathbb{E}\Big[\tfrac{\partial^2(-\ln p(\mathbf{x}\mid\boldsymbol{\theta}))}{\partial\boldsymbol{\theta}^\top\partial\boldsymbol{\theta}}\Big].
  \label{eq:crlb}
\end{equation}
\end{theorem}
\noindent CRLB 给出“给定测量后参数估计能有多确定”的根本极限。\textbf{高斯/批量例（达到 CRLB）}：设 $K$ 个独立测量 $\mathbf{x}_k\sim\mathcal{N}(\boldsymbol{\mu},\boldsymbol{\Sigma})$，估计均值 $\boldsymbol{\mu}$。联合负对数似然形如 $\tfrac12(\mathbf{x}-\mathbf{A}\boldsymbol{\mu})^\top\mathbf{B}^{-1}(\mathbf{x}-\mathbf{A}\boldsymbol{\mu})+\text{常数}$，其中 $\mathbf{A}=[\mathbf{I}\ \cdots\ \mathbf{I}]^\top$、$\mathbf{B}=\mathrm{diag}(\boldsymbol{\Sigma},\dots,\boldsymbol{\Sigma})$——\emph{正是批量最小二乘目标}的雏形。Fisher 信息 $\boldsymbol{\mathcal{I}}_{\boldsymbol{\mu}}=\mathbf{A}^\top\mathbf{B}^{-1}\mathbf{A}=K\boldsymbol{\Sigma}^{-1}$（信息可加，$K$ 倍逆协方差），故 CRLB $\mathrm{cov}(\hat{\boldsymbol{\mu}})\succeq\tfrac1K\boldsymbol{\Sigma}$。样本均值 $\hat{\boldsymbol{\mu}}=\tfrac1K\sum_k\mathbf{x}_{\text{meas},k}$ 无偏（$\mathbb{E}[\hat{\boldsymbol{\mu}}]=\boldsymbol{\mu}$），且其协方差用独立性算得 $\mathrm{cov}(\hat{\boldsymbol{\mu}})=\tfrac1{K^2}\sum_k\boldsymbol{\Sigma}=\tfrac1K\boldsymbol{\Sigma}$——\emph{恰好达到 CRLB}，是有效估计器。这预告了 \cref{thm:normal-eq} 的关键结论：线性高斯下，最小二乘解的协方差 $=$ 逆 Fisher 信息 $=$ CRLB，即\emph{无偏估计能达到的方差极限}。

\begin{pitfall}[概念误区：协方差“大”当成“没用”，以及不相关$=$独立]
两个易错点。\textbf{(1)} 协方差大不等于“无信息”——只有\emph{信息矩阵} $\boldsymbol{\Lambda}=\boldsymbol{\Sigma}^{-1}\to\mathbf{0}$ 才表示一无所知。\emph{现象}：笼统说“某方向方差大所以约束弱”，却忽略了它可能仍被其他观测强约束。\emph{根因}：方差是边缘不确定度，信息矩阵才是“知道多少”的直接度量。\emph{正确做法}：谈“某方向约束弱”应看 $\boldsymbol{\Lambda}$ 的对应特征值（小特征值 $=$ 弱约束方向）。\textbf{(2)} “不相关”一般\emph{弱于}“独立”，唯有\textbf{高斯}下两者等价（\cref{thm:gauss-indep}）——\emph{现象}：把高斯专属的便利无意识用到非高斯变量（如重投影残差经过 atan2 后已非高斯），错误地由“协方差为零”推“独立”。\emph{自检}：用到“不相关$\Rightarrow$独立”时，先确认变量确为联合高斯。
\end{pitfall}

\begin{practice}
\begin{enumerate}
  \item 由 \cref{eq:info-form} 证明：若 $\mathbf{x}\sim\mathcal{N}(\boldsymbol{\mu},\boldsymbol{\Sigma})$，则信息向量 $\boldsymbol{\eta}=\boldsymbol{\Lambda}\boldsymbol{\mu}$，且 $\|\mathbf{x}-\boldsymbol{\mu}\|_{\boldsymbol{\Sigma}}^2=\mathbf{x}^\top\boldsymbol{\Lambda}\mathbf{x}-2\boldsymbol{\eta}^\top\mathbf{x}+\text{常数}$。答案线索：直接展开二次型。
  \item 用 \cref{thm:schur-cond} 说明：观测后 $\mathbf{x}$ 的协方差一定\emph{不增}（提示：舒尔补 $\mathbf{S}\preceq\boldsymbol{\Sigma}_{xx}$，因减去 SPD 项）。这对应“信息只增不减”。
  \item （概念）举例说明两个变量不相关但不独立（如 $X\sim\mathcal{U}[-1,1]$、$Y=X^2$），并解释为何高斯由 \cref{thm:gauss-indep} 排除了这种情形。
  \item 推导 \cref{eq:info-marginal}：从 \cref{eq:joint-gauss} 出发，对联合信息矩阵做 $\boldsymbol{\Lambda}_{yy}$ 的舒尔补，验证它给出 $p(\mathbf{x})$ 的信息矩阵。答案线索：信息形式边缘$=$协方差形式条件的对偶。
  \item （跨小节）用 \cref{eq:smw} 证明 $\boldsymbol{\Sigma}\mathbf{C}^\top(\mathbf{C}\boldsymbol{\Sigma}\mathbf{C}^\top+\mathbf{R})^{-1}=(\boldsymbol{\Sigma}^{-1}+\mathbf{C}^\top\mathbf{R}^{-1}\mathbf{C})^{-1}\mathbf{C}^\top\mathbf{R}^{-1}$（卡尔曼增益两种写法），这正是 \cref{ch:eskf} 协方差/信息两形式 KF 等价的关键。
  \item 由 \cref{eq:gauss-mutual-info} 说明：若 $\boldsymbol{\Sigma}_{xy}=\mathbf{0}$（不相关，高斯下即独立），则互信息 $I=0$。再说明为什么“某候选观测的互信息越大、越值得去观测”——联系主动 SLAM。答案线索：$\mathbf{S}=\boldsymbol{\Sigma}_{xx}$ 时 $\det\boldsymbol{\Sigma}_{xx}/\det\mathbf{S}=1$，$\ln 1=0$。
  \item （推导）验证 \cref{eq:gauss-sample}：对 $\mathbf{z}=\boldsymbol{\mu}+\mathbf{V}\mathbf{x}_{\text{std}}$、$\boldsymbol{\Sigma}=\mathbf{V}\mathbf{V}^\top$，用期望直接算出 $\mathbb{E}[\mathbf{z}]=\boldsymbol{\mu}$、$\mathrm{cov}(\mathbf{z})=\boldsymbol{\Sigma}$，并解释它与白化 $\boldsymbol{\Sigma}^{-1/2}\mathbf{e}$ 互为逆操作。
\end{enumerate}
\end{practice}

% ════════════════════════════════════════════════════════════════
\section{最大似然与最大后验}
\label{sec:est-mle-map}

\paragraph{从后验到点估计。}
直接求后验分布 \cref{eq:posterior} 困难（归一化分母昂贵），但求一个\emph{最优点}使后验最大是可行的。这就是\textbf{最大后验估计}（MAP）。用贝叶斯 \cref{eq:bayes} 并丢掉与 $\mathbf{x}$ 无关的分母\cite{gaoxiang2019slam14,carlone2026handbook}：
\begin{equation}
  \mathbf{x}^\star_{\mathrm{MAP}}
  =\arg\max_{\mathbf{x}} p(\mathbf{x}\mid\mathbf{z},\mathbf{u})
  =\arg\max_{\mathbf{x}} \underbrace{p(\mathbf{z}\mid\mathbf{x},\mathbf{u})}_{\text{似然}}\,\underbrace{p(\mathbf{x}\mid\mathbf{u})}_{\text{先验}}.
  \label{eq:map}
\end{equation}
这里 (i) 用贝叶斯把后验写成似然 $\times$ 先验除以证据 $p(\mathbf{z}\mid\mathbf{u})$；(ii) 证据不依赖 $\mathbf{x}$，不影响 $\arg\max$，丢掉。结论：MAP 估计\emph{最大化似然与先验之积}。

\paragraph{无先验时退化为 MLE。}
若我们\emph{不知道}状态大致在何处，便没有先验。此时只最大化似然，得\textbf{最大似然估计}（MLE）：
\begin{equation}
  \mathbf{x}^\star_{\mathrm{MLE}}=\arg\max_{\mathbf{x}} p(\mathbf{z}\mid\mathbf{x}).
  \label{eq:mle}
\end{equation}
直观地讲，似然是“\emph{在当前状态下，会产生怎样的观测}”。既然观测已知，MLE 就是问：“\emph{在什么状态下，最可能产生我现在看到的数据}”——这就是最大似然的全部含义\cite{gaoxiang2019slam14}。

\begin{insight}
MAP $=$ MLE $+$ 先验。当先验为\emph{均匀分布}（“哪儿都一样可能”）时，先验项是常数，\cref{eq:map} 退化为 \cref{eq:mle}——\textbf{均匀先验下 MAP 与 MLE 同解}。SLAM 中的“先验”往往来自上一时刻估计与运动模型（运动因子）、或对首位姿的锚定；纯由图像做重建（SfM/BA 无运动模型）则常用 MLE。从历史看，最小二乘由 Gauss（1809）为天体轨道估计而创，后被证明在无分布假设下也最优（Gauss 1821/1823；Markov 1912 再发现，即高斯–马尔可夫定理，见 \cref{sec:est-ls}）。
\end{insight}

\paragraph{动机与反面：为什么必须区分。}
反面教材：若把一个\emph{本有先验}的问题当 MLE 解会怎样？见下方陷阱——会丢失对全局自由度（gauge）的约束，使解不唯一。正面：保留先验（哪怕只是对首位姿的弱锚定）就能消除这一病态。这与 \cref{sec:est-linear} 的“正规方程可逆性 $=$ 可观测性”是同一回事。

\begin{pitfall}[思维陷阱：丢掉先验导致 gauge 漂移]
\textbf{错误想法}：图省事，把带先验的问题当 MLE 解。\textbf{现象}：若整个问题对某个\emph{全局变换}不敏感（如所有位姿同时平移、整体尺度——所谓 gauge 自由度），缺少先验/锚定会让信息矩阵 $\boldsymbol{\Lambda}=\mathbf{H}^\top\mathbf{W}^{-1}\mathbf{H}$ \emph{奇异}，解不唯一、数值漂移、优化器报“不可解”或给出随机平移的轨迹。\textbf{根本原因}：仅相对观测（如纯里程计、纯 bearing）无法确定绝对位姿；零空间对应不被任何观测约束的状态组合。\textbf{正确做法}：保留先验，或显式加一个\emph{先验因子}锚定首位姿（见 \cref{sec:est-fg}）。\textbf{自检}：检查 $\mathbf{H}$ 是否列满秩 / $\boldsymbol{\Lambda}$ 是否正定；若有零特征值，其特征向量就是漂移方向。
\end{pitfall}

\begin{practice}
\begin{enumerate}
  \item 写出含一个高斯先验 $\mathbf{x}_0\sim\mathcal{N}(\check{\mathbf{x}}_0,\check{\boldsymbol{\Sigma}}_0)$ 的 MAP 目标，指出去掉先验项后它何时仍有唯一 MLE 解（提示：需可观测）。
  \item （概念）纯 bearing-only SLAM 有哪些 gauge 自由度？需要多少个、什么样的先验/绝对观测来锚定？答案线索：2D 含平移 $+$ 旋转 $+$（可能）尺度。
  \item 证明：若先验是均匀分布，则 \cref{eq:map} 与 \cref{eq:mle} 的 $\arg\max$ 相同。答案线索：常数项不影响 $\arg\max$。
\end{enumerate}
\end{practice}

% ════════════════════════════════════════════════════════════════
\section{从 MAP 到最小二乘}
\label{sec:est-ls}

这一节是全书地基：在高斯噪声下，把“最大化概率”一步步化成“最小化误差平方和”。

\paragraph{单次观测：负对数把高斯变成马氏二次型。}
观测 $\mathbf{z}=h(\mathbf{x})+\mathbf{v}$，$\mathbf{v}\sim\mathcal{N}(\mathbf{0},\boldsymbol{\Sigma}_{\mathbf{v}})$，故 $p(\mathbf{z}\mid\mathbf{x})=\mathcal{N}\big(h(\mathbf{x}),\boldsymbol{\Sigma}_{\mathbf{v}}\big)$。对数单调，\emph{最大化}似然等价于\emph{最小化负对数}。由 \cref{eq:neglog-gauss}，
\begin{equation}
  -\ln p(\mathbf{z}\mid\mathbf{x})
  =\tfrac12\big(\mathbf{z}-h(\mathbf{x})\big)^\top\boldsymbol{\Sigma}_{\mathbf{v}}^{-1}\big(\mathbf{z}-h(\mathbf{x})\big)
  +\underbrace{\tfrac12\ln\!\big((2\pi)^M\det\boldsymbol{\Sigma}_{\mathbf{v}}\big)}_{\text{与 }\mathbf{x}\text{ 无关}}.
  \label{eq:neglog}
\end{equation}
第一项以外的常数与 $\mathbf{x}$ 无关，可略。于是单次 MLE 就是
\begin{equation}
  \mathbf{x}^\star=\arg\min_{\mathbf{x}}\ \big(\mathbf{z}-h(\mathbf{x})\big)^\top\boldsymbol{\Sigma}_{\mathbf{v}}^{-1}\big(\mathbf{z}-h(\mathbf{x})\big)
  =\arg\min_{\mathbf{x}}\ \|\mathbf{z}-h(\mathbf{x})\|_{\boldsymbol{\Sigma}_{\mathbf{v}}}^2.
  \label{eq:single-ls}
\end{equation}
这正是一个\emph{加权最小二乘}：权重是\textbf{信息矩阵} $\boldsymbol{\Sigma}_{\mathbf{v}}^{-1}$。

\begin{insight}[信息矩阵即权重]
\cref{eq:single-ls} 把“误差 $\mathbf{z}-h(\mathbf{x})$”按 $\boldsymbol{\Sigma}_{\mathbf{v}}^{-1}$ 加权后求平方和。某次观测\emph{越准}，协方差 $\boldsymbol{\Sigma}_{\mathbf{v}}$ 越“小”、信息矩阵越“大”，该误差项在总目标里\emph{权重越高}——优化会优先满足它。这是“信息矩阵 $=$ 置信度权重”的精确含义。工程上，\emph{白化}（whitening）把这一步显式化：对 $\boldsymbol{\Sigma}_{\mathbf{v}}^{-1/2}$ 左乘残差，$\|\mathbf{e}\|_{\boldsymbol{\Sigma}_{\mathbf{v}}}^2=\|\boldsymbol{\Sigma}_{\mathbf{v}}^{-1/2}\mathbf{e}\|_2^2$\cite{carlone2026handbook}，把带权最小二乘化为普通最小二乘，并\emph{消去测量单位}（长度/角度），使不同传感器的残差能合进同一代价。
\end{insight}

\paragraph{批量：独立性把联合似然因式分解。}
通常假设各时刻输入与观测\emph{相互独立}，于是联合似然分解为各因子之积\cite{gaoxiang2019slam14,barfoot2024state}：
\begin{equation}
  p(\mathbf{z},\mathbf{u}\mid\mathbf{x})
  =\underbrace{p(\mathbf{x}_0\mid\check{\mathbf{x}}_0)}_{\text{先验因子}}\ \underbrace{\prod_{k}p(\mathbf{x}_k\mid\mathbf{x}_{k-1},\mathbf{u}_k)}_{\text{运动因子}}\ \underbrace{\prod_{k,j}p(\mathbf{z}_{k,j}\mid\mathbf{x}_k,\mathbf{y}_j)}_{\text{观测因子}}.
  \label{eq:factorize}
\end{equation}
定义两类\emph{残差}（误差）：
\begin{equation}
  \mathbf{e}_{\mathbf{u},k}=\mathbf{x}_k-f(\mathbf{x}_{k-1},\mathbf{u}_k)\ (\text{运动}),\qquad
  \mathbf{e}_{\mathbf{z},k,j}=\mathbf{z}_{k,j}-h(\mathbf{x}_k,\mathbf{y}_j)\ (\text{观测}).
  \label{eq:residuals}
\end{equation}

\begin{theorem}[高斯下 MAP $\Longleftrightarrow$ 加权最小二乘]\label{thm:map-ls}
在 \cref{eq:slam-model} 的独立零均值高斯噪声假设下，批量 MAP（含先验时）/ MLE（无先验时）等价于最小化
\begin{equation}
  \boxed{\ \min_{\mathbf{x}}\ J(\mathbf{x})
  =\underbrace{\|\mathbf{x}_0-\check{\mathbf{x}}_0\|_{\check{\boldsymbol{\Sigma}}_0}^2}_{\text{先验项}}
  +\sum_{k}\|\mathbf{e}_{\mathbf{u},k}\|_{\boldsymbol{\Sigma}_{\mathbf{w}}}^2
  +\sum_{k,j}\|\mathbf{e}_{\mathbf{z},k,j}\|_{\boldsymbol{\Sigma}_{\mathbf{v}}}^2\ }
  \label{eq:batch-ls}
\end{equation}
（无先验时去掉先验项；常数因子 $\tfrac12$ 不影响 $\arg\min$，本式已吸收）。
\end{theorem}
\begin{proof}
对 \cref{eq:factorize} 两边取负对数：积变和。每个因子是高斯，其负对数按 \cref{eq:neglog-gauss} 给出“半个平方马氏距离 $+$ 与 $\mathbf{x}$ 无关的常数”；常数可丢。先验因子 $p(\mathbf{x}_0\mid\check{\mathbf{x}}_0)=\mathcal{N}(\check{\mathbf{x}}_0,\check{\boldsymbol{\Sigma}}_0)$ 贡献 $\tfrac12\|\mathbf{x}_0-\check{\mathbf{x}}_0\|_{\check{\boldsymbol{\Sigma}}_0}^2$；运动因子 $p(\mathbf{x}_k\mid\mathbf{x}_{k-1},\mathbf{u}_k)=\mathcal{N}\big(f(\mathbf{x}_{k-1},\mathbf{u}_k),\boldsymbol{\Sigma}_{\mathbf{w}}\big)$ 贡献 $\tfrac12\|\mathbf{e}_{\mathbf{u},k}\|_{\boldsymbol{\Sigma}_{\mathbf{w}}}^2$；观测因子贡献 $\tfrac12\|\mathbf{e}_{\mathbf{z},k,j}\|_{\boldsymbol{\Sigma}_{\mathbf{v}}}^2$。求和并去掉公共因子 $\tfrac12$ 即 \cref{eq:batch-ls}。$\arg\max p\Leftrightarrow\arg\min(-\ln p)$，得证。
\end{proof}

\paragraph{Handbook 视角：因子图上同一结论。}
SLAM Handbook\cite{carlone2026handbook} 以\emph{因子图}语言给出等价命题：设每个因子正比于高斯 $\phi_i(\mathbf{x}_i)\propto\exp(-\tfrac12\|\mathbf{z}_i-h_i(\mathbf{x}_i)\|_{\boldsymbol{\Sigma}_i}^2)$，则
\begin{equation}
  \mathbf{x}^\star_{\mathrm{MAP}}=\arg\max_{\mathbf{x}}\prod_i\phi_i(\mathbf{x}_i)=\arg\min_{\mathbf{x}}\sum_i\|\mathbf{z}_i-h_i(\mathbf{x}_i)\|_{\boldsymbol{\Sigma}_i}^2.
  \label{eq:fg-ls-early}
\end{equation}
这与 \cref{eq:batch-ls} 是同一目标，只是把“先验/运动/观测”统一看作\emph{因子}，把残差求和看作\emph{遍历因子}。两书视角对照：Barfoot 从概率密度因式分解 $+$ 负对数严格推出（更严谨），Handbook 直接以因子图为对象（更直觉、更贴近 GTSAM/g2o 代码）。\cref{sec:est-fg} 把因子图作为统一框架展开。

\paragraph{SLAM 最小二乘的三条结构性质。}
\cref{eq:batch-ls} 不是普通最小二乘，它有三个对后端至关重要的结构\cite{gaoxiang2019slam14,carlone2026handbook}：
\begin{enumerate}
  \item \textbf{稀疏}：每个误差项只关联\emph{一两个}状态（运动项连 $\mathbf{x}_{k-1},\mathbf{x}_k$；观测项连 $\mathbf{x}_k,\mathbf{y}_j$）。虽然总维度极高，单项却很简单——这让问题天然稀疏（\cref{sec:est-linear,sec:est-fg}）。
  \item \textbf{李代数 $\Rightarrow$ 无约束}：若用李代数 / 右扰动（\cref{ch:lie}）表示位姿增量，\cref{eq:batch-ls} 是\emph{无约束}优化；若直接用旋转矩阵，则要带上 $\mathbf{R}^\top\mathbf{R}=\mathbf{I},\det\mathbf{R}=1$ 这些令人头大的约束。这正是李代数的价值。残差对位姿求导用右扰动 $\mathbf{T}\,\mathrm{Exp}(\delta\boldsymbol{\xi})$ 与右雅可比 $\mathbf{J}_r$。
  \item \textbf{二次型加权}：误差按信息矩阵加权，准的观测权重大（见上文洞见）。
\end{enumerate}

\paragraph{因子常欠定、目标常非凸。}
两条要紧观察\cite{carlone2026handbook}：(1) \cref{eq:fg-ls-early} 里单个因子\emph{通常对所连变量欠定}——测量维度低于状态维度，单因子只给一个无穷子集相同似然（如相机一个 2D 像素观测与投影到同一像点的整条 3D 射线一致），故须\emph{多因子融合}才能定姿（这与 \cref{eq:info-sum-G} 的 $M_k\le N$ 一致）；(2) $h_i$ 非线性使目标\emph{非凸}，初值差会陷局部极小，这催生了好初值策略与“可证明最优”求解器（留 \cref{ch:nlopt}）。

\paragraph{旁注：最小二乘的最优性不依赖高斯（Gauss–Markov）。}
即便噪声\emph{不}高斯，只要各误差零均值、不相关、方差已知（以正确加权），加权最小二乘仍是\textbf{最佳线性无偏估计}（BLUE）——没有别的线性无偏估计能给更小方差。这就是\emph{高斯–马尔可夫定理}\cite{barfoot2024state}。故“高斯下 MAP $=$ 最小二乘”与“一般噪声下最小二乘 $=$ BLUE”是两条\emph{独立}的、都支持最小二乘的理由。

\begin{pitfall}[概念误区：以为最小二乘“只是工程凑合”]
\textbf{错误想法}：“误差平方和只是一个方便的代价，没有原理依据。”\textbf{实际上}：在高斯噪声下它\emph{精确}等于 MAP（\cref{thm:map-ls}）——最小化它就是在求最可能状态；在一般零均值不相关噪声下它是 BLUE（高斯–马尔可夫）。\textbf{为什么重要}：理解这点才知道权重 $\boldsymbol{\Sigma}_i^{-1}$ 不可随意设——它是噪声协方差的逆，设错会系统性地偏置估计。\textbf{正确做法}：按传感器真实噪声标定 $\boldsymbol{\Sigma}_i$；若噪声有重尾/外点，改用鲁棒核（\cref{ch:nlopt}），它对应把高斯似然换成重尾似然后的 MAP。
\end{pitfall}

\begin{practice}
\begin{enumerate}
  \item 完整写出 \cref{eq:neglog} 到 \cref{eq:single-ls} 的每一步，指出哪一项因与 $\mathbf{x}$ 无关而被丢弃。
  \item 把一个\emph{先验} $\mathbf{x}_0\sim\mathcal{N}(\check{\mathbf{x}}_0,\check{\boldsymbol{\Sigma}}_0)$ 写成 \cref{eq:batch-ls} 里的一项，说明它如何“锚定”问题、消除 gauge 自由度。
  \item （概念）若把两次独立观测合并，用 \cref{prop:info-add} 说明其等效信息矩阵为何是两者之和。答案线索：负对数把乘积变求和，信息相加。
  \item （跨章综合）取 \cref{ch:lie} 的右扰动，对残差 $\mathbf{e}(\mathbf{T})=\mathbf{z}-\mathbf{T}\mathbf{p}$ 写出它对 $\delta\boldsymbol{\xi}$ 的雅可比，并把它代入 \cref{eq:batch-ls} 的一项——这就是后端对位姿求导的完整链条。
\end{enumerate}
\end{practice}

% ════════════════════════════════════════════════════════════════
\section{线性高斯：正规方程与信息矩阵}
\label{sec:est-linear}

当残差对状态\emph{线性}时，\cref{eq:batch-ls} 有闭式解，且一切结构（信息矩阵、协方差、可观测性、CRLB）都能写清楚。非线性情形（真实 SLAM）通过反复线性化归到这里，迭代求解的数值内核留 \cref{ch:nlopt}；本节给“每一步在解什么”。本节主线取自 Barfoot\cite{barfoot2024state} 第 3 章，逐式复现。

\subsection{批量 MAP 的堆叠形式与正规方程}
\label{sec:est-normal}

\paragraph{线性高斯模型。}
线性、时变的离散运动与观测模型\cite{barfoot2024state}
\begin{equation}
  \mathbf{x}_k=\mathbf{A}_{k-1}\mathbf{x}_{k-1}+\mathbf{v}_k+\mathbf{w}_k\ (k=1\dots K),\qquad
  \mathbf{y}_k=\mathbf{C}_k\mathbf{x}_k+\mathbf{n}_k\ (k=0\dots K),
  \label{eq:lg-models}
\end{equation}
$\mathbf{A}_k$ 转移矩阵、$\mathbf{C}_k$ 观测矩阵、$\mathbf{v}_k$ 确定输入、$\mathbf{w}_k\sim\mathcal{N}(\mathbf{0},\boldsymbol{\Sigma}_{\mathbf{w},k})$ 过程噪声、$\mathbf{n}_k\sim\mathcal{N}(\mathbf{0},\boldsymbol{\Sigma}_{\mathbf{v},k})$ 观测噪声、初始 $\mathbf{x}_0\sim\mathcal{N}(\check{\mathbf{x}}_0,\check{\boldsymbol{\Sigma}}_0)$；噪声与初始知识彼此且跨时刻\emph{互不相关}。

\paragraph{MAP 目标 $=$ 各马氏项之和。}
仿 \cref{thm:map-ls}，对联合似然取负对数得（各项均为平方马氏距离）
\begin{equation}
  J(\mathbf{x})=\tfrac12\|\mathbf{x}_0-\check{\mathbf{x}}_0\|_{\check{\boldsymbol{\Sigma}}_0}^2
  +\tfrac12\sum_{k=1}^{K}\|\mathbf{x}_k-\mathbf{A}_{k-1}\mathbf{x}_{k-1}-\mathbf{v}_k\|_{\boldsymbol{\Sigma}_{\mathbf{w},k}}^2
  +\tfrac12\sum_{k=0}^{K}\|\mathbf{y}_k-\mathbf{C}_k\mathbf{x}_k\|_{\boldsymbol{\Sigma}_{\mathbf{v},k}}^2.
  \label{eq:lg-J}
\end{equation}
把全部数据与状态\emph{堆叠}（lifted）：数据列 $\mathbf{z}$、状态列 $\mathbf{x}$、块矩阵 $\mathbf{H}$（上半先验/运动块、下半观测块）、块对角权 $\mathbf{W}$：
\begin{equation}
  \mathbf{z}=\begin{bmatrix}\check{\mathbf{x}}_0\\\mathbf{v}_{1:K}\\\mathbf{y}_{0:K}\end{bmatrix},\quad
  \mathbf{H}=\begin{bmatrix}
    \mathbf{I}&&&\\-\mathbf{A}_0&\mathbf{I}&&\\&\ddots&\ddots&\\&&-\mathbf{A}_{K-1}&\mathbf{I}\\\hline
    \mathbf{C}_0&&&\\&\mathbf{C}_1&&\\&&\ddots&\\&&&\mathbf{C}_K\end{bmatrix},\quad
  \mathbf{W}=\mathrm{diag}\big(\check{\boldsymbol{\Sigma}}_0,\boldsymbol{\Sigma}_{\mathbf{w},1:K},\boldsymbol{\Sigma}_{\mathbf{v},0:K}\big),
  \label{eq:lg-stack}
\end{equation}
于是 \cref{eq:lg-J} 紧凑地写成
\begin{equation}
  J(\mathbf{x})=\tfrac12(\mathbf{z}-\mathbf{H}\mathbf{x})^\top\mathbf{W}^{-1}(\mathbf{z}-\mathbf{H}\mathbf{x}),
  \label{eq:batch-quadratic}
\end{equation}
恰为 $\mathbf{x}$ 的凸二次型。

\begin{theorem}[正规方程]\label{thm:normal-eq}
\cref{eq:batch-quadratic} 的最小值点 $\hat{\mathbf{x}}$ 满足\emph{正规方程}
\begin{equation}
  \boxed{\ \underbrace{(\mathbf{H}^\top\mathbf{W}^{-1}\mathbf{H})}_{\textstyle\boldsymbol{\Lambda}\ \text{信息矩阵}}\,\hat{\mathbf{x}}
  =\underbrace{\mathbf{H}^\top\mathbf{W}^{-1}\mathbf{z}}_{\textstyle\boldsymbol{\eta}\ \text{信息向量}}\ }
  \label{eq:normal-eq}
\end{equation}
且其解的协方差为 $\hat{\boldsymbol{\Sigma}}=\boldsymbol{\Lambda}^{-1}=(\mathbf{H}^\top\mathbf{W}^{-1}\mathbf{H})^{-1}$。
\end{theorem}
\begin{proof}
$J$ 是二次型，梯度 $\partial J/\partial\mathbf{x}^\top=-\mathbf{H}^\top\mathbf{W}^{-1}(\mathbf{z}-\mathbf{H}\mathbf{x})$。令其在 $\hat{\mathbf{x}}$ 处为零即得 \cref{eq:normal-eq}。海森 $\partial^2 J/\partial\mathbf{x}\partial\mathbf{x}^\top=\mathbf{H}^\top\mathbf{W}^{-1}\mathbf{H}\succeq\mathbf{0}$，故为极小。协方差：由高斯后验 $p(\mathbf{x}\mid\mathbf{z})\propto\exp(-J)$，比对 \cref{eq:gaussian} 的二次型，其逆协方差正是海森 $\mathbf{H}^\top\mathbf{W}^{-1}\mathbf{H}$，故 $\hat{\boldsymbol{\Sigma}}=\boldsymbol{\Lambda}^{-1}$。另证（直接取期望）：把 $\mathbf{x}-\hat{\mathbf{x}}=(\mathbf{H}^\top\mathbf{W}^{-1}\mathbf{H})^{-1}\mathbf{H}^\top\mathbf{W}^{-1}\mathbf{s}$（$\mathbf{s}=\mathbf{H}\mathbf{x}-\mathbf{z}$ 为堆叠噪声、$\mathbb{E}[\mathbf{s}\mathbf{s}^\top]=\mathbf{W}$）代入 $\hat{\boldsymbol{\Sigma}}=\mathbb{E}[(\mathbf{x}-\hat{\mathbf{x}})(\cdots)^\top]=(\mathbf{H}^\top\mathbf{W}^{-1}\mathbf{H})^{-1}\mathbf{H}^\top\mathbf{W}^{-1}\,\mathbf{W}\,\mathbf{W}^{-1}\mathbf{H}(\mathbf{H}^\top\mathbf{W}^{-1}\mathbf{H})^{-1}=(\mathbf{H}^\top\mathbf{W}^{-1}\mathbf{H})^{-1}$，与前一致。
\end{proof}

\noindent 三点要记住：
\begin{itemize}
  \item \textbf{信息矩阵 $=$ 逆协方差 $=$ 海森 $=$ Fisher 信息}：$\boldsymbol{\Lambda}=\mathbf{H}^\top\mathbf{W}^{-1}\mathbf{H}$ 同时是“解的不确定度的逆”和“目标函数的曲率”。线性高斯下它还等于 Fisher 信息矩阵（\cref{thm:crlb} 的批量例），故该解恰好达到\emph{克拉默–拉奥下界}（CRLB）——无偏估计的方差极限\cite{barfoot2024state}。
  \item \textbf{左右两端都是累加}：$\mathbf{H}^\top\mathbf{W}^{-1}\mathbf{H}=\sum_i\mathbf{H}_i^\top\boldsymbol{\Sigma}_i^{-1}\mathbf{H}_i$，每条残差贡献一份信息——\cref{prop:info-add} 与 \cref{eq:info-sum-G} 的批量版。
  \item \textbf{稀疏来自马尔可夫性}：因每项只连一两个变量，$\mathbf{H}$ 极稀疏；按时间排序时 $\boldsymbol{\Lambda}=\mathbf{H}^\top\mathbf{W}^{-1}\mathbf{H}$ 是\emph{块三对角}的（\cref{sec:est-batch-recursive} 证）。如何\emph{利用}稀疏高效求解（Cholesky、变量消元、Schur 补）见 \cref{sec:est-fg} 与 \cref{ch:nlopt}；如何把它\emph{递归}化为滤波见 \cref{sec:est-batch-recursive} 与 \cref{ch:eskf}。
\end{itemize}

\subsection{存在唯一性 \texorpdfstring{$=$}{=} 可观测性}
\label{sec:est-observability}

$\hat{\mathbf{x}}$ 存在且唯一 $\iff$ $\boldsymbol{\Lambda}=\mathbf{H}^\top\mathbf{W}^{-1}\mathbf{H}$ 可逆 $\iff$ $\mathbf{H}$ 列满秩（因 $\mathbf{W}^{-1}\succ\mathbf{0}$，秩条件可去掉它，$\mathrm{rank}(\mathbf{H}^\top\mathbf{W}^{-1}\mathbf{H})=\mathrm{rank}\,\mathbf{H}=\dim\mathbf{x}$）。这对应控制论的\emph{可观测性}\cite{barfoot2024state}：

\paragraph{有初始先验时。}
若保留 $\check{\mathbf{x}}_0$，则 \cref{eq:lg-stack} 中 $\mathbf{H}$ 的上半含整套 $\{\mathbf{I},-\mathbf{A}_k\}$ 行阶梯块，只要 $\check{\boldsymbol{\Sigma}}_0\succ\mathbf{0}$、$\boldsymbol{\Sigma}_{\mathbf{w},k}\succ\mathbf{0}$ 就总有唯一解——先验已提供完整解，测量只是调整。这是\emph{充分非必要}条件。

\paragraph{无初始先验时（时不变化简）。}
去掉初始块，需观测把全部状态约束住。对时不变 $\mathbf{A}_k=\mathbf{A},\mathbf{C}_k=\mathbf{C}$、$K\gg N$，用 Cayley–Hamilton 定理（$\mathbf{A}$ 的 $\ge N$ 次幂可由 $\mathbf{I},\dots,\mathbf{A}^{N-1}$ 线性表出）把秩条件归结为\emph{可观测性矩阵}
\begin{equation}
  \mathcal{O}=\begin{bmatrix}\mathbf{C}\\\mathbf{C}\mathbf{A}\\\vdots\\\mathbf{C}\mathbf{A}^{N-1}\end{bmatrix}\quad\text{满秩}:\ \ \mathrm{rank}\,\mathcal{O}=N,
  \label{eq:observability}
\end{equation}
正是 Kalman（1960）的可观测性检验。存在唯一总条件（充分非必要）：$\boldsymbol{\Sigma}_{\mathbf{w},k}\succ\mathbf{0},\ \boldsymbol{\Sigma}_{\mathbf{v},k}\succ\mathbf{0},\ \mathrm{rank}\,\mathcal{O}=N$。\emph{可观测性定义}：系统可观测当且仅当初始状态能由有限时间内的测量唯一推断。

\begin{derivation}[从 $\mathbf{H}^\top$ 满秩到可观测性矩阵：逐步块行消元]
秩条件 $\mathrm{rank}\,\mathbf{H}^\top=N(K+1)$。\textbf{有先验时}，$\mathbf{H}^\top$ 含整套 $\{\mathbf{I},-\mathbf{A}_k^\top\}$ 构成的\emph{行阶梯}块（每块行有“前导 $\mathbf{I}$”），整体已满秩，故 $\check{\boldsymbol{\Sigma}}_0,\boldsymbol{\Sigma}_{\mathbf{w},k}\succ\mathbf{0}$ 即保证唯一解——先验已提供完整解。\textbf{去掉初始块}后，$\mathbf{H}^\top$ 只剩 $K+1$ 个块行（每块 $N$ 维），为运动块 $\{\mathbf{I},-\mathbf{A}_k^\top\}$ 与观测块 $\{\mathbf{C}_k^\top\}$ 的并排。三步\emph{保秩}行变换归约之。\textbf{(1) 置换}：把首块行移到末，其余 $\{\mathbf{I},-\mathbf{A}_k^\top\}$ 部分即标准行阶梯（置换不改秩）。\textbf{(2) 消元}：向末块行依次加“首块行的 $\mathbf{A}_0^\top$ 倍、次块行的 $\mathbf{A}_0^\top\mathbf{A}_1^\top$ 倍、$\dots$”，把末块行运动部分清零，只在观测列留下
\[
  \big[\,\mathbf{C}_0^\top\ \ \mathbf{A}_0^\top\mathbf{C}_1^\top\ \ \mathbf{A}_0^\top\mathbf{A}_1^\top\mathbf{C}_2^\top\ \ \cdots\ \ \mathbf{A}_0^\top\!\cdots\!\mathbf{A}_{K-1}^\top\mathbf{C}_K^\top\,\big].
\]
\textbf{(3) 分块定秩}：此时左下分块为零、左上是满秩 $NK$ 的行阶梯，故总秩 $=NK+\mathrm{rank}(\text{上面观测块})$；要 $\mathbf{H}^\top$ 满秩 $N(K+1)$ 当且仅当该观测块秩 $=N$。\textbf{时不变化简}：设 $\mathbf{A}_k=\mathbf{A},\mathbf{C}_k=\mathbf{C}$、$K\gg N$。由\emph{Cayley--Hamilton 定理}（$\mathbf{A}$ 满足自身特征方程，任意 $k\ge N$ 次幂 $(\mathbf{A}^\top)^{k}$ 可由 $\mathbf{I},\mathbf{A}^\top,\dots,(\mathbf{A}^\top)^{N-1}$ 线性表出），观测块中 $(\mathbf{A}^\top)^{k}\mathbf{C}^\top\,(k\ge N)$ 不增添线性无关列，秩条件\emph{坍缩}为 $\mathrm{rank}\big[\,\mathbf{C}^\top\ \ \mathbf{A}^\top\mathbf{C}^\top\ \ \cdots\ \ (\mathbf{A}^\top)^{N-1}\mathbf{C}^\top\,\big]=N$，两边转置（行秩 $=$ 列秩）即 \cref{eq:observability} 的 $\mathrm{rank}\,\mathcal{O}=N$。\hfill$\square$
\end{derivation}

\begin{insight}[质量–弹簧直觉]
把每个误差项看成一根弹簧，最优解 $\hat{\mathbf{x}}$ 是总弹性能 $J$ 最小的位形。\emph{可观测} $\iff$ 无法整体平移系统而不拉伸任何弹簧 $\iff$ 最低能态唯一。\emph{不可观测}（如整条小车链可一起平移而不改变任一弹簧能量）则最低能态是一整条“谷底”，解不唯一——这与 \cref{sec:est-mle-map} 的 gauge 漂移、信息矩阵奇异是同一现象的三种说法。
\end{insight}

\subsection{三条等价路径：优化 \texorpdfstring{$=$}{=} 贝叶斯 \texorpdfstring{$=$}{=} 信息}
\label{sec:est-three-paths}

\begin{theorem}[线性高斯下 MAP $=$ 贝叶斯后验均值]\label{thm:map-bayes}
对 \cref{eq:lg-models}，完整贝叶斯后验 $p(\mathbf{x}\mid\mathbf{v},\mathbf{y})$ 恰好是高斯，其均值 $=$ MAP 解 $\hat{\mathbf{x}}$、协方差 $=\hat{\boldsymbol{\Sigma}}=\boldsymbol{\Lambda}^{-1}$。
\end{theorem}
\begin{proof}
把整条轨迹 lifted：$\mathbf{x}=\mathbf{A}(\mathbf{v}+\mathbf{w})$（$\mathbf{A}$ 为下三角 lifted 转移矩阵），先验 $p(\mathbf{x}\mid\mathbf{v})=\mathcal{N}(\check{\mathbf{x}},\check{\boldsymbol{\Sigma}})=\mathcal{N}(\mathbf{A}\mathbf{v},\mathbf{A}\mathbf{Q}\mathbf{A}^\top)$（$\mathbf{Q}=\mathrm{diag}(\check{\boldsymbol{\Sigma}}_0,\boldsymbol{\Sigma}_{\mathbf{w},1:K})$）；测量 $\mathbf{y}=\mathbf{C}\mathbf{x}+\mathbf{n}$。先验状态与测量的联合密度
\[
  p(\mathbf{x},\mathbf{y}\mid\mathbf{v})=\mathcal{N}\!\left(\begin{bmatrix}\check{\mathbf{x}}\\\mathbf{C}\check{\mathbf{x}}\end{bmatrix},
  \begin{bmatrix}\check{\boldsymbol{\Sigma}}&\check{\boldsymbol{\Sigma}}\mathbf{C}^\top\\\mathbf{C}\check{\boldsymbol{\Sigma}}&\mathbf{C}\check{\boldsymbol{\Sigma}}\mathbf{C}^\top+\mathbf{R}\end{bmatrix}\right),\quad \mathbf{R}=\mathrm{diag}(\boldsymbol{\Sigma}_{\mathbf{v},0:K}).
\]
用条件化 \cref{eq:schur-cond} 直接写后验
\[
  p(\mathbf{x}\mid\mathbf{v},\mathbf{y})=\mathcal{N}\Big(\check{\mathbf{x}}+\check{\boldsymbol{\Sigma}}\mathbf{C}^\top(\mathbf{C}\check{\boldsymbol{\Sigma}}\mathbf{C}^\top+\mathbf{R})^{-1}(\mathbf{y}-\mathbf{C}\check{\mathbf{x}}),\ \check{\boldsymbol{\Sigma}}-\check{\boldsymbol{\Sigma}}\mathbf{C}^\top(\mathbf{C}\check{\boldsymbol{\Sigma}}\mathbf{C}^\top+\mathbf{R})^{-1}\mathbf{C}\check{\boldsymbol{\Sigma}}\Big).
\]
用 SMW \cref{eq:smw} 化为信息形式
\[
  p(\mathbf{x}\mid\mathbf{v},\mathbf{y})=\mathcal{N}\Big((\check{\boldsymbol{\Sigma}}^{-1}+\mathbf{C}^\top\mathbf{R}^{-1}\mathbf{C})^{-1}(\check{\boldsymbol{\Sigma}}^{-1}\check{\mathbf{x}}+\mathbf{C}^\top\mathbf{R}^{-1}\mathbf{y}),\ (\check{\boldsymbol{\Sigma}}^{-1}+\mathbf{C}^\top\mathbf{R}^{-1}\mathbf{C})^{-1}\Big),
\]
其均值满足 $(\check{\boldsymbol{\Sigma}}^{-1}+\mathbf{C}^\top\mathbf{R}^{-1}\mathbf{C})\hat{\mathbf{x}}=\check{\boldsymbol{\Sigma}}^{-1}\check{\mathbf{x}}+\mathbf{C}^\top\mathbf{R}^{-1}\mathbf{y}$。代入 $\check{\mathbf{x}}=\mathbf{A}\mathbf{v}$、$\check{\boldsymbol{\Sigma}}^{-1}=\mathbf{A}^{-\top}\mathbf{Q}^{-1}\mathbf{A}^{-1}$，并令 $\mathbf{H}=[\mathbf{A}^{-1};\mathbf{C}]$、$\mathbf{W}=\mathrm{diag}(\mathbf{Q},\mathbf{R})$、$\mathbf{z}=[\mathbf{v};\mathbf{y}]$，恰好化为 $(\mathbf{H}^\top\mathbf{W}^{-1}\mathbf{H})\hat{\mathbf{x}}=\mathbf{H}^\top\mathbf{W}^{-1}\mathbf{z}$，\emph{与 MAP 正规方程 \cref{eq:normal-eq} 完全相同}。根本原因：高斯后验的众数 $=$ 均值。
\end{proof}

\begin{insight}[殊途同归]
线性高斯问题有三条等价路径，给出\emph{同一}最优解 $\hat{\mathbf{x}}$ 与协方差 $\hat{\boldsymbol{\Sigma}}$：(i) \textbf{优化视角}——最小化 \cref{eq:batch-quadratic} 令梯度为零得正规方程；(ii) \textbf{贝叶斯视角}——对联合高斯做舒尔补条件化 \cref{eq:schur-cond}（\cref{thm:map-bayes}）；(iii) \textbf{信息视角}——\cref{prop:info-add} 把先验信息与各观测信息相加。三者一致，根本原因是\emph{高斯后验的众数 $=$ 均值}。非线性情形三者开始分化（mean $\ne$ mode、滤波须线性化），这是后续滤波/优化章的分水岭。
\end{insight}

\paragraph{最小可复现例：一维小车。}
体会上述结构的最小例子\cite{gaoxiang2019slam14,barfoot2024state}：小车沿直线运动，
\(
  x_k=x_{k-1}+u_k+w_k,\ w_k\sim\mathcal{N}(0,\sigma_w^2);\
  z_k=x_k+v_k,\ v_k\sim\mathcal{N}(0,\sigma_v^2).
\)
取 $k=1,2,3$、$x_0$ 已知。误差 $e_{u,k}=x_k-x_{k-1}-u_k$、$e_{z,k}=z_k-x_k$，目标
\(
  \min\sum_k e_{u,k}^2/\sigma_w^2+\sum_k e_{z,k}^2/\sigma_v^2.
\)
向量化 $\mathbf{z}-\mathbf{H}\mathbf{x}=\mathbf{e}\sim\mathcal{N}(\mathbf{0},\boldsymbol{\Sigma})$（$\boldsymbol{\Sigma}=\mathrm{diag}(\sigma_w^2,\dots,\sigma_v^2)$），由 \cref{thm:normal-eq} 得闭式
\begin{equation}
  \hat{\mathbf{x}}=(\mathbf{H}^\top\boldsymbol{\Sigma}^{-1}\mathbf{H})^{-1}\mathbf{H}^\top\boldsymbol{\Sigma}^{-1}\mathbf{z}.
  \label{eq:car-solution}
\end{equation}
取 $\sigma_w=\sigma_v=1$ 时可算出 $\boldsymbol{\Lambda}=\mathbf{H}^\top\mathbf{H}$ 是\emph{三对角}阵（主对角 $\{2,3,3,\dots,2\}$、次对角 $-1$），印证“每项只连相邻状态 $\Rightarrow$ 三对角”。此处线性，闭式可直接写出；真实 SLAM 的 $h$ 非线性，须迭代（\cref{ch:nlopt}）。

\begin{practice}
\begin{enumerate}
  \item 推导 \cref{eq:normal-eq}：从 \cref{eq:batch-quadratic} 求 $\partial J/\partial\mathbf{x}=\mathbf{0}$，每步写清。
  \item 对一维小车（$K=3$）写出 $\mathbf{H},\boldsymbol{\Sigma}$ 的具体分块，并证明 $\boldsymbol{\Lambda}=\mathbf{H}^\top\boldsymbol{\Sigma}^{-1}\mathbf{H}$ 为何是\emph{三对角}的——把它与“每项只连相邻状态”对应起来。
  \item （概念）若去掉“$x_0$ 已知”这一先验，$\boldsymbol{\Lambda}$ 会怎样？联系本节“可观测性 / gauge”讨论（提示：首行变为 $1$，仍满秩；若再去掉所有 $z_k$ 则不可观）。
  \item 验证 \cref{thm:map-bayes}：对一维小车，分别用“正规方程”和“舒尔补条件化”算 $\hat{x}_1$，确认两者数值相同。
  \item （跨章综合）说明 \cref{eq:gn-normal}（\cref{ch:nlopt} 的高斯牛顿增量方程）与本节正规方程 \cref{eq:normal-eq} 是同一形式——前者是后者在每个线性化点上的重复求解。
\end{enumerate}
\end{practice}

% ════════════════════════════════════════════════════════════════
\section{批量与递归：等价与稀疏}
\label{sec:est-batch-recursive}

上节给出批量解 $\hat{\mathbf{x}}=(\mathbf{H}^\top\mathbf{W}^{-1}\mathbf{H})^{-1}\mathbf{H}^\top\mathbf{W}^{-1}\mathbf{z}$。直接对 $\boldsymbol{\Lambda}$（维度 $N(K+1)$）求逆是 $O((NK)^3)$，对长轨迹不可接受。本节说明：信息矩阵的\emph{块三对角}稀疏（源于马尔可夫性）使批量可 $O(K)$ 解，且其\emph{前向递归}恰是卡尔曼滤波——\emph{递归是批量的高效特例，无近似}。这清偿了 \cref{sec:est-why} “批量更易讲清、理解批量也更易理解递归”的承诺；完整 KF 递推方程留 \cref{ch:eskf}，本节只证等价与稀疏的来由。

\subsection{块三对角稀疏来自马尔可夫性}
\label{sec:est-sparsity}

把状态按时间排序，\cref{eq:lg-stack} 的信息矩阵
\begin{equation}
  \mathbf{H}^\top\mathbf{W}^{-1}\mathbf{H}=\mathbf{A}^{-\top}\mathbf{Q}^{-1}\mathbf{A}^{-1}+\mathbf{C}^\top\mathbf{R}^{-1}\mathbf{C}
  =\begin{bmatrix}*&*&&\\ *&*&*&\\ &*&*&\ddots\\ &&\ddots&\ddots\end{bmatrix}\ \text{（块三对角）},
  \label{eq:tridiag}
\end{equation}
其中 lifted 转移矩阵 $\mathbf{A}$（\cref{thm:map-bayes} 证明里 $\mathbf{x}=\mathbf{A}(\mathbf{v}+\mathbf{w})$ 的那个）及其逆 $\mathbf{A}^{-1}$ 的显式分块形对照鲜明——$\mathbf{A}$ \emph{稠密}下三角（含转移阵连乘），$\mathbf{A}^{-1}$ 却只剩\emph{两条对角}（以 $K=3$ 示意）：
\begin{equation}
  \mathbf{A}=\begin{bmatrix}
    \mathbf{I}&\mathbf{0}&\mathbf{0}&\mathbf{0}\\
    \mathbf{A}_0&\mathbf{I}&\mathbf{0}&\mathbf{0}\\
    \mathbf{A}_1\mathbf{A}_0&\mathbf{A}_1&\mathbf{I}&\mathbf{0}\\
    \mathbf{A}_2\mathbf{A}_1\mathbf{A}_0&\mathbf{A}_2\mathbf{A}_1&\mathbf{A}_2&\mathbf{I}
  \end{bmatrix},\qquad
  \mathbf{A}^{-1}=\begin{bmatrix}
    \mathbf{I}&\mathbf{0}&\mathbf{0}&\mathbf{0}\\
    -\mathbf{A}_0&\mathbf{I}&\mathbf{0}&\mathbf{0}\\
    \mathbf{0}&-\mathbf{A}_1&\mathbf{I}&\mathbf{0}\\
    \mathbf{0}&\mathbf{0}&-\mathbf{A}_2&\mathbf{I}
  \end{bmatrix}.
  \label{eq:lifted-Ainv}
\end{equation}
$\mathbf{A}$ 第一列是 $\mathbf{A}_{k-1}\!\cdots\!\mathbf{A}_0$ 的连乘（一般稠密），而 $\mathbf{A}^{-1}$ 第 $k$ 行块为 $[\,-\mathbf{A}_{k-1}\ \ \mathbf{I}\,]$，仅\emph{主对角与下一对角}非零——正对应运动方程 $\mathbf{x}_k-\mathbf{A}_{k-1}\mathbf{x}_{k-1}=\mathbf{v}_k+\mathbf{w}_k$ 只耦合相邻两态（这就是稀疏的代数根；\cref{eq:lg-stack} 的 $\mathbf{H}$ 上半即由 $\mathbf{A}^{-1}$ 排成）。配合 $\mathbf{Q}^{-1},\mathbf{R}^{-1}$ 块对角、$\mathbf{C}^\top\mathbf{R}^{-1}\mathbf{C}$ 块对角，三者相加给出\emph{块三对角}。

\begin{insight}[稀疏的物理根源：马尔可夫性]
$\mathbf{A}^{-1}$ 之所以只有两条对角非零，是因为运动模型满足\emph{马尔可夫性质}——$\mathbf{x}_k$ 只直接依赖 $\mathbf{x}_{k-1}$，不直接依赖更早状态。于是误差项只耦合\emph{相邻}时刻，信息矩阵只在相邻块间有非零——这就是块三对角。SLAM 的稀疏性最终都可追溯到这种“局部性”：运动只连相邻位姿，观测只连一个位姿与一个路标。\emph{稀疏不是巧合，而是概率图局部结构的代数投影}（\cref{sec:est-fg} 用因子图把这一点画出来）。
\end{insight}

\paragraph{$O(K)$ 求解：稀疏 Cholesky 平滑。}
块三对角 $\boldsymbol{\Lambda}=\mathbf{L}\mathbf{L}^\top$ 的 Cholesky 因子 $\mathbf{L}$ 仍是块下双对角，分解 $O(N^3(K+1))$；前向解 $\mathbf{L}\mathbf{d}=\mathbf{H}^\top\mathbf{W}^{-1}\mathbf{z}$（前代）、后向解 $\mathbf{L}^\top\hat{\mathbf{x}}=\mathbf{d}$（回代）各 $O(N^3(K+1))$。故批量方程\emph{线性于轨迹长度}求解。这一“前向 $+$ 后向 pass”的递归实现就是经典\emph{固定区间平滑器}（Rauch–Tung–Striebel, RTS），它\emph{无近似}地实现完整批量解\cite{barfoot2024state}。关键论断是：\emph{平滑器 $=$ 批量解的高效（稀疏）实现，而非另一种近似方法}。下面把这套“前向 $+$ 后向 pass”\emph{逐式}写到块一级（\cref{alg:est-cholesky-smoother}），再补上从同一因子 $\mathbf{L}$ 直接\emph{回收边际协方差}的后向递推（\cref{alg:est-cov-recovery}）；这两套递推与卡尔曼滤波 / 经典 RTS 平滑（\cref{ch:eskf} 的 \cref{eq:rts}、\cref{eq:rts-cov}）逐字等价，证据在两个算法盒后给出。

\paragraph{块前向/后向递归：把 $\mathbf{L}\mathbf{L}^\top=\boldsymbol{\Lambda}$ 拆到块一级。}
按时间排序，把块三对角信息矩阵的下三角 Cholesky 因子 $\mathbf{L}$（块下双对角，见 \cref{eq:tridiag}）的非零子块记作
\begin{equation}
  \mathbf{L}=\begin{bmatrix}
    \mathbf{L}_0&&&&\\
    \mathbf{L}_{1,0}&\mathbf{L}_1&&&\\
    &\mathbf{L}_{2,1}&\mathbf{L}_2&&\\
    &&\ddots&\ddots&\\
    &&&\mathbf{L}_{K,K-1}&\mathbf{L}_K
  \end{bmatrix},
  \label{eq:cholesky-blocks}
\end{equation}
其中 $\mathbf{L}_k$ 是第 $k$ 个对角块（方阵），$\mathbf{L}_{k,k-1}$ 是其下一条次对角块。把 \cref{eq:lg-stack} 的 $\mathbf{H},\mathbf{W}$ 代入 $\mathbf{H}^\top\mathbf{W}^{-1}\mathbf{H}=\mathbf{L}\mathbf{L}^\top$ 并\emph{逐块比对}（左乘右乘后第 $(k,k)$、$(k,k\!-\!1)$ 块相等），即得块递归。为简洁先定义两个\emph{中间信息量}（沿用本书 $\boldsymbol{\Lambda}_k$ 记信息块、Barfoot 记之为 $\mathbf{I}_k$\cite{barfoot2024state}）：
\begin{equation}
  \boldsymbol{\Lambda}_0:=\check{\boldsymbol{\Sigma}}_0^{-1}+\mathbf{C}_0^\top\boldsymbol{\Sigma}_{\mathbf{v},0}^{-1}\mathbf{C}_0,\qquad
  \mathbf{q}_0:=\check{\boldsymbol{\Sigma}}_0^{-1}\check{\mathbf{x}}_0+\mathbf{C}_0^\top\boldsymbol{\Sigma}_{\mathbf{v},0}^{-1}\mathbf{y}_0,
  \label{eq:cholesky-init}
\end{equation}
$\boldsymbol{\Lambda}_k$ 聚合“到 $k$ 为止流入 $\mathbf{x}_k$ 的全部信息”、$\mathbf{q}_k$ 是配套信息向量。

\begin{algo}[Cholesky 平滑器：块前向 $+$ 块后向逐式]\label{alg:est-cholesky-smoother}
对线性高斯模型 \cref{eq:lg-models}，按 \cref{eq:cholesky-init} 初始化 $\{\boldsymbol{\Lambda}_0,\mathbf{q}_0\}$，执行：\\[2pt]
\textbf{前向 pass}（$k=1\dots K$，同时算出 $\mathbf{L}$ 与 $\mathbf{d}=\mathbf{L}^{-1}\mathbf{H}^\top\mathbf{W}^{-1}\mathbf{z}$ 的块）：
\begin{subequations}\label{eq:cholesky-forward}
\begin{align}
  \mathbf{L}_{k-1}\mathbf{L}_{k-1}^\top&=\boldsymbol{\Lambda}_{k-1}+\mathbf{A}_{k-1}^\top\boldsymbol{\Sigma}_{\mathbf{w},k}^{-1}\mathbf{A}_{k-1}
   &&\text{（对左端做稠密 Cholesky 得 }\mathbf{L}_{k-1}\text{）},\label{eq:chol-fwd-Lk}\\
  \mathbf{L}_{k-1}\mathbf{d}_{k-1}&=\mathbf{q}_{k-1}-\mathbf{A}_{k-1}^\top\boldsymbol{\Sigma}_{\mathbf{w},k}^{-1}\mathbf{v}_k
   &&\text{（前代解出 }\mathbf{d}_{k-1}\text{）},\label{eq:chol-fwd-dk}\\
  \mathbf{L}_{k,k-1}\mathbf{L}_{k-1}^\top&=-\boldsymbol{\Sigma}_{\mathbf{w},k}^{-1}\mathbf{A}_{k-1}
   &&\text{（解出次对角块 }\mathbf{L}_{k,k-1}\text{）},\label{eq:chol-fwd-Lsub}\\
  \boldsymbol{\Lambda}_k&=-\mathbf{L}_{k,k-1}\mathbf{L}_{k,k-1}^\top+\boldsymbol{\Sigma}_{\mathbf{w},k}^{-1}+\mathbf{C}_k^\top\boldsymbol{\Sigma}_{\mathbf{v},k}^{-1}\mathbf{C}_k,\label{eq:chol-fwd-info}\\
  \mathbf{q}_k&=-\mathbf{L}_{k,k-1}\mathbf{d}_{k-1}+\boldsymbol{\Sigma}_{\mathbf{w},k}^{-1}\mathbf{v}_k+\mathbf{C}_k^\top\boldsymbol{\Sigma}_{\mathbf{v},k}^{-1}\mathbf{y}_k.\label{eq:chol-fwd-q}
\end{align}
\end{subequations}
末步取 $\hat{\mathbf{x}}_K=\boldsymbol{\Lambda}_K^{-1}\mathbf{q}_K$（即 $\mathbf{L}_K\mathbf{L}_K^\top=\boldsymbol{\Lambda}_K$、$\mathbf{L}_K\mathbf{d}_K=\mathbf{q}_K$、$\mathbf{L}_K^\top\hat{\mathbf{x}}_K=\mathbf{d}_K$）。\\[2pt]
\textbf{后向 pass}（$k=K\dots 1$，回代解 $\mathbf{L}^\top\hat{\mathbf{x}}=\mathbf{d}$）：
\begin{equation}
  \mathbf{L}_{k-1}^\top\,\hat{\mathbf{x}}_{k-1}=-\mathbf{L}_{k,k-1}^\top\,\hat{\mathbf{x}}_k+\mathbf{d}_{k-1}.
  \label{eq:cholesky-backward}
\end{equation}
全程仅需 Cholesky 分解、矩阵乘加、前/后代解三角系统，总复杂度 $O(N^3(K+1))$。
\end{algo}

\noindent\textbf{逐式从何而来（块比对）。}\;以前几块为例：$\mathbf{H}^\top\mathbf{W}^{-1}\mathbf{H}$ 的第 $(0,0)$ 块是 $\check{\boldsymbol{\Sigma}}_0^{-1}+\mathbf{C}_0^\top\boldsymbol{\Sigma}_{\mathbf{v},0}^{-1}\mathbf{C}_0+\mathbf{A}_0^\top\boldsymbol{\Sigma}_{\mathbf{w},1}^{-1}\mathbf{A}_0$（来自先验、$\mathbf{y}_0$ 观测、$\mathbf{x}_1$ 的运动项三处），它须等于 $(\mathbf{L}\mathbf{L}^\top)_{00}=\mathbf{L}_0\mathbf{L}_0^\top$，故 $\mathbf{L}_0\mathbf{L}_0^\top=\boldsymbol{\Lambda}_0+\mathbf{A}_0^\top\boldsymbol{\Sigma}_{\mathbf{w},1}^{-1}\mathbf{A}_0$——即 \cref{eq:chol-fwd-Lk} 在 $k=1$ 的形态，对它做一次\emph{小而稠密}的 Cholesky 即得 $\mathbf{L}_0$。第 $(1,0)$ 块 $-\boldsymbol{\Sigma}_{\mathbf{w},1}^{-1}\mathbf{A}_0$ 须等于 $\mathbf{L}_{1,0}\mathbf{L}_0^\top$，给出 \cref{eq:chol-fwd-Lsub}；再代回第 $(1,1)$ 块得 $\mathbf{L}_1\mathbf{L}_1^\top=\boldsymbol{\Lambda}_1+\mathbf{A}_1^\top\boldsymbol{\Sigma}_{\mathbf{w},2}^{-1}\mathbf{A}_1$，其中 $\boldsymbol{\Lambda}_1=-\mathbf{L}_{1,0}\mathbf{L}_{1,0}^\top+\boldsymbol{\Sigma}_{\mathbf{w},1}^{-1}+\mathbf{C}_1^\top\boldsymbol{\Sigma}_{\mathbf{v},1}^{-1}\mathbf{C}_1$，正是 \cref{eq:chol-fwd-info}。$\mathbf{d}$ 的递推同理来自 $\mathbf{L}\mathbf{d}=\mathbf{H}^\top\mathbf{W}^{-1}\mathbf{z}$ 的块比对（$\mathbf{H}^\top\mathbf{W}^{-1}\mathbf{z}$ 的第 $k$ 块是 $\check{\boldsymbol{\Sigma}}_0^{-1}\check{\mathbf{x}}_0/\,\boldsymbol{\Sigma}_{\mathbf{w},k}^{-1}\mathbf{v}_k$ 类信息向量项，减去耦合到 $\mathbf{x}_{k+1}$ 的 $\mathbf{A}_k^\top\boldsymbol{\Sigma}_{\mathbf{w},k+1}^{-1}\mathbf{v}_{k+1}$）。如此从 $k=0$ 一路解到 $K$，单次前向 pass 算全 $\mathbf{L},\mathbf{d}$ 块；再单次后向回代 \cref{eq:cholesky-backward} 算全 $\hat{\mathbf{x}}_k$ 块——这就完整复现了 \cref{thm:normal-eq} 的批量解 $\hat{\mathbf{x}}=(\mathbf{H}^\top\mathbf{W}^{-1}\mathbf{H})^{-1}\mathbf{H}^\top\mathbf{W}^{-1}\mathbf{z}$，无任何近似。

\begin{insight}[前向 pass $=$ 卡尔曼滤波，前向 $+$ 后向 $=$ RTS 平滑]
\cref{eq:cholesky-forward} 的五式映射 $\{\mathbf{q}_{k-1},\boldsymbol{\Lambda}_{k-1}\}\mapsto\{\mathbf{q}_k,\boldsymbol{\Lambda}_k\}$，再加 \cref{eq:cholesky-backward} 一条后向式，\emph{逐字等价}于经典 RTS 平滑器（\cref{ch:eskf} 的 \cref{eq:rts}、\cref{eq:rts-cov}）；其中\emph{五条前向式单独}就等价于卡尔曼滤波。等价的桥梁是 SMW 引理 \cref{eq:smw}：令前向滤波协方差 $\hat{\boldsymbol{\Sigma}}_{k,f}:=\boldsymbol{\Lambda}_k^{-1}$，对 \cref{eq:chol-fwd-Lk}、\cref{eq:chol-fwd-Lsub} 代入 \cref{eq:chol-fwd-info} 并用 SMW，得
\[
  \boldsymbol{\Lambda}_k=\big(\mathbf{A}_{k-1}\boldsymbol{\Lambda}_{k-1}^{-1}\mathbf{A}_{k-1}^\top+\boldsymbol{\Sigma}_{\mathbf{w},k}\big)^{-1}+\mathbf{C}_k^\top\boldsymbol{\Sigma}_{\mathbf{v},k}^{-1}\mathbf{C}_k
  \;\Longleftrightarrow\;
  \begin{cases}\check{\boldsymbol{\Sigma}}_{k,f}=\mathbf{A}_{k-1}\hat{\boldsymbol{\Sigma}}_{k-1,f}\mathbf{A}_{k-1}^\top+\boldsymbol{\Sigma}_{\mathbf{w},k},\\ \hat{\boldsymbol{\Sigma}}_{k,f}^{-1}=\check{\boldsymbol{\Sigma}}_{k,f}^{-1}+\mathbf{C}_k^\top\boldsymbol{\Sigma}_{\mathbf{v},k}^{-1}\mathbf{C}_k,\end{cases}
\]
即 KF 的预测协方差与信息形式校正（与 \cref{eq:kf-info-form} 同）；对 \cref{eq:chol-fwd-dk}、\cref{eq:chol-fwd-q} 同样代入并用 SMW，$\mathbf{q}_k$ 的递推化为 $\check{\mathbf{x}}_{k,f}=\mathbf{A}_{k-1}\hat{\mathbf{x}}_{k-1,f}+\mathbf{v}_k$ 与信息形式均值校正——正是 KF 预测–校正。后向回代 \cref{eq:cholesky-backward} 则化为 RTS 的后向修正 $\hat{\mathbf{x}}_{k-1}=\hat{\mathbf{x}}_{k-1,f}+\mathbf{G}_{k-1}(\hat{\mathbf{x}}_k-\check{\mathbf{x}}_{k,f})$（增益 $\mathbf{G}_{k-1}=\hat{\boldsymbol{\Sigma}}_{k-1,f}\mathbf{A}_{k-1}^\top\check{\boldsymbol{\Sigma}}_{k,f}^{-1}$）。一言以蔽之：\emph{Cholesky 是“矩阵分解”的语言、RTS 是“增益递推”的语言，二者是同一稀疏批量解的两种写法}（完整代数化简见 \cref{ch:eskf}）。
\end{insight}

\paragraph{从 Cholesky 因子直接回收边际协方差。}
\cref{alg:est-cholesky-smoother} 给出了均值 $\hat{\mathbf{x}}$，却没顺手给\emph{每个时刻的后验边际协方差} $\hat{\boldsymbol{\Sigma}}_k$。天真做法是把 $\hat{\boldsymbol{\Sigma}}=\boldsymbol{\Lambda}^{-1}=\mathbf{L}^{-\top}\mathbf{L}^{-1}$ 整块求逆再抽对角块——但 $\mathbf{L}^{-1}$ 的下三角\emph{稠密}（$\mathbf{L}^{-1}$ 第 $(k,j)$ 块含 $\mathbf{L}_k^{-1}\mathbf{L}_{k,k-1}\cdots\mathbf{L}_{j+1,j}\mathbf{L}_j^{-1}$ 这类连乘），故 $\hat{\boldsymbol{\Sigma}}$ 一般\emph{满阵}，$O((NK)^2)$ 存不下、更别说算。所幸我们\emph{只要主对角块（及紧邻一条次对角块）}，它们可由一条 $O(K)$ 的\emph{后向递推}从已算好的 $\mathbf{L}$ 块直接得到\cite{barfoot2024state}。

\begin{algo}[平滑器协方差回收（从 $\mathbf{L}$ 后向递推）]\label{alg:est-cov-recovery}
承 \cref{alg:est-cholesky-smoother} 已得全部 $\{\mathbf{L}_k,\mathbf{L}_{k,k-1}\}$。要后验协方差 $\hat{\boldsymbol{\Sigma}}$ 的主对角块 $\hat{\boldsymbol{\Sigma}}_k$ 与次对角块 $\hat{\boldsymbol{\Sigma}}_{k,k-1}$，自末端
\begin{equation}
  \hat{\boldsymbol{\Sigma}}_K=\mathbf{L}_K^{-\top}\mathbf{L}_K^{-1}
  \label{eq:cov-recovery-init}
\end{equation}
起，按 $k=K\dots 1$ 后向递推
\begin{subequations}\label{eq:cov-recovery}
\begin{align}
  \hat{\boldsymbol{\Sigma}}_{k,k-1}&=-\hat{\boldsymbol{\Sigma}}_k\,\mathbf{L}_{k,k-1}\mathbf{L}_{k-1}^{-1},\label{eq:cov-recovery-cross}\\
  \hat{\boldsymbol{\Sigma}}_{k-1}&=\mathbf{L}_{k-1}^{-\top}\big(\mathbf{I}+\mathbf{L}_{k,k-1}^\top\,\hat{\boldsymbol{\Sigma}}_k\,\mathbf{L}_{k,k-1}\big)\mathbf{L}_{k-1}^{-1}.\label{eq:cov-recovery-diag}
\end{align}
\end{subequations}
此过程可\emph{并入} \cref{alg:est-cholesky-smoother} 的后向 pass，不改变 $O(K)$ 总复杂度（仅常数略增）。
\end{algo}

\noindent\textbf{递推来历。}\;由 $\hat{\boldsymbol{\Sigma}}=\mathbf{L}^{-\top}\mathbf{L}^{-1}$ 逐块比对：$\mathbf{L}^{-1}$ 仍块下三角，其对角块为 $\mathbf{L}_k^{-1}$、紧邻次对角块为 $-\mathbf{L}_k^{-1}\mathbf{L}_{k,k-1}\mathbf{L}_{k-1}^{-1}$。把 $\hat{\boldsymbol{\Sigma}}=\mathbf{L}^{-\top}\mathbf{L}^{-1}$ 的第 $(k\!-\!1,k\!-\!1)$ 块展开，只有 $\mathbf{L}^{-1}$ 第 $k\!-\!1$ 列的两块（$\mathbf{L}_{k-1}^{-1}$ 与来自第 $k$ 行的 $-\mathbf{L}_k^{-1}\mathbf{L}_{k,k-1}\mathbf{L}_{k-1}^{-1}$，后者经 $\mathbf{L}^{-\top}$ 配对回 $\hat{\boldsymbol{\Sigma}}_k$）贡献，整理即 \cref{eq:cov-recovery-diag}；第 $(k,k\!-\!1)$ 块同理给 \cref{eq:cov-recovery-cross}。\emph{这正是稀疏矩阵“选择求逆”（Takahashi 递推）的块特化}：只回收因子非零模式所覆盖的那几条对角，避开稠密满阵。次对角块 $\hat{\boldsymbol{\Sigma}}_{k,k-1}$（相邻时刻的互协方差）在连续时间插值（\cref{sec:est-steam-query}）、平滑后一致性检验中要用。把 \cref{eq:cov-recovery} 用 SMW 进一步代入前向滤波量，即化为经典 RTS 后向协方差 \cref{eq:rts-cov} $\hat{\boldsymbol{\Sigma}}_{k-1}=\hat{\boldsymbol{\Sigma}}_{k-1,f}+\mathbf{G}_{k-1}(\hat{\boldsymbol{\Sigma}}_k-\check{\boldsymbol{\Sigma}}_{k,f})\mathbf{G}_{k-1}^\top$（\cref{ch:eskf}），再次印证两种语言同源。

\subsection{卡尔曼滤波 \texorpdfstring{$=$}{=} 批量 MAP 的前向递归}
\label{sec:est-kf-from-map}

\begin{theorem}[KF 是批量 MAP 的前向递归特例]\label{thm:kf-from-map}
对线性高斯模型 \cref{eq:lg-models}，把批量 MAP 按时间因式分解、对前向部分逐时刻边缘化历史，得到的递归恰是卡尔曼滤波的预测–校正两步；其结果与批量平滑在“仅用 $\le k$ 数据”意义下完全一致。
\end{theorem}
\begin{proof}[证明思路（完整递推留 \cref{ch:eskf}）]
设已有 $k-1$ 的前向估计 $\{\hat{\mathbf{x}}_{k-1},\hat{\boldsymbol{\Sigma}}_{k-1}\}$（基于 $\le k-1$ 全部数据）。要并入新数据 $\{\mathbf{v}_k,\mathbf{y}_k\}$ 得 $\{\hat{\mathbf{x}}_k,\hat{\boldsymbol{\Sigma}}_k\}$。\emph{用马尔可夫性}，把 $\{\hat{\mathbf{x}}_{k-1},\hat{\boldsymbol{\Sigma}}_{k-1}\}$ 作为 $\le k-1$ 全部数据的等价代表，构一个三块小 MAP：先验块 $\mathcal{N}(\hat{\mathbf{x}}_{k-1},\hat{\boldsymbol{\Sigma}}_{k-1})$、运动块（连 $\mathbf{x}_{k-1},\mathbf{x}_k$）、观测块（连 $\mathbf{x}_k$）。解此小 MAP 并把不再关心的 $\mathbf{x}_{k-1}$ \emph{边缘化}（信息形式下即对 $\mathbf{x}_{k-1}$ 块做舒尔补，\cref{eq:info-marginal}），用 SMW \cref{eq:smw} 化简，得
\begin{align}
  &\check{\boldsymbol{\Sigma}}_k=\mathbf{A}_{k-1}\hat{\boldsymbol{\Sigma}}_{k-1}\mathbf{A}_{k-1}^\top+\boldsymbol{\Sigma}_{\mathbf{w},k},\quad
  \check{\mathbf{x}}_k=\mathbf{A}_{k-1}\hat{\mathbf{x}}_{k-1}+\mathbf{v}_k&&\text{（预测步）},\notag\\
  &\hat{\boldsymbol{\Sigma}}_k^{-1}=\check{\boldsymbol{\Sigma}}_k^{-1}+\mathbf{C}_k^\top\boldsymbol{\Sigma}_{\mathbf{v},k}^{-1}\mathbf{C}_k,\quad
  \hat{\boldsymbol{\Sigma}}_k^{-1}\hat{\mathbf{x}}_k=\check{\boldsymbol{\Sigma}}_k^{-1}\check{\mathbf{x}}_k+\mathbf{C}_k^\top\boldsymbol{\Sigma}_{\mathbf{v},k}^{-1}\mathbf{y}_k&&\text{（校正步，信息形式）}.\label{eq:kf-info-form}
\end{align}
定义卡尔曼增益 $\mathbf{K}_k=\check{\boldsymbol{\Sigma}}_k\mathbf{C}_k^\top(\mathbf{C}_k\check{\boldsymbol{\Sigma}}_k\mathbf{C}_k^\top+\boldsymbol{\Sigma}_{\mathbf{v},k})^{-1}$（由 SMW 把信息形式 \cref{eq:kf-info-form} 转协方差形式），校正步即 $\hat{\mathbf{x}}_k=\check{\mathbf{x}}_k+\mathbf{K}_k(\mathbf{y}_k-\mathbf{C}_k\check{\mathbf{x}}_k)$、$\hat{\boldsymbol{\Sigma}}_k=(\mathbf{I}-\mathbf{K}_k\mathbf{C}_k)\check{\boldsymbol{\Sigma}}_k$。这正是卡尔曼滤波。\emph{新息}（innovation）$\mathbf{y}_k-\mathbf{C}_k\check{\mathbf{x}}_k$ 是实际与预期测量之差，$\mathbf{K}_k$ 权衡它相对预测的贡献。等价性：前向递归在每个 $k$ 给出“仅用 $\le k$ 数据的最优估计”，与批量平滑的前向 pass 逐字相同（去掉平滑后向 pass 即得因果的滤波）。
\end{proof}

\begin{insight}[滤波 $=$ 边缘化历史]
卡尔曼滤波“看起来”比批量省，是因为它在每一步\emph{把历史状态边缘化掉}、只保留当前状态的高斯摘要 $\{\hat{\mathbf{x}}_k,\hat{\boldsymbol{\Sigma}}_k\}$。这正是 \cref{thm:kf-from-map} 里那一步“对 $\mathbf{x}_{k-1}$ 做舒尔补”。代价是：边缘化在信息形式下是舒尔补，会在被消变量的邻居之间引入新边——若状态里有路标，边缘化一个位姿会让它见过的所有路标\emph{两两耦合}，使信息矩阵\emph{变稠密}。这就是下节“平滑稀疏、滤波稠密”的根，也是现代 SLAM 后端多走平滑/因子图的原因。
\end{insight}

\paragraph{增益优化视角与 BLUE。}
卡尔曼增益还有第三种来历\cite{barfoot2024state}：设校正 $\hat{\mathbf{x}}_k=\check{\mathbf{x}}_k+\mathbf{K}_k(\mathbf{y}_k-\mathbf{C}_k\check{\mathbf{x}}_k)$、增益 $\mathbf{K}_k$ 待定，估计误差协方差取 Joseph 形 $\mathbb{E}[\hat{\mathbf{e}}_k\hat{\mathbf{e}}_k^\top]=(\mathbf{I}-\mathbf{K}_k\mathbf{C}_k)\check{\boldsymbol{\Sigma}}_k(\mathbf{I}-\mathbf{K}_k\mathbf{C}_k)^\top+\mathbf{K}_k\boldsymbol{\Sigma}_{\mathbf{v},k}\mathbf{K}_k^\top$，最小化其迹 $J(\mathbf{K}_k)=\tfrac12\mathrm{tr}\,\mathbb{E}[\hat{\mathbf{e}}_k\hat{\mathbf{e}}_k^\top]$（即 MMSE）。借两条\emph{迹–导数恒等式}（\cref{ch:matderiv}）
\begin{equation}
  \frac{\partial\,\mathrm{tr}(\mathbf{X}\mathbf{Y})}{\partial\mathbf{X}}=\mathbf{Y}^\top,\qquad
  \frac{\partial\,\mathrm{tr}(\mathbf{X}\mathbf{Z}\mathbf{X}^\top)}{\partial\mathbf{X}}=2\mathbf{X}\mathbf{Z}\ \ (\mathbf{Z}=\mathbf{Z}^\top),
  \label{eq:trace-deriv}
\end{equation}
把 Joseph 形展开成 $\mathbf{K}_k$ 的二次型（$\check{\boldsymbol{\Sigma}}_k,\boldsymbol{\Sigma}_{\mathbf{v},k}$ 对称），对 $\mathbf{K}_k$ 求导置零
\begin{equation}
  \frac{\partial J(\mathbf{K}_k)}{\partial\mathbf{K}_k}=-(\mathbf{I}-\mathbf{K}_k\mathbf{C}_k)\check{\boldsymbol{\Sigma}}_k\mathbf{C}_k^\top+\mathbf{K}_k\boldsymbol{\Sigma}_{\mathbf{v},k}=\mathbf{0},
  \label{eq:gain-stationary}
\end{equation}
整理得 $\mathbf{K}_k(\mathbf{C}_k\check{\boldsymbol{\Sigma}}_k\mathbf{C}_k^\top+\boldsymbol{\Sigma}_{\mathbf{v},k})=\check{\boldsymbol{\Sigma}}_k\mathbf{C}_k^\top$，即同一 $\mathbf{K}_k=\check{\boldsymbol{\Sigma}}_k\mathbf{C}_k^\top(\mathbf{C}_k\check{\boldsymbol{\Sigma}}_k\mathbf{C}_k^\top+\boldsymbol{\Sigma}_{\mathbf{v},k})^{-1}$。\textbf{重要论断}：加权最小二乘/KF 在 MMSE 意义下最优，只要测量误差独立、零均值、方差已知——\emph{误差不必高斯}，且无其他无偏估计能更好（BLUE，高斯–马尔可夫定理，见 \cref{sec:est-ls}）。

\begin{table}[H]\centering\small
\caption{批量平滑 vs 递归滤波（本章给等价，完整 KF 见 \cref{ch:eskf}）}
\label{tab:est-batch-vs-recursive}
\begin{tabular}{lll}
\toprule
维度 & 批量平滑（smoothing） & 递归滤波（filtering） \\
\midrule
保留变量 & 整条轨迹（全部位姿/路标） & 仅当前状态（边缘化历史） \\
因果性 & 非因果（用全部数据） & 因果（仅用 $\le k$ 数据） \\
信息矩阵 & \textbf{稀疏且保持稀疏}（块三对角/箭头） & \textbf{随时间变稠密}（舒尔补引入新边） \\
求解 & 稀疏 Cholesky/QR，$O(K)$ & 预测–校正递归，每步 $O(N^3)$ \\
最优性 & 完整 MAP（无近似） & 等于批量的前向 pass（线性高斯下无近似） \\
适用 & 事后处理、回环、可扩展到百万变量 & 在线、实时、内存恒定 \\
\bottomrule
\end{tabular}
\end{table}

\paragraph{前沿一瞥：连续时间与高斯过程回归。}
若状态是时间的连续函数 $\mathbf{x}(t)\sim\mathcal{GP}(\check{\mathbf{x}}(t),\check{\boldsymbol{\Sigma}}(t,t'))$，把估计看作\emph{高斯过程回归}：用线性时不变运动模型（如白噪声加速度“常速度”模型 $\ddot{\mathbf{p}}=\mathbf{w}$）构造先验，其离散化在测量时刻给出与批量离散方法\emph{完全相同}的块三对角系统，故同样可用稀疏 Cholesky/RTS 高效解；查询任意非测量时刻 $\tau$ 的状态是 $O(1)$ 插值（常速度模型下恰为三次 Hermite 多项式插值）\cite{barfoot2024state}。这把“批量/递归/连续时间”统一为“在不同先验下做 GP 回归”，是连续时间 SLAM（STEAM）的基础。本章 \cref{sec:est-steam} 把这一思想在 $\SE(3)$ 上\emph{完整展开}（局部 GP 缝合、白噪声加速度先验、$O(1)$ 插值/外推与协方差恢复）。

\begin{pitfall}[思维陷阱：以为“滤波比平滑省所以更优”]
\textbf{错误想法}：“滤波每步只算当前状态，比批量平滑省，所以更好。”\textbf{现象}：在含路标/回环的 SLAM 里用滤波（如 EKF-SLAM），信息矩阵随时间\emph{变稠密}（\cref{tab:est-batch-vs-recursive}），单步代价从 $O(N^3)$ 涨到与地图规模平方乃至立方相关，反而比稀疏平滑慢且不可扩展；且滤波的一次线性化\emph{不可回溯}，早期误差被永久固化。\textbf{根本原因}：滤波 $=$ 边缘化历史，边缘化在信息形式下是舒尔补，会引入新边使矩阵稠密（\cref{thm:kf-from-map} 后的洞见）。\textbf{正确做法}：后端用平滑/因子图保持稀疏（\cref{sec:est-fg}），并按需重新线性化；纯递推、低维、强实时（如 IMU 姿态）才用滤波。\textbf{自检}：观察信息矩阵非零元随时间的增长——若快速稠密化，说明该用平滑。
\end{pitfall}

\begin{practice}
\begin{enumerate}
  \item 证明 $\mathbf{A}^{-1}$（lifted 转移矩阵之逆）仅主对角与下一对角非零（提示：$\mathbf{A}$ 是单位下三角，逆仍下三角，且 $\mathbf{x}_k$ 只依赖 $\mathbf{x}_{k-1}$）。由此说明 \cref{eq:tridiag} 的块三对角。
  \item 从 \cref{eq:kf-info-form} 出发，用 SMW \cref{eq:smw} 推出协方差形式的卡尔曼增益与 $\hat{\boldsymbol{\Sigma}}_k=(\mathbf{I}-\mathbf{K}_k\mathbf{C}_k)\check{\boldsymbol{\Sigma}}_k$。
  \item （概念）在一个含 1 个路标、$K$ 个位姿的滤波里，边缘化第一个位姿后，剩余信息矩阵的稀疏结构如何变化？为何这预示 EKF-SLAM 的稠密化？
  \item （跨章综合）把 \cref{ch:eskf} 的 KF 预测–校正与本节 \cref{eq:kf-info-form} 对照，指出二者是同一组方程的协方差形式与信息形式。
\end{enumerate}
\end{practice}

% ════════════════════════════════════════════════════════════════
\section{因子图：现代统一视角}
\label{sec:est-fg}

前面把 SLAM 化为最小二乘，并看到稀疏来自概率图的局部结构。\emph{因子图}给这套最小二乘一个图形语言，它是当代 SLAM 后端（GTSAM、g2o、Ceres）的通用范式与“普通话”\cite{dellaert2017factor,carlone2026handbook}。本节沿“因子图 $\to$ MAP $\to$ 线性化 $\to$ 稀疏分解 $\to$ 变量消元 $\to$ 增量平滑”这条链条，把前面零散建立的概率与稀疏结论收束成一个统一框架：因子图既是后端的\emph{建模语言}（声明变量与因子），又精确对应一个可稀疏求解的最小二乘问题，其变量消元就是稀疏矩阵分解的图论侧写。

\subsection{从局部性到因子图}
\label{sec:est-fg-def}

\paragraph{局部性 $\Rightarrow$ 后验可因式分解。}
\cref{eq:factorize} 已表明：因每个测量只涉及少数变量（\emph{局部性}），高维后验可写成许多\emph{小定义域}因子之积。以一个玩具例（机器人三位姿 $\mathbf{p}_1,\mathbf{p}_2,\mathbf{p}_3$、两路标 $\boldsymbol{\ell}_1,\boldsymbol{\ell}_2$、对 $\mathbf{p}_1$ 一个绝对测量、对路标做 bearing 观测）为例\cite{carlone2026handbook}，贝叶斯展开后验
\begin{equation}
  p(\mathbf{x}\mid\mathbf{z})\propto
  \underbrace{p(\mathbf{p}_1)p(\mathbf{p}_2\mid\mathbf{p}_1)p(\mathbf{p}_3\mid\mathbf{p}_2)}_{\text{运动链（马尔可夫）}}
  \underbrace{p(\boldsymbol{\ell}_1)p(\boldsymbol{\ell}_2)}_{\text{路标先验}}
  \underbrace{p(\mathbf{z}_1\mid\mathbf{p}_1)}_{\text{绝对测量}}
  \underbrace{p(\mathbf{z}_2\mid\mathbf{p}_1,\boldsymbol{\ell}_1)p(\mathbf{z}_3\mid\mathbf{p}_2,\boldsymbol{\ell}_1)p(\mathbf{z}_4\mid\mathbf{p}_3,\boldsymbol{\ell}_2)}_{\text{bearing 观测}}.
  \label{eq:toy-factorize}
\end{equation}
因子图就是把这种分解\emph{画出来}的工具：每个因子一个节点，只连它依赖的变量。

\begin{definition}[因子图]\label{def:factor-graph}
因子图是一个\emph{二部图} $F=(\mathcal{U},\mathcal{V},\mathcal{E})$：\emph{变量节点} $x_j\in\mathcal{V}$（待估的位姿、路标，画作圆），\emph{因子节点} $\phi_i\in\mathcal{U}$（每条信息一个，画作方块/小黑点），边 $e_{ij}\in\mathcal{E}$ 只连因子与其涉及的变量。记 $\mathcal{X}(\phi_i)$ 为因子 $\phi_i$ 邻接的变量集、$\mathbf{x}_i$ 为对该集的一次赋值。整图定义全局函数（未归一化后验）的分解
\begin{equation}
  \phi(\mathbf{x})=\prod_i\phi_i(\mathbf{x}_i)\ \propto\ p(\mathbf{z}\mid\mathbf{x})\,p(\mathbf{x}).
  \label{eq:fg-product}
\end{equation}
测量\emph{不画}成节点——它们已知，被吸收进因子（如观测因子 $\phi_i(\mathbf{x}_k,\mathbf{y}_j)\propto p(\mathbf{z}_{k,j}\mid\mathbf{x}_k,\mathbf{y}_j)$）。因子值只需\emph{正比于}对应密度，不依赖状态的归一化常数可省。
\end{definition}

\begin{figure}[ht]
\centering
\begin{tikzpicture}[
  >=Stealth, node distance=14mm,
  var/.style={circle, draw=main!70!black, fill=main!8, minimum size=8mm, font=\small, inner sep=0pt},
  lmk/.style={circle, draw=third!70!black, fill=third!10, minimum size=8mm, font=\small, inner sep=0pt},
  fac/.style={rectangle, draw=second!60!black, fill=second!60!black, minimum size=2.1mm, inner sep=1.6pt},
  ed/.style={gray!55!black}]
  \node[var] (p1) {$\mathbf{x}_1$};
  \node[var, right=of p1] (p2) {$\mathbf{x}_2$};
  \node[var, right=of p2] (p3) {$\mathbf{x}_3$};
  \node[lmk, above=11mm of p1] (l1) {$\mathbf{y}_1$};
  \node[lmk, above=11mm of p3] (l2) {$\mathbf{y}_2$};
  \node[fac, left=8mm of p1] (fp) {}; \draw[ed] (fp)--(p1);
  \node[fac] (o1) at ($(p1)!0.5!(p2)$) {}; \draw[ed] (p1)--(o1) (o1)--(p2);
  \node[fac] (o2) at ($(p2)!0.5!(p3)$) {}; \draw[ed] (p2)--(o2) (o2)--(p3);
  \node[fac] (m1) at ($(p1)!0.5!(l1)$) {}; \draw[ed] (p1)--(m1) (m1)--(l1);
  \node[fac] (m2) at ($(p2)!0.45!(l1)$) {}; \draw[ed] (p2)--(m2) (m2)--(l1);
  \node[fac] (m3) at ($(p3)!0.5!(l2)$) {}; \draw[ed] (p3)--(m3) (m3)--(l2);
\end{tikzpicture}
\caption{一个小型 SLAM 因子图：圆为变量（位姿 $\mathbf{x}_k$、路标 $\mathbf{y}_j$），方块为因子（左端孤立方块 $=$ 先验，位姿间 $=$ 里程计/运动，位姿–路标间 $=$ 观测）。测量不出现为节点。（仿 \cite{dellaert2017factor,carlone2026handbook} 重绘）}
\label{fig:factor-graph}
\end{figure}

\begin{theorem}[因子图上的 MAP $=$ 非线性最小二乘]\label{thm:fg-map}
设每个因子正比于高斯 $\phi_i(\mathbf{x}_i)\propto\exp\!\big(-\tfrac12\|\mathbf{z}_i-h_i(\mathbf{x}_i)\|_{\boldsymbol{\Sigma}_i}^2\big)$，则
\begin{equation}
  \mathbf{x}^\star_{\mathrm{MAP}}
  =\arg\max_{\mathbf{x}}\prod_i\phi_i(\mathbf{x}_i)
  =\arg\min_{\mathbf{x}}\sum_i\|\mathbf{z}_i-h_i(\mathbf{x}_i)\|_{\boldsymbol{\Sigma}_i}^2.
  \label{eq:fg-ls}
\end{equation}
\end{theorem}
\begin{proof}
对 \cref{eq:fg-product} 取负对数：乘积变求和，每个高斯因子贡献半个平方马氏距离（\cref{eq:neglog-gauss}），常数丢弃，$\arg\max\prod=\arg\min\sum$。这是 \cref{thm:map-ls} 的图论版。
\end{proof}
\noindent 最小化 \cref{eq:fg-ls} 即\emph{传感器融合}：把若干观测因子（与先验）拢到一起定出 MAP 解。因子图作为一种\emph{语言}，能可视化各种 SLAM 变体（\cref{tab:est-slam-variants}），给出结构洞见，并充当后端的通用接口。

\begin{table}[H]\centering\small
\caption{因子图统一的 SLAM 变体}
\label{tab:est-slam-variants}
\begin{tabular}{lll}
\toprule
变体 & 变量 & 与玩具例的区别 \\
\midrule
Landmark-based SLAM & 位姿 $+$ 路标 & 标准型（\cref{fig:factor-graph}） \\
光束法平差 BA & 位姿 $+$ 路标 & \emph{无}运动先验/运动模型（纯视觉） \\
位姿图优化 PGO & 仅位姿 & \emph{无}路标，位姿间含回环约束 \\
STEAM（连续时间） & 高阶状态（含速度 $\boldsymbol{\varpi}$） & 用光滑连续时间运动先验（GP，见 \cref{sec:est-steam}） \\
\bottomrule
\end{tabular}
\end{table}

\subsection{SLAM 的各类因子}
\label{sec:est-factors}

\cref{tab:factors} 列出常见因子及其残差。它们都是 \cref{eq:fg-ls} 求和中的一项，结构一致、只是连接的变量与 $h_i$ 不同。

\begin{table}[H]\centering\small
\caption{SLAM 因子图中的常见因子（残差 $\mathbf{e}_i$，代价 $\|\mathbf{e}_i\|_{\boldsymbol{\Sigma}_i}^2$）}
\label{tab:factors}
\begin{tabular}{lll}
\toprule
因子 & 连接变量 & 残差 $\mathbf{e}_i$ \\
\midrule
先验（锚定 gauge） & 单个 $\mathbf{x}_1$ & $\mathbf{x}_1\ominus\check{\mathbf{x}}_1$ \\
里程计 & $\mathbf{x}_k,\mathbf{x}_{k+1}$ & $\mathbf{o}_k-(\mathbf{x}_{k+1}\ominus\mathbf{x}_k)$ \\
运动（含控制 $\mathbf{u}_k$） & $\mathbf{x}_k,\mathbf{x}_{k+1}$ & $\mathbf{x}_{k+1}-f(\mathbf{x}_k,\mathbf{u}_k)$ \\
观测 / 重投影 & $\mathbf{x}_k,\mathbf{y}_j$ & $\mathbf{z}_{k,j}-h(\mathbf{x}_k,\mathbf{y}_j)$ \\
回环 & $\mathbf{x}_a,\mathbf{x}_b$（非相邻） & $\mathbf{o}_{ab}-(\mathbf{x}_b\ominus\mathbf{x}_a)$ \\
\bottomrule
\end{tabular}
\end{table}

\noindent 其中 $\ominus$ 是位姿的流形相减（\cref{ch:lie} 的 $\boxminus$/右扰动）。回环因子与里程计同型，只是连接\emph{非相邻}位姿——这一条边往往把累积漂移“拉回”，是 SLAM 区别于纯里程计的关键。运动模型先验（条件高斯 $p(\mathbf{p}_{t+1}\mid\mathbf{p}_t,\mathbf{u}_t)\propto\exp(-\tfrac12\|\mathbf{p}_{t+1}-g(\mathbf{p}_t,\mathbf{u}_t)\|_{\boldsymbol{\Sigma}_{\mathbf{w}}}^2)$）与里程计因子（$\propto\exp(-\tfrac12\|\mathbf{o}_t-(\mathbf{p}_{t+1}-\mathbf{p}_t)\|_{\boldsymbol{\Sigma}}^2)$）可并存\cite{carlone2026handbook}。

\subsection{线性化、白化与稀疏分解}
\label{sec:est-fg-linear}

\paragraph{线性化与白化。}
$h_i$ 非线性，从初值 $\mathbf{x}_i^0$ 一阶展开 $h_i(\mathbf{x}_i^0+\boldsymbol{\delta}_i)\approx h_i(\mathbf{x}_i^0)+\mathbf{H}_i\boldsymbol{\delta}_i$（$\mathbf{H}_i=\partial h_i/\partial\mathbf{x}_i|_{\mathbf{x}_i^0}$ 为测量雅可比，$\boldsymbol{\delta}_i$ 为状态更新；位姿处用右扰动 $\boldsymbol{\delta}_i=\delta\boldsymbol{\xi}_i$），\cref{eq:fg-ls} 化为关于 $\boldsymbol{\delta}$ 的线性最小二乘
\begin{equation}
  \boldsymbol{\delta}^\star=\arg\min_{\boldsymbol{\delta}}\sum_i\big\|\big(\mathbf{z}_i-h_i(\mathbf{x}_i^0)\big)-\mathbf{H}_i\boldsymbol{\delta}_i\big\|_{\boldsymbol{\Sigma}_i}^2.
  \label{eq:fg-linearized}
\end{equation}
\emph{白化}：左乘 $\boldsymbol{\Sigma}_i^{-1/2}$ 消去协方差与单位\cite{carlone2026handbook}，令 $\mathbf{A}_i=\boldsymbol{\Sigma}_i^{-1/2}\mathbf{H}_i$、$\mathbf{b}_i=\boldsymbol{\Sigma}_i^{-1/2}(\mathbf{z}_i-h_i(\mathbf{x}_i^0))$，得标准最小二乘 $\boldsymbol{\delta}^\star=\arg\min_{\boldsymbol{\delta}}\|\mathbf{A}\boldsymbol{\delta}-\mathbf{b}\|_2^2$，$\mathbf{A}$ 是大而稀疏的块矩阵、其块结构\emph{镜像因子图}（每个因子对应一个块行、每个变量对应一个块列）。

\paragraph{正规方程与平方根（$\sqrt{}$SAM）。}
满秩时唯一解由正规方程 $(\mathbf{A}^\top\mathbf{A})\boldsymbol{\delta}^\star=\mathbf{A}^\top\mathbf{b}$ 给出。常做法是因式分解\emph{信息矩阵} $\boldsymbol{\Lambda}=\mathbf{A}^\top\mathbf{A}=\mathbf{R}^\top\mathbf{R}$（Cholesky，$\mathbf{R}$ 上三角，$n^3/3$ flops），再前代 $\mathbf{R}^\top\mathbf{y}=\mathbf{A}^\top\mathbf{b}$、回代 $\mathbf{R}\boldsymbol{\delta}^\star=\mathbf{y}$\cite{carlone2026handbook}。更稳的是\emph{不形成} $\boldsymbol{\Lambda}$，直接对 $\mathbf{A}$（连同 RHS）做 QR：$\mathbf{A}=\mathbf{Q}\big[\begin{smallmatrix}\mathbf{R}\\\mathbf{0}\end{smallmatrix}\big]$，$\mathbf{Q}^\top\mathbf{b}=\big[\begin{smallmatrix}\mathbf{d}\\\mathbf{e}\end{smallmatrix}\big]$，由 $\mathbf{Q}$ 正交得 $\|\mathbf{A}\boldsymbol{\delta}-\mathbf{b}\|_2^2=\|\mathbf{R}\boldsymbol{\delta}-\mathbf{d}\|_2^2+\|\mathbf{e}\|_2^2$，解 $\mathbf{R}\boldsymbol{\delta}^\star=\mathbf{d}$（$\|\mathbf{e}\|_2^2$ 即残差平方和）。QR 与 Cholesky 给出\emph{相同}的 $\mathbf{R}$（至多对角符号），因 $\mathbf{A}^\top\mathbf{A}=\mathbf{R}^\top\mathbf{Q}^\top\mathbf{Q}\mathbf{R}=\mathbf{R}^\top\mathbf{R}$。QR 更精确稳定（$2(m-n/3)n^2$，比 Cholesky 慢约 2 倍）。这族“把信息矩阵或雅可比分解为平方根”的方法称\emph{平方根 SAM}（$\sqrt{}$SAM）\cite{dellaert2017factor}。

\paragraph{稀疏性、fill-in 与变量排序。}
$\mathbf{A}$ 块稀疏 $\Rightarrow\boldsymbol{\Lambda}=\mathbf{A}^\top\mathbf{A}$ 稀疏；$\boldsymbol{\Lambda}$ 的稀疏模式定义一个\emph{无向图} $G$：两变量若曾在同一因子中共现则有边（$n$ 元因子在其变量间诱导一个团）。稀疏分解远快于稠密，但\emph{变量排序戏剧性地影响 fill-in}（分解因子里超出原稀疏模式的额外非零元）：任意排序给出\emph{相同}的 MAP 解，但 fill-in 不同。\emph{求最小 fill-in 排序是 NP-hard}，实用启发式是 \emph{COLAMD}（列近似最小度）\cite{carlone2026handbook}。一个实测对比：某模拟 SLAM 的 Cholesky 因子 $\mathbf{R}$ 在自然排序（位姿在前、路标在后）下有 9399 个非零元，COLAMD 重排后只剩 4168 个——\emph{解完全相同，速度大不同}。

\subsection{变量消元 \texorpdfstring{$=$}{=} 矩阵分解}
\label{sec:est-elimination}

\begin{theorem}[变量消元 $=$ 稀疏矩阵分解]\label{thm:elimination}
在线性测量 $+$ 高斯噪声下，把因子图\emph{逐变量消元}转成贝叶斯网的算法，\emph{等价于}对信息矩阵/雅可比做稀疏 Cholesky/QR 分解；二者是同一通用算法的特例。
\end{theorem}
\begin{proof}[过程（含线性高斯条件的显式形式）]
消去变量 $x_j$：(1) 移除所有与 $x_j$ 相邻的因子，乘成中间因子 $\psi(x_j,\mathbf{s}_j)$（$\mathbf{s}_j$ 是 $x_j$ 的\emph{分隔集}——消元后 $x_j$ 条件依赖的变量），即把相关 $\mathbf{A}_i$ 累积进更大块矩阵 $\bar{\mathbf{A}}_j$：$\psi(x_j,\mathbf{s}_j)=\exp(-\tfrac12\|\bar{\mathbf{A}}_j[x_j;\mathbf{s}_j]-\bar{\mathbf{b}}_j\|_2^2)$；(2) 对增广 $[\bar{\mathbf{A}}_j\mid\bar{\mathbf{b}}_j]$ 做\emph{部分 QR}，得
\[
  [\bar{\mathbf{A}}_j\mid\bar{\mathbf{b}}_j]=\mathbf{Q}\begin{bmatrix}\mathbf{R}_j&\mathbf{T}_j&\mathbf{d}_j\\&\tilde{\mathbf{A}}_\tau&\tilde{\mathbf{b}}_\tau\end{bmatrix},
\]
$\mathbf{Q}$ 正交不改范数，故 $\psi$ 因式分解为
\[
  \psi(x_j,\mathbf{s}_j)=\underbrace{\exp(-\tfrac12\|\mathbf{R}_j x_j+\mathbf{T}_j\mathbf{s}_j-\mathbf{d}_j\|_2^2)}_{p(x_j\mid\mathbf{s}_j)}\ \underbrace{\exp(-\tfrac12\|\tilde{\mathbf{A}}_\tau\mathbf{s}_j-\tilde{\mathbf{b}}_\tau\|_2^2)}_{\tau(\mathbf{s}_j)\ \text{新因子}}.
\]
第一项是被消变量上的\emph{线性高斯条件} $p(x_j\mid\mathbf{s}_j)$，第二项是分隔集上的新因子，返回约简因子图。所有变量消完即得贝叶斯网分解 $p(\mathbf{x})=\prod_j p(x_j\mid\mathbf{s}_j)$。\emph{接回高斯条件化}：条件 $p(x_j\mid\mathbf{s}_j)$ 的均值 $\boldsymbol{\mu}_j=\mathbf{R}_j^{-1}(\mathbf{d}_j-\mathbf{T}_j\mathbf{s}_j)$（线性依赖分隔集）、协方差 $\boldsymbol{\Sigma}_j=(\mathbf{R}_j^\top\mathbf{R}_j)^{-1}$——这正是 \cref{thm:schur-cond} 高斯条件化的平方根/消元侧写法。逐变量消元累积出的上三角 $\mathbf{R}$ 就是 Cholesky 因子，故消元 $=$ 分解。
\end{proof}
\noindent\textbf{回代求 MAP}：消元完成后按\emph{逆消元序}回代，每步 $x_j^\star=\mathbf{R}_j^{-1}(\mathbf{d}_j-\mathbf{T}_j\mathbf{s}_j^\star)$（分隔集 MAP 此时已知，即条件均值）。\emph{洞见}：与上三角矩阵关联的图模型是\emph{贝叶斯网}（有向无环，对应“上三角”）；正如因子图是稀疏雅可比的图形体现，贝叶斯网揭示 Cholesky 因子的稀疏结构。

\begin{derivation}[一个可手算的标量消元例：消去路标如何产生 fill-in]
\label{der:est-elim-numeric}
抽象的“部分 QR”落到具体数字才看得真切。取 \cref{fig:factor-graph} 的一个最小子图，把全部变量压成\emph{标量}、测量取\emph{线性}，于是部分 QR 退化为可口算的配方。设两位姿 $x_1,x_2$ 与一路标 $\ell$，三个因子：对 $x_1$ 的先验 $\phi_a$、$x_1$–$x_2$ 的里程计 $\phi_b$、以及\emph{两条}对 $\ell$ 的测量 $\phi_c$（从 $x_1$）、$\phi_d$（从 $x_2$）。白化后（除以各自标准差，\cref{eq:fg-linearized}）残差取
\[
  \phi_a:\ x_1-0,\quad
  \phi_b:\ x_2-x_1-u,\quad
  \phi_c:\ \ell-x_1-r_1,\quad
  \phi_d:\ \ell-x_2-r_2,
\]
为使算术干净，设四个因子\emph{等权}（白化后信息均为 $1$）。变量序取 $[\ell;\,x_1;\,x_2]$（\emph{路标在前}），白化雅可比与右端为
\[
  \mathbf{A}=\begin{bmatrix}0&1&0\\0&-1&1\\1&-1&0\\1&0&-1\end{bmatrix},\qquad
  \mathbf{b}=\begin{bmatrix}0\\u\\r_1\\r_2\end{bmatrix}
  \quad(\text{行 }=\phi_a,\phi_b,\phi_c,\phi_d).
\]
\textbf{消去 $\ell$（第一列）}：按 \cref{thm:elimination} 第 (1) 步，只取\emph{第一列非零}的行——即 $\phi_c,\phi_d$——拼成中间因子 $\psi(\ell,x_1,x_2)$，其增广块
\[
  [\bar{\mathbf{A}}\mid\bar{\mathbf{b}}]=\begin{bmatrix}1&-1&0&\mid&r_1\\1&0&-1&\mid&r_2\end{bmatrix}.
\]
对第一列做\emph{部分 QR}：长度 $\sqrt{1^2+1^2}=\sqrt2$，一步 Householder/Givens 把两行旋成“上三角行 $+$ 余项行”。直接配方更快：两行平方和按 $\ell$ 配方，
\[
  (\ell-x_1-r_1)^2+(\ell-x_2-r_2)^2
  =2\Big(\ell-\underbrace{\tfrac{(x_1+r_1)+(x_2+r_2)}{2}}_{=\,\mu_\ell(x_1,x_2)}\Big)^2
  +\tfrac12\big((x_2+r_2)-(x_1+r_1)\big)^2 .
\]
\emph{第一项}给被消变量上的条件 $p(\ell\mid x_1,x_2)$：$\mathbf{R}_\ell=\sqrt2$、条件均值 $\mu_\ell=\tfrac12(x_1+x_2+r_1+r_2)$、条件方差 $\sigma_\ell^2=\mathbf{R}_\ell^{-2}=\tfrac12$（\emph{两条独立测量把 $\ell$ 的信息翻倍、方差减半}，正是 \cref{prop:info-add} 的信息相加）。\emph{第二项}\textbf{不再含 $\ell$}，是返回约简因子图的\emph{新因子}
\[
  \tau(x_1,x_2)=\exp\!\Big(-\tfrac12\cdot\tfrac12\big((x_2-x_1)-(r_1-r_2)\big)^2\Big),
\]
它是一条\emph{原图没有}的 $x_1$–$x_2$ 直接约束——这就是\textbf{fill-in}：消去共同路标 $\ell$ 后，曾通过 $\ell$ 间接耦合的 $x_1,x_2$ 之间长出一条新边（信息 $\tfrac12$，相当于“以 $\ell$ 为中介测得的相对位姿”）。它与原里程计 $\phi_b$ 同型，故二者在约简图里合并：$x_1$–$x_2$ 的等效信息 $=1+\tfrac12=\tfrac32$。
继续按 $x_1,x_2$ 消元即得贝叶斯网分解 $p=p(\ell\mid x_1,x_2)\,p(x_1\mid x_2)\,p(x_2)$，与 \cref{fig:factor-graph} 玩具例的结构（练习 4）一致。\hfill$\square$
\end{derivation}

\begin{insight}[消元顺序决定 fill-in：为何路标常\emph{先}消]
\cref{der:est-elim-numeric} 里“消 $\ell$ 在 $x_1,x_2$ 间生出一条边”是一般规律的缩影：消去一个变量，会在它\emph{所有邻居之间}补满一个团（其分隔集上的稠密块）。路标通常被许多位姿观测，但\emph{两两位姿原本不直接相连}；若\emph{先消位姿}，每消一个位姿就会把它看到的所有路标连成团，fill-in 迅速膨胀。反之\emph{先消路标}——路标的邻居（看到它的几个位姿）数目小且本就沿轨迹相邻——补出的团很小。这正是 \cref{sec:est-fg-linear} 里 COLAMD“按最小度选下一个消元变量”的直觉，也是 \cref{ch:nlopt} BA 中“先 Schur 消路标、留下稠密但小的位姿约简系统”的同一件事：\emph{消元顺序不改 MAP 解，却决定稀疏分解的代价}。
\end{insight}

\subsection{平滑 vs 滤波；增量平滑（Bayes 树 / iSAM）}
\label{sec:est-smooth-filter}

\paragraph{平滑 vs 滤波：信息矩阵稀疏还是稠密。}
\cref{tab:factors} 的因子图保留\emph{整条轨迹}的全部变量，称\emph{平滑}（smoothing）；其信息矩阵 $\boldsymbol{\Lambda}$ \textbf{稀疏且保持稀疏}。\emph{滤波}（filtering）只保留当前状态、把历史\emph{边缘化}掉——边缘化在信息形式下是舒尔补（\cref{eq:info-marginal}），会在被消变量的邻居间引入新边，使 $\boldsymbol{\Lambda}$ 随时间\emph{变稠密}\cite{carlone2026handbook}（与 \cref{thm:kf-from-map} 后的洞见、\cref{tab:est-batch-vs-recursive} 一致）。这解释了为何现代后端多走平滑/因子图（稀疏可扩展到百万级变量），而把滤波视为“边缘化历史”的特例——\cref{ch:eskf} 详谈。历史上 Mariner 10 金星任务（1969）的平方根信息滤波（SRIF）就用平方根换取数值稳定\cite{carlone2026handbook}。

\paragraph{增量平滑：Bayes 树与 iSAM。}
在线 SLAM 每来一个新测量就要更新解。朴素地每次重做整套消元浪费大量计算。\emph{Bayes 树}\cite{carlone2026handbook} 把消元得到的（弦化）贝叶斯网的\emph{团}组织成一棵有向树：每个团 $c_k=\mathbf{f}_k:\mathbf{s}_k$（前沿变量 $\mathbf{f}_k$ $+$ 与父团的分隔集 $\mathbf{s}_k$），定义条件 $p(\mathbf{f}_k\mid\mathbf{s}_k)$，整树给出 $p(\mathbf{x})=\prod_k p(\mathbf{f}_k\mid\mathbf{s}_k)$。\emph{增量更新}对应\emph{编辑 Bayes 树}：新因子只影响“含其变量的团”到“根”之间的路径——把这部分转回因子图、加新因子、重新消元、再接回未受影响的子树。两条性质保证“只有树的顶部受影响”：(1) 信息只向上传到根（按逆消元序收集）；(2) 因子信息只在其第一个变量被消时进入。据此发展出机器人学最先进的增量非线性 MAP 方法 \emph{iSAM2}\cite{carlone2026handbook}——新测量常只影响树的一小部分，故高效（大回环除外），并配合\emph{受影响因子的选择性重新线性化}与 constrained COLAMD 排序处理非线性。iSAM1/iSAM2 已用于百万级变量问题，实现于 GTSAM 库。这属现代后端专题，完整算法见 \cref{ch:vio,ch:lidar_slam} 与原综述\cite{dellaert2017factor}，此处给出统一视角与去向。

\begin{pitfall}[概念误区：把测量画成因子图的节点]
\textbf{错误想法}：因子图里既有变量、又把观测值 $\mathbf{z}_i$ 画成节点。\textbf{现象}：得到的图与库（GTSAM/g2o）的数据结构对不上，且分不清“已知 vs 待估”。\textbf{根本原因}：因子图的\emph{节点只有待估变量}（位姿、路标）；测量是\emph{已知量}，藏在因子内部（\cref{def:factor-graph}），把测量当节点会得到一张贝叶斯网而非因子图。\textbf{正确做法}：你要\emph{求}的量才是变量节点；你\emph{测}到的量决定因子的残差与权重。\textbf{自检}：每个变量节点都应是优化变量；每个因子都应能写成一项 $\|\mathbf{e}_i\|_{\boldsymbol{\Sigma}_i}^2$。
\end{pitfall}

\begin{practice}
\begin{enumerate}
  \item 对 \cref{fig:factor-graph}，按 \cref{eq:fg-product} 写出后验的因子乘积，并据 \cref{thm:fg-map} 写出对应的最小二乘目标。
  \item 在该图上加一条 $\mathbf{x}_1$–$\mathbf{x}_3$ 的回环因子，画出新图，说明它如何额外约束轨迹、如何减小累积漂移。
  \item （概念）解释“边缘化 $\mathbf{x}_1$”为何会在 $\mathbf{x}_2$ 与 $\mathbf{y}_1$ 间引入新边，从而破坏稀疏——联系 \cref{eq:info-marginal} 的舒尔补与平滑/滤波之别。
  \item 按 \cref{thm:elimination}，对 \cref{eq:toy-factorize} 的玩具例用排序 $\boldsymbol{\ell}_1,\boldsymbol{\ell}_2,\mathbf{p}_1,\mathbf{p}_2,\mathbf{p}_3$ 走一遍变量消元，写出每步的分隔集与新因子，并对照得到的贝叶斯网分解 $p=p(\boldsymbol{\ell}_1\mid\mathbf{p}_1,\mathbf{p}_2)p(\boldsymbol{\ell}_2\mid\mathbf{p}_3)p(\mathbf{p}_1\mid\mathbf{p}_2)p(\mathbf{p}_2\mid\mathbf{p}_3)p(\mathbf{p}_3)$。
  \item （跨章综合）说明变量消元中“消去路标得到关于位姿的约简系统”如何对应 \cref{ch:nlopt} 的 BA 箭头矩阵 Schur 消元——二者是同一稀疏消元的两种叙述。
\end{enumerate}
\end{practice}

\subsection{消息传递：高斯置信传播（GBP，选读）}
\label{sec:est-gbp}

变量消元（\cref{thm:elimination}）给出\emph{精确} MAP，但它要一个\emph{全局}消元序、本质\emph{串行}：必须按序逐个消、每步依赖前面的结果，且消元会产生 fill-in。还有一条性质截然不同的路——\textbf{消息传递}：每个节点\emph{只与邻居}通信，反复迭代直至收敛，无需全局排序、天然\emph{并行/分布式}。当因子图上一切都是高斯时，这套置信传播（belief propagation）特化为\textbf{高斯置信传播}（Gaussian belief propagation, GBP）\cite{davison2019futuremapping2,ortiz2021gbp}。它与本章的因子图/信息形式无缝衔接，恰好用来对照“精确 vs 近似、串行 vs 分布式”两种推断。

\paragraph{两类消息与置信，全用信息形式相加/舒尔补。}
GBP 在变量节点与因子节点间来回传\emph{高斯消息}，一律用信息形式 $(\boldsymbol{\eta},\boldsymbol{\Lambda})$ 表示（\cref{eq:info-form}）。设变量 $x_i$、与之相邻的因子集 $\mathcal{N}(i)$：
\begin{itemize}
  \item \textbf{变量置信} $b(x_i)$ $=$ 所有入射因子消息之\emph{积}，在信息形式下即\emph{相加}（\cref{prop:info-add}）：
  \begin{equation}
    \boldsymbol{\Lambda}_{b(i)}=\sum_{s\in\mathcal{N}(i)}\boldsymbol{\Lambda}_{s\to i},\qquad
    \boldsymbol{\eta}_{b(i)}=\sum_{s\in\mathcal{N}(i)}\boldsymbol{\eta}_{s\to i}.
    \label{eq:gbp-belief}
  \end{equation}
  \item \textbf{变量到因子}消息 $x_i\to\phi_s$ $=$ 置信\emph{除去}来自该因子的那一条（“别人告诉我的”）：$\boldsymbol{\Lambda}_{i\to s}=\sum_{t\in\mathcal{N}(i)\setminus s}\boldsymbol{\Lambda}_{t\to i}$，$\boldsymbol{\eta}_{i\to s}$ 同理。
  \item \textbf{因子到变量}消息 $\phi_s\to x_i$ $=$ 把因子自身的高斯势\emph{乘上}其它相邻变量传来的消息（信息相加），再把那些\emph{其它}变量\emph{边缘化}掉。这“先乘后边缘”正是本章已备好的两件工具：相乘 $=$ 信息相加，边缘化 $=$ 对被边缘块做舒尔补（\cref{eq:info-marginal}）。记因子联合信息阵按“目标 $i$ / 其它 $o$”分块 $\boldsymbol{\Lambda}=\big[\begin{smallmatrix}\boldsymbol{\Lambda}_{ii}&\boldsymbol{\Lambda}_{io}\\\boldsymbol{\Lambda}_{oi}&\boldsymbol{\Lambda}_{oo}\end{smallmatrix}\big]$、$\boldsymbol{\eta}=[\boldsymbol{\eta}_i;\boldsymbol{\eta}_o]$（已并入入射消息），则
  \begin{equation}
    \boldsymbol{\Lambda}_{s\to i}=\boldsymbol{\Lambda}_{ii}-\boldsymbol{\Lambda}_{io}\boldsymbol{\Lambda}_{oo}^{-1}\boldsymbol{\Lambda}_{oi},\qquad
    \boldsymbol{\eta}_{s\to i}=\boldsymbol{\eta}_i-\boldsymbol{\Lambda}_{io}\boldsymbol{\Lambda}_{oo}^{-1}\boldsymbol{\eta}_o.
    \label{eq:gbp-factor-msg}
  \end{equation}
\end{itemize}
节点按任意次序（同步或异步）反复执行这三式，直到置信稳定。每步只动\emph{局部}量——这正是它能分布式跑的原因。非线性因子在当前估计处重新线性化（如 BA），鲁棒因子按核函数降权调 $\boldsymbol{\Lambda}$\cite{davison2019futuremapping2}，机理与 \cref{sec:est-fg-linear}、\cref{ch:nlopt} 一致。

\begin{insight}[GBP 与变量消元：同一信息代数的两种调度]
\cref{eq:gbp-factor-msg} 的舒尔补与变量消元 \cref{thm:elimination} 用的是\emph{同一套}高斯条件化/边缘化代数（都落到 \cref{thm:schur-cond}、\cref{eq:info-marginal}），区别只在\emph{调度}与\emph{是否精确}：
\begin{itemize}
  \item 在\emph{无环}（树状）因子图上，GBP 一遍前后向消息就给出\textbf{精确}边缘——此时它\emph{等价于}一次变量消元 $+$ 回代，只是把“全局消元序”换成了“沿树传消息”。
  \item 在\emph{有环}（loopy）图上——SLAM 因回环几乎总有环——GBP 变成\emph{迭代近似}：可证\emph{若收敛，则各变量的均值收敛到精确 MAP；但协方差（精度）一般偏乐观（过自信）}\cite{weiss2001correctness}。实践中用\emph{阻尼}（消息按比例混合新旧值）改善收敛。
\end{itemize}
对照本章主线：变量消元/Bayes 树（\cref{sec:est-smooth-filter}）走\emph{精确、串行、需全局序}，GBP 走\emph{近似、并行、纯局部}。二者不是对立而是一谱之两端——精确推断要全局协调，分布式推断拿“精确协方差”换“可扩展与并行”。
\end{insight}

\begin{insight}[为什么有人关心 GBP：硬件与多机（一种前瞻观点，选读）]
近年一种观点把 GBP 视为\emph{空间智能}（Spatial AI）走向专用并行硬件的天然算法\cite{davison2019futuremapping2}。\textbf{(1) 存算一体}：因子图节点与“核间消息传递”可直接映射到图处理器（如 Graphcore IPU——上千核、各持本地内存、核间高速互联），把整张图\emph{存在芯片上}、就地计算，免去与外部内存的搬运；在这类硬件上做光束法平差，有报道比 CPU 快约 $30\times$\cite{ortiz2020bundleipu}。\textbf{(2) 异步多机}：把一张大因子图\emph{切}给一群机器人各持一片、经 Wi-Fi 等\emph{即兴异步}通信传跨界消息，可做分布式多机定位（“机器人网”），并对掉线、传感器失效鲁棒\cite{murai2024robotweb}。这些都把因子图从“一次性求解的问题”变成“持续就地更新的活体表示”。本书对此点到为止；它与本章“平滑 vs 滤波”“精确 vs 近似”的张力一脉相承，详尽展开属系统/前沿专题。
\end{insight}

\paragraph{为什么是“现在”才押注分布式：算力的风向变了（选读）。}
上面这套“把图铺到硬件上”的想法，并非一直成立——它是被\emph{算力的风向}逼出来的\cite{davison2026computational}。SLAM 早年所处的，是单核 CPU 时钟频率每一两年翻一番的年代，串行地跑一遍消元、再跑一遍回代，性能会自动随硬件水涨船高。这条捷径如今断了：晶体管密度（摩尔定律）虽仍在涨，但“晶体管缩小而功率密度不变”的丹纳德缩放（Dennard scaling）已失效——纳米尺度的晶体管漏电、发热，把单核时钟频率死死压在数 GHz 的“功率墙”下。于是处理器的演进方向，从“更快的单核”转向“更多、更专、更并行的核”。对要塞进眼镜、靠电池供电的嵌入式空间 AI 而言，功耗是头等约束，异构、专用、大规模并行几乎是唯一出路。这正呼应了 AI 领域那条著名的“苦涩教训”（Bitter Lesson）——\emph{能随算力扩展的通用方法，长远看总是胜出}。若真信“低功耗的算力与存储会越来越分散”，那么该押注的算法，就得是\emph{天生分布式、天生局部}的那一类——而 GBP 的纯局部消息传递，恰好是为这种硬件量身定做的。

\begin{insight}[一个对偶：世界模型图 $=$ 信息矩阵稀疏性的硬件投影（选读）]
为什么因子图能这样干净地铺到一片多核芯片上？根子在本章反复出现的\emph{信息矩阵稀疏性}。回忆 \cref{sec:est-fim-laplacian}：SLAM 的信息阵 $\boldsymbol\Lambda$ 不是稠密的，它的非零结构\emph{恰好就是问题的图}——谁与谁有共同观测，$\boldsymbol\Lambda$ 对应块才非零。本书一路把这条稀疏性当作\emph{算法效率}之源：稀疏 Cholesky、Schur 补、Bayes 树，无一不是在“吃”这块稀疏。空间智能的硬件构想，不过是把同一条稀疏性\emph{投影到物理空间}\cite{davison2026computational}：地图里“邻近”的特征/位姿强相关、连强边，“遥远”的弱相关、可剪成弱边乃至断开——于是一张本就稀疏的世界模型图，被原样\emph{铺}到一片“核 $+$ 本地内存 $+$ 近邻互联”的硬件上，每个核只存图的一小片、只与拓扑上的近邻通信。\textbf{“算法的稀疏图”与“硬件的连接图”这两张图能对上}，正是低功耗的关键：数据不必在核与远端内存之间来回长途搬运（这恰是前面“比特 $\times$ 毫米”那条度量所计的代价）。一句话——信息矩阵的稀疏性，在算法侧表现为“可被高效求解”，在硬件侧就表现为“可被分布式铺开”；二者是同一个结构的两面。
\end{insight}

\paragraph{硬件全景：不止处理器，还有传感器与云（选读）。}
把视野放宽，一套嵌入式空间 AI 的硬件版图大致分三层\cite{davison2026computational}。\textbf{处理器}一侧，是从 CPU 到 GPU、再到更激进的图处理器（如前述 IPU，乃至 SpiNNaker 这类“百万核、重通信”的类脑机）的谱系；商用设备里已能见到为低功耗实时感知定制的专用芯片（如早期 HoloLens 的 HPU、Apple Vision Pro 的 R1 协处理器），把 SLAM、眼动、手势等都搬到端上跑、无需外接电池。\textbf{传感器}一侧，相机的概念正被拓宽：事件相机（event camera）只传\emph{亮度变化}、以极低码率编码视频信息，是“只报与预测不符之处”这一思想的雏形；像平面内处理（in-plane processing，如 SCAMP 系视觉芯片）则把特征检测、光流这类\emph{近传感器}计算直接做在像素阵列上，主旨同样是\emph{压低传感器到主处理器的数据搬运}。\textbf{云}是第三层：云端算力可近似看作无限且廉价，但与设备间的通信（尤其传视频）在功耗与延迟上都不便宜，且空间、水下、地下等场景未必有稳定连接。三层的共同主线只有一句——\emph{尽一切可能减少数据搬运}，而这与“把稀疏因子图就地铺开”是同一套逻辑。

\paragraph{把“学习”也收进同一张图：持续学习与 GBP-Learning（选读）。}
最后一个引申，把这条硬件主线接回 \cref{ch:learning} 的立场。那一章反复申明：学习不该\emph{取代}母方程，而该作为\emph{可换的零件}插进同一张因子图——“换零件，不换母方程”。沿这条线推到极致，会得到一个颇激进的设想\cite{davison2026computational}：连\emph{学习本身}也用因子图与 GBP 来做。其要旨是，在贝叶斯机器学习的视角下，“估计状态”与“学习参数”\emph{本无本质区别}——都是同一套概率推断律，差别只在结构与规模；于是可以把一个过参数化的学习结构（仿神经网络的层次）直接\emph{搭成因子图}：可学参数成了图里的\emph{变量}（推断其值即“学权重”），非线性因子扮演“软开关”的神经元，整张图照样用标准 GBP 来推断。这样一来，“训练”与“推断”、“离线学”与“在线用”之间那道人为的界线被抹平，学习得以\emph{持续}、\emph{自监督}地随系统运行而发生——一台机器人可以一边重建场景、一边就地学出该场景的平滑或分割模型，而这一切与几何估计共用同一套分布式、可铺到硬件上的机理。这是否可行尚是开放问题，但它精确地落在本书一以贯之的立场上：\emph{把学习与估计熔进同一张图、同一条推断律}，而非让二者各说各话。

% ════════════════════════════════════════════════════════════════
\section{信息矩阵就是图的拉普拉斯：精度的极限}
\label{sec:est-fim-laplacian}

\cref{sec:est-crlb} 给了 CRLB 的一般形式：无偏估计的协方差被 Fisher 信息阵之逆下界，$\mathrm{cov}(\hat{\boldsymbol{\theta}})\succeq\boldsymbol{\mathcal{I}}_{\boldsymbol{\theta}}^{-1}$，且线性高斯下最小二乘解恰达此界（\cref{thm:normal-eq}）。本节追问一个更结构化的问题：SLAM 的 Fisher 信息阵\emph{长什么样}？答案出人意料地干净——它由\emph{问题的图}（谁观测谁）决定，本质就是位姿图的\textbf{拉普拉斯矩阵}\cite{carlone2026handbook}。这把“信息矩阵”从一个抽象的代数对象，升级为一个可\emph{从图结构直接读出}的量，并直接服务于主动 SLAM 与测量剪枝：\emph{该往哪看、哪些测量值得留}。

\paragraph{前置：图的关联矩阵与（加权）拉普拉斯。}
位姿图的每条边是一次\emph{相对}测量（里程计/回环），把 $n$ 个位姿当节点、$m$ 条测量当有向边。定义\emph{关联矩阵} $\mathbf{A}_{\mathcal{G}}\in\mathbb{R}^{m\times n}$：第 $e$ 行对应边 $e=(i\to j)$，在第 $i$ 列填 $-1$、第 $j$ 列填 $+1$、余为 $0$。则\emph{图拉普拉斯}（标量、无权）
\begin{equation}
  \boldsymbol{\mathcal{L}}=\mathbf{A}_{\mathcal{G}}^\top\mathbf{A}_{\mathcal{G}},\qquad
  \boldsymbol{\mathcal{L}}_{ij}=\begin{cases}\deg(i),&i=j,\\ -1,&\{i,j\}\in\mathcal{E},\\ 0,&\text{否则},\end{cases}
  \label{eq:fim-laplacian-def}
\end{equation}
对角是节点的\emph{度}、非对角是“相邻则 $-1$”。给每条边一个信息权 $w_e>0$（测量越准越大），得\emph{加权拉普拉斯} $\boldsymbol{\mathcal{L}}_w=\mathbf{A}_{\mathcal{G}}^\top\mathbf{W}\mathbf{A}_{\mathcal{G}}$，$\mathbf{W}=\mathrm{diag}(w_1,\dots,w_m)$。$\boldsymbol{\mathcal{L}}$ 的两条性质后文要用：最小特征值恒为 $0$（特征向量 $\mathbf{1}$，因每行和为零——\emph{相对}测量定不了绝对位置，正是 \cref{sec:est-observability} 的 gauge 零空间）；其\emph{次小}特征值（\emph{代数连通度}，Fiedler 值）度量图连得有多牢。锚定一个位姿（删去对应行/列）得\emph{约简}拉普拉斯，满秩可逆。

\begin{theorem}[已知朝向的位姿图：FIM $=$ 加权拉普拉斯 $\otimes\,\mathbf{I}$]\label{thm:fim-laplacian}
设各位姿\emph{朝向已知}、只估\emph{位置} $\mathbf{t}_i\in\mathbb{R}^3$，第 $k$ 条相对测量 $\mathbf{z}_k=\mathbf{R}_{i_k}^\top(\mathbf{t}_{j_k}-\mathbf{t}_{i_k})+\mathbf{n}_k$，$\mathbf{n}_k\sim\mathcal{N}(\mathbf{0},w_k^{-1}\mathbf{I}_3)$（各向同性）。则全体位置的 Fisher 信息阵为加权（约简）拉普拉斯与 $\mathbf{I}_3$ 的克罗内克积：
\begin{equation}
  \boldsymbol{\mathcal{I}}=\boldsymbol{\mathcal{L}}_w\otimes\mathbf{I}_3,\qquad \boldsymbol{\mathcal{L}}_w=\mathbf{A}_{\mathcal{G}}^\top\mathbf{W}\mathbf{A}_{\mathcal{G}}.
  \label{eq:fim-laplacian}
\end{equation}
\end{theorem}
\begin{proof}
线性高斯下 FIM $=\mathbf{J}^\top\boldsymbol{\Sigma}^{-1}\mathbf{J}$（\cref{thm:crlb} 的高斯特例，亦即 \cref{eq:info-project} 的信息投影）。把测量堆叠：$\mathbf{z}=\mathbf{R}^\top(\mathbf{A}_{\mathcal{G}}\otimes\mathbf{I}_3)\,\mathbf{t}+\mathbf{n}$，其中 $\mathbf{R}=\mathrm{blkdiag}(\mathbf{R}_{i_1},\dots,\mathbf{R}_{i_m})$ 收集各测量所在位姿的（已知）朝向，故雅可比 $\mathbf{J}=\mathbf{R}^\top(\mathbf{A}_{\mathcal{G}}\otimes\mathbf{I}_3)$，噪声信息 $\boldsymbol{\Sigma}^{-1}=\mathbf{W}\otimes\mathbf{I}_3$。于是
\[
  \boldsymbol{\mathcal{I}}=\mathbf{J}^\top\boldsymbol{\Sigma}^{-1}\mathbf{J}
  =(\mathbf{A}_{\mathcal{G}}\otimes\mathbf{I}_3)^\top\,\mathbf{R}\,(\mathbf{W}\otimes\mathbf{I}_3)\,\mathbf{R}^\top\,(\mathbf{A}_{\mathcal{G}}\otimes\mathbf{I}_3).
\]
两步收尾：(i) 各向同性噪声的信息块\emph{旋转不变}，$\mathbf{R}(\mathbf{W}\otimes\mathbf{I}_3)\mathbf{R}^\top=\mathbf{W}\otimes\mathbf{I}_3$（每个 $3\times3$ 块是 $w_k\mathbf{R}_{i_k}\mathbf{R}_{i_k}^\top=w_k\mathbf{I}_3$），\emph{已知朝向被消掉}；(ii) 用克罗内克混合积律 $(\mathbf{A}\otimes\mathbf{I})(\mathbf{W}\otimes\mathbf{I})(\mathbf{A}^\top\otimes\mathbf{I})=(\mathbf{A}\mathbf{W}\mathbf{A}^\top)\otimes\mathbf{I}$（\cref{eq:matderiv-kron-mixed}），即得 \cref{eq:fim-laplacian}。
\end{proof}
\noindent 这条等式漂亮在：$\boldsymbol{\mathcal{I}}$ 的\emph{谱}就是 $\boldsymbol{\mathcal{L}}_w$ 的谱（每个特征值重数 $3$）——精度的全部信息浓缩在一个\emph{标量}图拉普拉斯里，维数比 FIM 小一个 $\mathbf{I}_3$ 的因子。一般位姿图（朝向\emph{未知}、全 $\SE(3)$ 状态）的 FIM 不再这么干净，但仍是“\emph{逐边}的拉普拉斯型求和”\cite{carlone2026handbook}：
\begin{equation}
  \boldsymbol{\mathcal{I}}(\mathbf{x})=\sum_{k=1}^{m}\boldsymbol{\mathcal{L}}_k\otimes\Big(\mathrm{Ad}(\mathbf{T}_{i_k}^{-1})^\top\boldsymbol{\Sigma}_k^{-1}\mathrm{Ad}(\mathbf{T}_{i_k}^{-1})\Big),
  \qquad \boldsymbol{\mathcal{L}}_k=\mathbf{a}_k\mathbf{a}_k^\top,
  \label{eq:fim-laplacian-general}
\end{equation}
$\mathbf{a}_k$ 是约简关联矩阵的第 $k$ 列、$\mathrm{Ad}(\cdot)$ 是 $\SE(3)$ 伴随（\cref{ch:lie}），它把各测量的信息搬到全局系——正是这个伴随项带来\emph{轨迹依赖}、破坏了 \cref{eq:fim-laplacian} 的洁净分解；令各边 $\boldsymbol{\Sigma}_k^{-1}=w_k\mathbf{I}$、朝向已知即退回 \cref{eq:fim-laplacian}（$\sum_k w_k\boldsymbol{\mathcal{L}}_k=\boldsymbol{\mathcal{L}}_w$）。

\begin{insight}[把信息阵当图来读：最优实验设计与“信息永不损害”]
\cref{thm:fim-laplacian} 让“估计有多准”变成“图连得多牢”，于是 SLAM 的精度问题接上了\emph{最优实验设计}的成熟语言。FIM 是矩阵，要排序好坏须压成标量，三个经典准则都是它的\emph{谱函数}：
\begin{itemize}
  \item \textbf{D-最优}：$\det\boldsymbol{\mathcal{I}}$（不确定椭球\emph{体积}最小，与 \cref{eq:gauss-mutual-info} 的互信息同源）；
  \item \textbf{A-最优}：$\mathrm{tr}(\boldsymbol{\mathcal{I}}^{-1})$（\emph{平均}方差最小）；
  \item \textbf{E-最优}：$\lambda_{\min}(\boldsymbol{\mathcal{I}})$ 最大（\emph{最坏}方向方差最小，对应 Fiedler 值/代数连通度）。
\end{itemize}
两条直接推论。\textbf{(1) 矩阵树定理}：由 \cref{eq:fim-laplacian}，$\det\boldsymbol{\mathcal{I}}=(\det\boldsymbol{\mathcal{L}}_w)^3$，而\emph{基尔霍夫矩阵树定理}说约简加权拉普拉斯的行列式 $=$ 图的\emph{加权生成树数目}\cite{carlone2026handbook}——于是 \emph{D-最优 $\Leftrightarrow$ 加权生成树最多 $\Leftrightarrow$ 图连得最牢}，把“连通度决定精度”这一直觉\emph{精确化}了。\textbf{(2) 信息永不损害}：每加一条测量边，就给 FIM 加一个半正定块（\cref{eq:fim-laplacian-general} 多一项），故 CRLB（按 Loewner 序）\emph{只降不升}——这正是 \cref{thm:schur-cond} 后“观测使协方差收缩”在全图尺度的回声。但\emph{加哪条}大不相同：闭一个\emph{大}回环（连起图上相距很远的两节点）比闭\emph{小}环更能提升连通度、更增信息。这把抽象的“信息矩阵”落到可操作的决策上——\emph{主动 SLAM}（规划轨迹去最大化预期 FIM）与\emph{测量剪枝}（终身 SLAM 里只留最有信息的边）都据此用拉普拉斯谱做代理，省去反复求解整个 SLAM\cite{carlone2026handbook}。
\end{insight}

\begin{note}[本节边界]
本节只谈“最优解\emph{有多准}”（精度极限，由图拉普拉斯刻画）。与之配对的另一问题——“能否\emph{不靠初值、可证全局最优}地\emph{解出}位姿图”（半定松弛 / SE-Sync）——属\emph{求解器}范畴，与 GN/LM 同处 \cref{ch:nlopt}（见该章可证最优一节）。本章分工与全书一致：\cref{ch:slam_est} 讲“要解什么、解得多准”，\cref{ch:nlopt} 讲“怎么解”。
\end{note}

\begin{practice}
\begin{enumerate}
  \item 对一条 $4$ 节点的链（$1$–$2$–$3$–$4$，三条等权 $w=1$ 的里程计边），写出关联矩阵 $\mathbf{A}_{\mathcal{G}}$ 与拉普拉斯 $\boldsymbol{\mathcal{L}}$，验证 $\boldsymbol{\mathcal{L}}\mathbf{1}=\mathbf{0}$；再加一条 $1$–$4$ 回环边，比较两图的代数连通度（次小特征值）变大还是变小。
  \item 用 \cref{eq:matderiv-kron-mixed} 补全 \cref{thm:fim-laplacian} 证明里 (ii) 的混合积律一步，并说明为何各向同性噪声是“朝向被消掉”的前提。
  \item （概念）由“信息永不损害”说明：为何闭\emph{大}回环通常比闭\emph{小}回环更降低整体不确定度？用生成树数目 / 代数连通度给直觉。
  \item （跨节）把 \cref{eq:fim-laplacian} 与 \cref{eq:matderiv-fim-reparam}（FIM 重参数化）联系：若把位置参数换成\emph{相对}坐标，FIM 如何随关联矩阵变换？
\end{enumerate}
\end{practice}

% ════════════════════════════════════════════════════════════════
\section{连续时间轨迹估计与 STEAM}
\label{sec:est-steam}

\paragraph{动机：传感器并不在同一拍子上敲门。}
迄今全章都把轨迹看成一串\emph{离散}状态 $\{\mathbf{x}_k\}$，每个 $k$ 对应一个测量时刻。这对同步相机够用，却在现代传感器栈前露怯：滚动快门逐行曝光、激光雷达逐点扫描、IMU 高频而相机低频、多传感器各有时钟——\emph{测量根本不发生在同一拍子上}。强行把它们对齐到公共时刻要么插值要么丢数据。一个更彻底的视角是：把状态当成时间的\emph{连续函数} $\mathbf{x}(t)$，需要哪个时刻的位姿就\emph{查询}哪个。

把轨迹写成连续函数有\emph{两条}路\cite{carlone2026handbook}。\textbf{参数法（样条）}：把位姿写成一组\emph{已知}时间基函数的加权和 $\mathbf{x}(t)=\sum_k\boldsymbol{\Psi}_k(t)\mathbf{c}_k$，待估的是有限个系数（控制点）；基函数取\emph{局部支撑}（如 B 样条），保证因子图仍稀疏。\textbf{非参数法（高斯过程，GP）}：用核函数刻画轨迹分布，由一个线性 SDE\emph{运动先验}生成；待估的是若干估计时刻的状态，先验自带正则。两条路都把“测量时刻 / 估计时刻 / 查询时刻”解耦，这是连续时间方法应对高频异步传感器的共同威力之源（\cref{sec:est-steam-query} 详述）。本书主线走 GP 这一支——它把运动先验显式写成因子、与观测在同一 MAP 里权衡，并自然给出查询协方差；而 GP 与样条在 $\SE(3)$ 上殊途同归（\cref{sec:est-steam-spline} 点明取舍）。本节把 GP 主线在 $\SE(3)$ 上完整落地，得到\textbf{同时轨迹估计与建图}（Simultaneous Trajectory Estimation And Mapping, STEAM）\cite{anderson2015steam}——它是 SLAM 的连续时间变体：解一次批量 MAP 后，能以 $O(1)$ 代价查询\emph{任意}时刻的位姿、速度\emph{及其协方差}。

\paragraph{反面：朴素地把 GP 搬上李群会卡住。}
向量空间里“连续时间估计 $=$ 高斯过程（GP）回归”早有定论（\cref{ch:nlopt} 附录的向量空间 derivation 给了骨架）：用线性时不变随机微分方程（SDE）作先验，状态 $\mathbf{x}(t)$ 就是一个 GP，批量解仍块三对角、可 $O(1)$ 插值。但位姿活在 $\SE(3)$ 上，最自然的连续先验
\begin{equation}
  \dot{\mathbf{T}}(t)=\boldsymbol{\varpi}(t)^\wedge\mathbf{T}(t),\qquad
  \dot{\boldsymbol{\varpi}}(t)=\mathbf{w}(t),\qquad
  \mathbf{w}\sim\mathcal{GP}\big(\mathbf{0},\,\mathbf{Q}_C\,\delta(t-t')\big)
  \label{eq:steam-global-sde}
\end{equation}
是\emph{非线性}的（$\mathbf{T}\in\SE(3)$、$\boldsymbol{\varpi}\in\mathbb{R}^6$，第一式右端含 $\boldsymbol{\varpi}^\wedge\mathbf{T}$）。这里 $\boldsymbol{\varpi}=[\boldsymbol{\nu};\boldsymbol{\omega}]$ 是机体系广义速度（线速度 $+$ 角速度，记号见 \cref{thm:se3-kinematics}），白噪声 $\mathbf{w}$ 注入广义\emph{加速度}：无噪声时 $\boldsymbol{\varpi}$ 恒定——这就是“白噪声加速度”（white-noise-on-acceleration, \textbf{WNOA}）或“常速度”先验。麻烦在于：非线性 SDE \emph{无法}像线性 SDE 那样闭式随机积分，\cref{ch:eskf} 的非线性高斯工具也不直接适用。直接在全局变量上做 GP 会卡死在“积不出转移函数”这一步。

\paragraph{历史与思路：用一串\emph{局部} GP 把轨迹缝起来。}
解法由 Barfoot、Tong 与 Särkkä 提出\cite{barfoot2014steam}、并由 Anderson 与 Barfoot 系统化为 STEAM\cite{anderson2015steam}：既然全局非线性、局部却近线性，那就\emph{分段}。沿轨迹选一串估计时刻 $t_0<t_1<\dots<t_K$，在每段 $[t_k,t_{k+1}]$ 引入一个\emph{局部}李代数变量 $\boldsymbol{\xi}_k(t)\in\se(3)$，把全局位姿\emph{复合}在该段起点上：
\begin{equation}
  \mathbf{T}(t)=\exp\!\big(\boldsymbol{\xi}_k(t)^\wedge\big)\,\mathbf{T}(t_k),
  \qquad\Longleftrightarrow\qquad
  \boldsymbol{\xi}_k(t)=\ln\!\big(\mathbf{T}(t)\mathbf{T}(t_k)^{-1}\big)^\vee.
  \label{eq:steam-local}
\end{equation}
只要该段内转动不过大，$\boldsymbol{\xi}_k(t)$ 就是全局位姿的良好局部坐标。妙处在于：\emph{在局部变量上}定义 GP，可把先验取成线性 SDE，于是随机积分有闭式。这与 \cref{thm:so3-kinematics}/\cref{thm:se3-kinematics} 同源——那里把连续旋转/位姿的非线性运动学用李代数坐标率 $\dot{\boldsymbol{\xi}}=\boldsymbol{\mathcal{J}}_r^{-1}\boldsymbol{\varpi}$ 线性化为切空间里的向量动力学，正是这里“化非线性为局部线性”的几何根。

\subsection{局部 LTI 先验与闭式随机积分}
\label{sec:est-steam-prior}

我们把 WNOA 先验\emph{间接}定义在局部变量上：令局部坐标的二阶导受白噪声驱动，
\begin{equation}
  \ddot{\boldsymbol{\xi}}_k(t)=\mathbf{w}_k(t),\qquad \mathbf{w}_k\sim\mathcal{GP}\big(\mathbf{0},\,\mathbf{Q}_C\,\delta(t-t')\big),
  \label{eq:steam-local-sde}
\end{equation}
$\mathbf{Q}_C\in\mathbb{R}^{6\times6}$ 是\emph{功率谱密度}（power spectral density, PSD），调节六自由度上的平滑度。引入 $\boldsymbol{\psi}_k:=\dot{\boldsymbol{\xi}}_k$，把 \cref{eq:steam-local-sde} 写成\emph{一阶线性}（LTI）SDE，状态 $\boldsymbol{\gamma}_k:=[\boldsymbol{\xi}_k;\boldsymbol{\psi}_k]\in\mathbb{R}^{12}$：
\begin{equation}
  \dot{\boldsymbol{\gamma}}_k(t)
  =\underbrace{\begin{bmatrix}\mathbf{0}&\mathbf{I}\\\mathbf{0}&\mathbf{0}\end{bmatrix}}_{\mathbf{A}}\boldsymbol{\gamma}_k(t)
  +\underbrace{\begin{bmatrix}\mathbf{0}\\\mathbf{I}\end{bmatrix}}_{\mathbf{L}}\mathbf{w}_k(t).
  \label{eq:steam-lti}
\end{equation}
（注意 $\boldsymbol{\psi}_k=\dot{\boldsymbol{\xi}}_k$ 是\emph{局部坐标率}，与机体系广义速度 $\boldsymbol{\varpi}$ 经 \cref{thm:se3-kinematics} 的 $\boldsymbol{\psi}_k=\boldsymbol{\mathcal{J}}_r(\boldsymbol{\xi}_k)^{-1}\boldsymbol{\varpi}$ 相联，二者并不相等。）线性 SDE 可\emph{闭式}随机积分（这正是切换到局部变量的目的）：$\boldsymbol{\gamma}_k(t)$ 是一个 GP，
\begin{equation}
  \boldsymbol{\gamma}_k(t)\sim\mathcal{GP}\Big(\underbrace{\boldsymbol{\Phi}(t,t_k)\check{\boldsymbol{\gamma}}_k(t_k)}_{\text{均值}},\ \underbrace{\boldsymbol{\Phi}(t,t_k)\check{\mathbf{P}}_k\boldsymbol{\Phi}(t,t_k)^\top+\mathbf{Q}(t-t_k)}_{\text{协方差}}\Big),
  \label{eq:steam-gp}
\end{equation}
其中\emph{转移函数}与\emph{累积过程噪声}由 $\mathbf{A},\mathbf{L}$ 解得（标量化 PSD 写 $\mathbf{Q}_C$）：
\begin{equation}
  \boldsymbol{\Phi}(t,s)=\exp\!\big(\mathbf{A}(t-s)\big)=\begin{bmatrix}\mathbf{I}&(t-s)\mathbf{I}\\\mathbf{0}&\mathbf{I}\end{bmatrix},\quad
  \mathbf{Q}(\Delta t)=\begin{bmatrix}\tfrac13\Delta t^3\,\mathbf{Q}_C&\tfrac12\Delta t^2\,\mathbf{Q}_C\\[2pt]\tfrac12\Delta t^2\,\mathbf{Q}_C&\Delta t\,\mathbf{Q}_C\end{bmatrix}.
  \label{eq:steam-Phi-Q}
\end{equation}

\begin{derivation}[闭式 $\boldsymbol{\Phi}$ 与 $\mathbf{Q}$ 的来历]
转移函数解 $\dot{\boldsymbol{\Phi}}=\mathbf{A}\boldsymbol{\Phi}$、$\boldsymbol{\Phi}(s,s)=\mathbf{I}$。因 $\mathbf{A}^2=\mathbf{0}$（幂零），矩阵指数级数\emph{只剩两项}：$\boldsymbol{\Phi}(t,s)=\mathbf{I}+\mathbf{A}(t-s)$，即 \cref{eq:steam-Phi-Q} 左式——上三角块 $(t-s)\mathbf{I}$ 表达“位置 $=$ 旧位置 $+$ 速度 $\times$ 时间”。累积过程噪声是 LTI SDE 的标准结果（连续 $\to$ 离散的过程噪声积分，与 \cref{ch:eskf} 一致）：
\[
  \mathbf{Q}(\Delta t)=\int_0^{\Delta t}\boldsymbol{\Phi}(\Delta t,s)\,\mathbf{L}\mathbf{Q}_C\mathbf{L}^\top\,\boldsymbol{\Phi}(\Delta t,s)^\top\,\mathrm{d}s .
\]
代入 $\boldsymbol{\Phi}(\Delta t,s)\mathbf{L}=[(\Delta t-s)\mathbf{I};\,\mathbf{I}]$，逐块积分 $\int_0^{\Delta t}(\Delta t-s)^2\mathrm{d}s=\tfrac13\Delta t^3$、$\int_0^{\Delta t}(\Delta t-s)\mathrm{d}s=\tfrac12\Delta t^2$、$\int_0^{\Delta t}\mathrm{d}s=\Delta t$，即得 \cref{eq:steam-Phi-Q} 右式。\hfill$\square$
\end{derivation}

\begin{insight}[GP 先验的稀疏性来自马尔可夫性]
通常 GP 回归要在\emph{所有}待估时刻两两之间建核矩阵（稠密）。但由 \cref{eq:steam-lti} 这个 LTI SDE 构造的 GP，其\emph{逆核}（精度矩阵）天然稀疏：状态 $\boldsymbol{\gamma}_k$ 是\emph{马尔可夫}的——$t$ 时刻只通过 $\boldsymbol{\Phi}$ 依赖\emph{紧邻}的前一时刻，不直接依赖更早。于是先验只在\emph{相邻}估计时刻间产生约束，精度矩阵 $\boldsymbol{\Lambda}_{\text{prior}}=\mathbf{F}^{-\top}\mathbf{Q}^{-1}\mathbf{F}^{-1}$ 是\textbf{块三对角}的（\cref{eq:steam-Finv} 给出 $\mathbf{F}^{-1}$ 的块下双对角形）。这与本章 \cref{sec:est-batch-recursive} 里“离散运动链 $\Rightarrow$ 块三对角”\emph{同一根源}（马尔可夫性 $\to$ 局部耦合 $\to$ 带状信息阵），只是先验从离散运动模型换成了连续时间 SDE。SDE 阶数越高（如下文 WNOJ）带越宽，但仍是常数带宽。
\end{insight}

\subsection{运动先验因子与线性化}
\label{sec:est-steam-factor}

要把这先验放进 MAP，须把它写成代价项（先验的负对数似然）。\cref{eq:steam-gp} 的 GP 给出相邻状态的条件分布，先在\emph{局部}变量上定义\emph{二元因子}误差（首时刻为锚定的一元因子）：把 $\boldsymbol{\gamma}_{k-1}(t)$ 在 $t_k$ 的实际值与“由 $t_{k-1}$ 经转移函数 $\boldsymbol{\Phi}$ 推到 $t_k$ 的预测值”相减，
\begin{equation}
  \mathbf{e}_{v,k}=
  \begin{cases}
    \begin{bmatrix}-\ln\!\big(\mathbf{T}(t_0)\check{\mathbf{T}}_0^{-1}\big)^\vee\\[2pt]\check{\boldsymbol{\varpi}}_0-\boldsymbol{\varpi}(t_0)\end{bmatrix}, & k=0,\\[12pt]
    \boldsymbol{\Phi}(t_k,t_{k-1})\big(\boldsymbol{\gamma}_{k-1}(t_{k-1})-\check{\boldsymbol{\gamma}}_{k-1}(t_{k-1})\big)-\big(\boldsymbol{\gamma}_{k-1}(t_k)-\check{\boldsymbol{\gamma}}_{k-1}(t_k)\big), & k>0,
  \end{cases}
  \label{eq:steam-error-local}
\end{equation}
其结构正是 \cref{eq:steam-gp} 的\emph{GP 转移条件残差}：若局部状态 $\boldsymbol{\gamma}_{k-1}(t_k)$ 恰好等于由 $\boldsymbol{\gamma}_{k-1}(t_{k-1})$ 经 $\boldsymbol{\Phi}(t_k,t_{k-1})$ 传播得到的均值，误差即为零；偏离越大，先验负对数似然（即代价 $J_{v,k}$，见下）越高——这把“相邻两态是否符合常速度先验”量化成一个二元因子（GP 的稀疏马尔可夫性保证只连相邻两态）。最后用 \cref{eq:steam-local}（$\boldsymbol{\xi}_k=\ln(\mathbf{T}\mathbf{T}_k^{-1})^\vee$）与局部运动学 $\boldsymbol{\psi}=\boldsymbol{\mathcal{J}}_r^{-1}\boldsymbol{\varpi}$ 把局部变量换回\emph{全局}变量 $\{\mathbf{T}_k,\boldsymbol{\varpi}_k\}$，得可直接优化的等价形式：
\begin{equation}
  \mathbf{e}_{v,k}=
  \begin{cases}
    \begin{bmatrix}-\ln\!\big(\mathbf{T}_0\check{\mathbf{T}}_0^{-1}\big)^\vee\\[2pt]\check{\boldsymbol{\varpi}}_0-\boldsymbol{\varpi}_0\end{bmatrix}, & k=0,\\[14pt]
    \begin{bmatrix}(t_k-t_{k-1})\,\boldsymbol{\varpi}_{k-1}-\ln\!\big(\mathbf{T}_k\mathbf{T}_{k-1}^{-1}\big)^\vee\\[4pt]
    \boldsymbol{\varpi}_{k-1}-\boldsymbol{\mathcal{J}}_r\!\big(\ln(\mathbf{T}_k\mathbf{T}_{k-1}^{-1})^\vee\big)^{-1}\boldsymbol{\varpi}_k\end{bmatrix}, & k>0,
  \end{cases}
  \label{eq:steam-error}
\end{equation}
代价 $J_v=\sum_{k=0}^{K}\tfrac12\,\mathbf{e}_{v,k}^\top\mathbf{Q}_k^{-1}\mathbf{e}_{v,k}$，其中 $\mathbf{Q}_k=\mathbf{Q}(t_k-t_{k-1})$（$k>0$）、$\mathbf{Q}_0=\check{\mathbf{P}}_0$ 为初态协方差。\cref{eq:steam-error} 的两行各有清晰含义：第一行说“两位姿之差应等于‘速度 $\times$ 时间’”（常速度先验），第二行说“两端速度应一致”（经 $\boldsymbol{\mathcal{J}}_r^{-1}$ 把 $\boldsymbol{\varpi}_k$ 搬到 $t_{k-1}$ 的局部坐标率上比较，用的正是 \cref{eq:steam-local-sde} 的 $\boldsymbol{\psi}=\boldsymbol{\mathcal{J}}_r^{-1}\boldsymbol{\varpi}$）。该误差\emph{只含全局变量} $\{\mathbf{T}_k,\boldsymbol{\varpi}_k\}$，可直接优化。

\begin{note}[与 Barfoot 记号的对照]
本节把 Barfoot\cite{barfoot2024state} 第 11 章的\emph{左}雅可比 $\boldsymbol{\mathcal{J}}$、左扰动改写为本书\emph{右}扰动主线：位姿扰动取 $\mathbf{T}_k=\exp(\boldsymbol{\epsilon}_k^\wedge)\mathbf{T}_{\mathrm{op},k}$（本书右扰动约定 \cref{ch:lie}），$\boldsymbol{\mathcal{J}}_r$ 是 \cref{eq:se3-kinematics} 的右雅可比，$\boldsymbol{\mathcal{T}}=\mathrm{Ad}(\mathbf{T})$ 为 $6\times6$ 伴随。误差结构、稀疏性与原书一致，仅雅可比的左右惯例不同。Barfoot 把信息阵记作 $\mathbf{I}$，本书一律记 $\boldsymbol{\Lambda}$。
\end{note}

\paragraph{线性化得块下双对角的 $\mathbf{F}^{-1}$。}
为用高斯牛顿，在工作点 $\mathbf{x}_{\mathrm{op}}$ 处线性化。位姿用右扰动 $\mathbf{T}_k=\exp(\boldsymbol{\epsilon}_k^\wedge)\mathbf{T}_{\mathrm{op},k}$、速度用加性扰动 $\boldsymbol{\varpi}_k=\boldsymbol{\varpi}_{\mathrm{op},k}+\boldsymbol{\eta}_k$，堆为 $\boldsymbol{\varepsilon}_k=[\boldsymbol{\epsilon}_k;\boldsymbol{\eta}_k]\in\mathbb{R}^{12}$。\emph{第一步（位姿差的 BCH 一阶）}：代入扰动到 $\ln(\mathbf{T}_k\mathbf{T}_{k-1}^{-1})^\vee$，先用伴随把两端扰动并到一处 $\exp(\boldsymbol{\epsilon}_k^\wedge)\mathbf{T}_{\mathrm{op},k}\mathbf{T}_{\mathrm{op},k-1}^{-1}\exp(-\boldsymbol{\epsilon}_{k-1}^\wedge)\approx\exp\!\big((\boldsymbol{\epsilon}_k-\boldsymbol{\mathcal{T}}_{\mathrm{op},k,k-1}\boldsymbol{\epsilon}_{k-1})^\wedge\big)\mathbf{T}_{\mathrm{op},k}\mathbf{T}_{\mathrm{op},k-1}^{-1}$，再用 BCH 一阶 $\ln(\exp(\boldsymbol{\delta}^\wedge)\mathbf{T})^\vee\approx\ln(\mathbf{T})^\vee+\boldsymbol{\mathcal{J}}_r^{-1}\boldsymbol{\delta}$：
\begin{equation}
  \ln\!\big(\mathbf{T}_k\mathbf{T}_{k-1}^{-1}\big)^\vee\approx\ln\!\big(\mathbf{T}_{\mathrm{op},k}\mathbf{T}_{\mathrm{op},k-1}^{-1}\big)^\vee+\boldsymbol{\mathcal{J}}_{\mathrm{op},k,k-1}^{-1}\big(\boldsymbol{\epsilon}_k-\boldsymbol{\mathcal{T}}_{\mathrm{op},k,k-1}\boldsymbol{\epsilon}_{k-1}\big),
  \label{eq:steam-bch}
\end{equation}
记 $\boldsymbol{\mathcal{T}}_{\mathrm{op},k,k-1}=\mathrm{Ad}(\mathbf{T}_{\mathrm{op},k}\mathbf{T}_{\mathrm{op},k-1}^{-1})$、$\boldsymbol{\mathcal{J}}_{\mathrm{op},k,k-1}=\boldsymbol{\mathcal{J}}_r(\ln(\mathbf{T}_{\mathrm{op},k}\mathbf{T}_{\mathrm{op},k-1}^{-1})^\vee)$。\emph{第二步（逆右雅可比的 Taylor）}：误差第二行含 $\boldsymbol{\mathcal{J}}_r(\cdot)^{-1}$，其内部已由 \cref{eq:steam-bch} 线性化，故对逆右雅可比做一阶 Taylor $\boldsymbol{\mathcal{J}}_r^{-1}(\mathbf{x})\approx\mathbf{I}-\tfrac12\mathbf{x}^\curlywedge$：
\begin{equation}
  \boldsymbol{\mathcal{J}}_r\!\big(\ln(\mathbf{T}_k\mathbf{T}_{k-1}^{-1})^\vee\big)^{-1}\approx\boldsymbol{\mathcal{J}}_{\mathrm{op},k,k-1}^{-1}-\tfrac12\Big(\boldsymbol{\mathcal{J}}_{\mathrm{op},k,k-1}^{-1}\big(\boldsymbol{\epsilon}_k-\boldsymbol{\mathcal{T}}_{\mathrm{op},k,k-1}\boldsymbol{\epsilon}_{k-1}\big)\Big)^\curlywedge.
  \label{eq:steam-Jinv-taylor}
\end{equation}
\emph{$\tfrac12\boldsymbol{\varpi}^\curlywedge$ 项的由来}：把 \cref{eq:steam-Jinv-taylor} 代入第二行的 $-\boldsymbol{\mathcal{J}}_r(\cdot)^{-1}\boldsymbol{\varpi}_k$（取 $\boldsymbol{\varpi}_k\approx\boldsymbol{\varpi}_{\mathrm{op},k}$，扰动$\times$扰动二阶项舍去），出现 $+\tfrac12\big(\boldsymbol{\mathcal{J}}_{\mathrm{op},k,k-1}^{-1}(\boldsymbol{\epsilon}_k-\boldsymbol{\mathcal{T}}\boldsymbol{\epsilon}_{k-1})\big)^\curlywedge\boldsymbol{\varpi}_{\mathrm{op},k}$；用 $6\times6$ 伴随（cross）算子的反对称性 $\mathbf{a}^\curlywedge\mathbf{b}=-\mathbf{b}^\curlywedge\mathbf{a}$（李括号反对称，见 \cref{ch:lie}）改写为 $-\tfrac12\boldsymbol{\varpi}_{\mathrm{op},k}^\curlywedge\boldsymbol{\mathcal{J}}_{\mathrm{op},k,k-1}^{-1}(\boldsymbol{\epsilon}_k-\boldsymbol{\mathcal{T}}_{\mathrm{op},k,k-1}\boldsymbol{\epsilon}_{k-1})$——这恰是 \cref{eq:steam-FE} 中 $\mathbf{E}_k,\mathbf{F}_{k-1}$ 第二行块里 $\tfrac12\boldsymbol{\varpi}_{\mathrm{op},k}^\curlywedge\boldsymbol{\mathcal{J}}^{-1}(\cdots)$ 项的来历。把 \cref{eq:steam-bch,eq:steam-Jinv-taylor} 代回 \cref{eq:steam-error} 的两行即得
\begin{equation}
  \mathbf{e}_{v,k}(\mathbf{x})\approx\mathbf{e}_{v,k}(\mathbf{x}_{\mathrm{op}})+\mathbf{F}_{k-1}\boldsymbol{\varepsilon}_{k-1}-\mathbf{E}_k\boldsymbol{\varepsilon}_k,
  \label{eq:steam-lin}
\end{equation}
其中（$k>0$，$\boldsymbol{\mathcal{T}}_{\mathrm{op},k,k-1},\boldsymbol{\mathcal{J}}_{\mathrm{op},k,k-1}$ 同 \cref{eq:steam-bch}）
\begin{equation}
  \mathbf{F}_{k-1}=\begin{bmatrix}\boldsymbol{\mathcal{J}}_{\mathrm{op},k,k-1}^{-1}\boldsymbol{\mathcal{T}}_{\mathrm{op},k,k-1}&(t_k-t_{k-1})\mathbf{I}\\[2pt]\tfrac12\boldsymbol{\varpi}_{\mathrm{op},k}^\curlywedge\boldsymbol{\mathcal{J}}_{\mathrm{op},k,k-1}^{-1}\boldsymbol{\mathcal{T}}_{\mathrm{op},k,k-1}&\mathbf{I}\end{bmatrix},\quad
  \mathbf{E}_k=\begin{bmatrix}\boldsymbol{\mathcal{J}}_{\mathrm{op},k,k-1}^{-1}&\mathbf{0}\\[2pt]\tfrac12\boldsymbol{\varpi}_{\mathrm{op},k}^\curlywedge\boldsymbol{\mathcal{J}}_{\mathrm{op},k,k-1}^{-1}&\boldsymbol{\mathcal{J}}_{\mathrm{op},k,k-1}^{-1}\end{bmatrix}.
  \label{eq:steam-FE}
\end{equation}
堆叠全部时刻，运动先验的雅可比是\emph{块下双对角}的（与本章离散运动先验同结构，故沿用记号 $\mathbf{F}^{-1}$）：
\begin{equation}
  \mathbf{F}^{-1}=\begin{bmatrix}\mathbf{E}_0&&&\\-\mathbf{F}_0&\mathbf{E}_1&&\\&\ddots&\ddots&\\&&-\mathbf{F}_{K-1}&\mathbf{E}_K\end{bmatrix},\quad
  \mathbf{Q}=\mathrm{diag}(\check{\mathbf{P}}_0,\mathbf{Q}_1,\dots,\mathbf{Q}_K),
  \label{eq:steam-Finv}
\end{equation}
先验代价的二次近似 $J_v\approx\tfrac12(\mathbf{e}_v(\mathbf{x}_{\mathrm{op}})-\mathbf{F}^{-1}\delta\mathbf{x}_1)^\top\mathbf{Q}^{-1}(\mathbf{e}_v(\mathbf{x}_{\mathrm{op}})-\mathbf{F}^{-1}\delta\mathbf{x}_1)$，$\delta\mathbf{x}_1=[\boldsymbol{\varepsilon}_0;\dots;\boldsymbol{\varepsilon}_K]$。

\subsection{STEAM：把运动先验与观测拼成批量 MAP}
\label{sec:est-steam-map}

现在加入对路标的观测，得到完整 STEAM 问题。待估状态\emph{同时}含位姿、广义速度与路标：
\begin{equation}
  \mathbf{x}=\{\mathbf{T}_0,\boldsymbol{\varpi}_0,\dots,\mathbf{T}_K,\boldsymbol{\varpi}_K,\ \mathbf{p}_1,\dots,\mathbf{p}_M\}.
  \label{eq:steam-state}
\end{equation}
每个观测 $\mathbf{y}_{jk}=\mathbf{s}(\mathbf{T}_k\mathbf{p}_j)+\mathbf{n}_{jk}$（$\mathbf{s}$ 为非线性传感器模型，如立体相机；$\mathbf{n}_{jk}\sim\mathcal{N}(\mathbf{0},\boldsymbol{\Sigma}_{\mathbf{v},jk})$），与离散 SLAM 完全相同。线性化观测误差 $\mathbf{e}_{y,jk}(\mathbf{x}_{\mathrm{op}})=\mathbf{y}_{jk}-\mathbf{s}(\mathbf{T}_{\mathrm{op},k}\mathbf{p}_{\mathrm{op},j})$，其对位姿/速度块 $\delta\mathbf{x}_1$ 与对路标块 $\delta\mathbf{x}_2$ 的雅可比为
\begin{equation}
  \mathbf{G}_{1,jk}=\big[\;\mathbf{S}_{jk}(\mathbf{T}_{\mathrm{op},k}\mathbf{p}_{\mathrm{op},j})^\odot\quad \mathbf{0}\;\big],\qquad
  \mathbf{G}_{2,jk}=\mathbf{S}_{jk}\mathbf{T}_{\mathrm{op},k}\mathbf{D},
  \label{eq:steam-G}
\end{equation}
$\mathbf{S}_{jk}=\partial\mathbf{s}/\partial(\mathbf{T}\mathbf{p})$ 为传感器雅可比，$(\cdot)^\odot$ 是把齐次点变为 $4\times6$ 矩阵的算子（\cref{ch:lie} 的 $\partial(\mathbf{T}\tilde{\mathbf{p}})/\partial\delta\boldsymbol{\xi}=(\mathbf{T}\tilde{\mathbf{p}})^\odot$），$\mathbf{D}$ 为齐次坐标的 $4\times3$ 投影。\textbf{与离散 SLAM 唯一的差别}：$\mathbf{G}_{1,jk}$ 末尾多出一个零块——因为速度 $\boldsymbol{\varpi}_k$ \emph{不参与}对路标的观测（观测只看位姿与路标）。把先验与观测拼起来，
\begin{equation}
  J(\mathbf{x})=J_v(\mathbf{x})+J_y(\mathbf{x})\approx J(\mathbf{x}_{\mathrm{op}})-\mathbf{b}^\top\delta\mathbf{x}+\tfrac12\delta\mathbf{x}^\top\mathbf{A}\,\delta\mathbf{x},
  \label{eq:steam-J}
\end{equation}
\begin{equation}
  \mathbf{A}=\mathbf{H}^\top\mathbf{W}^{-1}\mathbf{H},\quad
  \mathbf{H}=\begin{bmatrix}\mathbf{F}^{-1}&\mathbf{0}\\\mathbf{G}_1&\mathbf{G}_2\end{bmatrix},\quad
  \mathbf{W}=\begin{bmatrix}\mathbf{Q}&\mathbf{0}\\\mathbf{0}&\boldsymbol{\Sigma}_{\mathbf{v}}\end{bmatrix},
  \label{eq:steam-A}
\end{equation}
正规方程 $\mathbf{A}\,\delta\mathbf{x}^\star=\mathbf{b}$ 与本章 \cref{eq:normal-eq} \emph{同形}。解出后按各自规则更新——位姿右乘 $\mathbf{T}_{\mathrm{op},k}\leftarrow\exp(\boldsymbol{\epsilon}_k^{\star\wedge})\mathbf{T}_{\mathrm{op},k}$、速度加性 $\boldsymbol{\varpi}_{\mathrm{op},k}\leftarrow\boldsymbol{\varpi}_{\mathrm{op},k}+\boldsymbol{\eta}_k^\star$、路标 $\mathbf{p}_{\mathrm{op},j}\leftarrow\mathbf{p}_{\mathrm{op},j}+\mathbf{D}\boldsymbol{\zeta}_j^\star$——迭代至收敛。\emph{运动先验的可解性意义}：真实轨迹中常出现某位姿\emph{只观测到一个}路标（如螺旋下降时每帧只瞥见一个特征），仅凭这一个观测无法定姿；此时正是 WNOA 运动先验（\cref{eq:steam-error} 的二元因子）把相邻位姿“拉”在常速度轨迹上、补足缺失的约束，才使整个 STEAM 问题可解——这与 \cref{sec:est-observability} “先验补足可观测性、消除 gauge 自由度”是同一回事。

\begin{proposition}[STEAM 不破坏 SLAM 的稀疏结构]\label{prop:steam-sparse}
按位姿/速度（块 1）与路标（块 2）分块，\cref{eq:steam-A} 的 $\mathbf{A}$ 仍是\emph{箭头矩阵}，且
\begin{equation}
  \mathbf{A}=\begin{bmatrix}\mathbf{A}_{11}&\mathbf{A}_{12}\\\mathbf{A}_{12}^\top&\mathbf{A}_{22}\end{bmatrix}
  =\begin{bmatrix}\mathbf{F}^{-\top}\mathbf{Q}^{-1}\mathbf{F}^{-1}+\mathbf{G}_1^\top\boldsymbol{\Sigma}_{\mathbf{v}}^{-1}\mathbf{G}_1&\mathbf{G}_1^\top\boldsymbol{\Sigma}_{\mathbf{v}}^{-1}\mathbf{G}_2\\\mathbf{G}_2^\top\boldsymbol{\Sigma}_{\mathbf{v}}^{-1}\mathbf{G}_1&\mathbf{G}_2^\top\boldsymbol{\Sigma}_{\mathbf{v}}^{-1}\mathbf{G}_2\end{bmatrix},
  \label{eq:steam-arrow}
\end{equation}
其中路标块 $\mathbf{A}_{22}=\mathbf{G}_2^\top\boldsymbol{\Sigma}_{\mathbf{v}}^{-1}\mathbf{G}_2$ \emph{块对角}（各路标互不直接耦合），位姿/速度块 $\mathbf{A}_{11}=\underbrace{\mathbf{F}^{-\top}\mathbf{Q}^{-1}\mathbf{F}^{-1}}_{\text{先验，块三对角}}+\underbrace{\mathbf{G}_1^\top\boldsymbol{\Sigma}_{\mathbf{v}}^{-1}\mathbf{G}_1}_{\text{观测}}$ \emph{块三对角}。
\end{proposition}
\begin{proof}
$\mathbf{A}_{12},\mathbf{A}_{22}$ 的形式与离散 SLAM（\cref{ch:nlopt} 的 BA 箭头矩阵）逐字相同：路标只被各自的观测因子触及，故 $\mathbf{A}_{22}$ 块对角、$\mathbf{A}_{12}$ 稀疏。$\mathbf{A}_{11}$ 的先验项 $\mathbf{F}^{-\top}\mathbf{Q}^{-1}\mathbf{F}^{-1}$ 由 \cref{eq:steam-Finv} 的块下双对角 $\mathbf{F}^{-1}$ 与块对角 $\mathbf{Q}^{-1}$ 相乘，结果块三对角（块下双对角阵的“$\mathbf{B}^\top\mathbf{B}$”型乘积只在主对角与相邻对角非零）；观测项 $\mathbf{G}_1^\top\boldsymbol{\Sigma}_{\mathbf{v}}^{-1}\mathbf{G}_1$ 块对角（每个观测只连一个位姿）。二者之和仍块三对角。\cref{eq:steam-G} 末尾的零块不引入新耦合，故结构与离散 SLAM 一致。
\end{proof}

\begin{insight}[因此 STEAM 求解与离散 SLAM 同样高效]
\cref{prop:steam-sparse} 是 STEAM 实用性的关键：连续时间先验\emph{没有}让矩阵变稠密。于是本章/\,\cref{ch:nlopt} 的整套稀疏武器照搬——若路标多，先 Schur 消元路标块 $\mathbf{A}_{22}$（块对角，逆廉价）得位姿/速度上的约简系统；若位姿/速度多，则直接对块三对角的 $\mathbf{A}_{11}$ 做稀疏 Cholesky（此时优于 Schur，因无需构造稠密的 $\mathbf{A}_{11}^{-1}$）。把稀疏 GP 先验嵌入 iSAM 的增量框架（\cref{sec:est-smooth-filter} 的 Bayes 树）也已实现\cite{anderson2015steam}。一句话：\emph{STEAM 把“连续时间查询”的能力\textbf{免费}加到了离散 SLAM 上，渐近复杂度不变}。
\end{insight}

\subsection{任意时刻查询：GP 插值与外推}
\label{sec:est-steam-query}

连续时间方法相对离散的\emph{核心卖点}在此：解完测量时刻的状态后，可查询\emph{任意}时刻 $\tau$ 的位姿、速度及协方差，单次查询 $O(1)$。

\paragraph{均值插值（$t_k\le\tau<t_{k+1}$）。}
因状态马尔可夫，查询 $\tau$ \emph{只需}两端 $t_k,t_{k+1}$。GP 后验插值（向量空间通式，\cref{ch:nlopt} 附录）给出局部状态
\begin{equation}
  \hat{\boldsymbol{\gamma}}_k(\tau)=\boldsymbol{\Lambda}(\tau)\,\hat{\boldsymbol{\gamma}}_k(t_k)+\boldsymbol{\Psi}(\tau)\,\hat{\boldsymbol{\gamma}}_k(t_{k+1}),
  \label{eq:steam-interp}
\end{equation}
其中插值权由先验的 $\boldsymbol{\Phi},\mathbf{Q}$ 组成（$\mathbf{Q}_\tau=\mathbf{Q}(\tau-t_k)$、$\mathbf{Q}_{k+1}=\mathbf{Q}(t_{k+1}-t_k)$）：
\begin{equation}
  \boldsymbol{\Lambda}(\tau)=\boldsymbol{\Phi}(\tau,t_k)-\mathbf{Q}_\tau\,\boldsymbol{\Phi}(t_{k+1},\tau)^\top\mathbf{Q}_{k+1}^{-1}\boldsymbol{\Phi}(t_{k+1},t_k),\qquad
  \boldsymbol{\Psi}(\tau)=\mathbf{Q}_\tau\,\boldsymbol{\Phi}(t_{k+1},\tau)^\top\mathbf{Q}_{k+1}^{-1}.
  \label{eq:steam-LambdaPsi}
\end{equation}
把局部量按 \cref{eq:steam-local} 换回全局位姿（$\boldsymbol{\Lambda}=[\boldsymbol{\Lambda}_1;\boldsymbol{\Lambda}_2]$、$\boldsymbol{\Psi}=[\boldsymbol{\Psi}_1;\boldsymbol{\Psi}_2]$ 按位姿/速度分块）：
\begin{equation}
  \hat{\mathbf{T}}(\tau)=\exp\!\Big(\big(\boldsymbol{\Lambda}_1(\tau)\hat{\boldsymbol{\gamma}}_k(t_k)+\boldsymbol{\Psi}_1(\tau)\hat{\boldsymbol{\gamma}}_k(t_{k+1})\big)^\wedge\Big)\hat{\mathbf{T}}_k,
  \label{eq:steam-Tquery}
\end{equation}
速度 $\hat{\boldsymbol{\varpi}}(\tau)$ 由 \cref{eq:steam-local-sde} 的 $\boldsymbol{\varpi}=\boldsymbol{\mathcal{J}}_r(\boldsymbol{\xi})\boldsymbol{\psi}$ 关系恢复。在常速度（WNOA）先验下，\cref{eq:steam-interp} 恰好退化为\emph{三次 Hermite 多项式插值}（位置三次、速度二次）——这给了“为什么连续时间查询常被实现成三次样条”一个\emph{先验层面}的解释。

\paragraph{协方差插值。}
均值之外还要查询时刻的不确定度。收敛后正规方程左端 $\mathbf{A}=\hat{\mathbf{P}}^{-1}$ 是全轨迹估计的逆协方差（拉普拉斯近似，\cref{sec:nlopt-cov}）。查询 $\tau$ 的协方差\emph{只需}两端状态的边缘块
\begin{equation}
  \mathrm{cov}(\delta\mathbf{x}_{k:k+1}^\star)=\begin{bmatrix}\hat{\mathbf{P}}(t_k,t_k)&\hat{\mathbf{P}}(t_k,t_{k+1})\\\hat{\mathbf{P}}(t_{k+1},t_k)&\hat{\mathbf{P}}(t_{k+1},t_{k+1})\end{bmatrix},
  \label{eq:steam-cov-bracket}
\end{equation}
它可\emph{利用箭头结构高效抽取}——这正是\emph{选择性稀疏逆}（selected inverse / Takahashi 方程）的用武之地：只回收 $\hat{\mathbf{P}}=\boldsymbol{\Lambda}^{-1}$ 中\emph{需要的}块，绝不显式求整个稠密逆（\cref{sec:nlopt-cov} 的“按需回代特定块”即此，连续时间下因要反复查询而尤为关键）。拿到两端边缘后，把原本连接 $t_k,t_{k+1}$ 的单个先验因子\emph{替换}为经过 $\tau$ 的两段先验因子（\cref{eq:steam-FE} 取相应时间差），对仅含 $\{t_k,\tau,t_{k+1}\}$ 三状态的小问题装配 $36\times36$ 逆协方差——它由\emph{三部分}相加：(i) 主解给的两端边缘当先验、(ii) 经 $\tau$ 串联的两段新先验因子、(iii) \emph{减去}原来直连 $t_k,t_{k+1}$ 的那个先验因子（它已含在两端边缘内，须扣除以免重复计入）：
\begin{equation}
  \hat{\mathbf{P}}_{\mathrm{interp}}^{-1}=
  \underbrace{\mathbf{M}_1^\top\!\begin{bmatrix}\hat{\mathbf{P}}(t_k,t_k)&\hat{\mathbf{P}}(t_k,t_{k+1})\\\hat{\mathbf{P}}(t_{k+1},t_k)&\hat{\mathbf{P}}(t_{k+1},t_{k+1})\end{bmatrix}^{-1}\!\mathbf{M}_1}_{\text{(i) 两端边缘作先验}}
  +\underbrace{\mathbf{M}_2^\top\!\begin{bmatrix}\mathbf{Q}_{\tau,k}&\mathbf{0}\\\mathbf{0}&\mathbf{Q}_{k+1,\tau}\end{bmatrix}^{-1}\!\mathbf{M}_2}_{\text{(ii) 经 }\tau\text{ 的两段新因子}}
  -\underbrace{\begin{bmatrix}-\mathbf{F}_k^\top\\\mathbf{0}\\\mathbf{E}_{k+1}^\top\end{bmatrix}\!\mathbf{Q}_{k+1}^{-1}\!\begin{bmatrix}-\mathbf{F}_k&\mathbf{0}&\mathbf{E}_{k+1}\end{bmatrix}}_{\text{(iii) 扣除原直连因子}},
  \label{eq:steam-cov-assembly}
\end{equation}
其中选择阵 $\mathbf{M}_1=\big[\begin{smallmatrix}\mathbf{I}&\mathbf{0}&\mathbf{0}\\\mathbf{0}&\mathbf{0}&\mathbf{I}\end{smallmatrix}\big]$（按 $t_k,t_{k+1}$ 取块、跳过 $\tau$），两段因子雅可比 $\mathbf{M}_2=\big[\begin{smallmatrix}-\mathbf{F}_{\tau,k}&\mathbf{E}_{\tau,k}&\mathbf{0}\\\mathbf{0}&-\mathbf{F}_{k+1,\tau}&\mathbf{E}_{k+1,\tau}\end{smallmatrix}\big]$（$\mathbf{F}_{\tau,k},\mathbf{E}_{\tau,k}$ 同 \cref{eq:steam-FE} 取时间 $\tau,k$，$\mathbf{Q}_{\tau,k}=\mathbf{Q}(\tau-t_k),\ \mathbf{Q}_{k+1,\tau}=\mathbf{Q}(t_{k+1}-\tau)$）。由于均值已在上一段插得，这里只需装好这一次高斯牛顿的左端逆协方差（拉普拉斯近似），无须再迭代。再对 $t_k,t_{k+1}$ 两块做\emph{边缘化}（信息形式舒尔补，\cref{eq:info-marginal}），剩 $\tau$ 块即所求：
\begin{equation}
  \hat{\mathbf{P}}(\tau,\tau)^{-1}
  =\mathbf{E}_{\tau,k}^\top\mathbf{Q}_{\tau,k}^{-1}\mathbf{E}_{\tau,k}+\mathbf{F}_{k+1,\tau}^\top\mathbf{Q}_{k+1,\tau}^{-1}\mathbf{F}_{k+1,\tau}
  -\mathbf{U}^\top\boldsymbol{\Xi}^{-1}\mathbf{U},
  \label{eq:steam-cov-interp}
\end{equation}
其中 $\tau$ 与两端的耦合块 $\mathbf{U}=\big[\begin{smallmatrix}\mathbf{F}_{\tau,k}^\top\mathbf{Q}_{\tau,k}^{-1}\mathbf{E}_{\tau,k}\\\mathbf{E}_{k+1,\tau}^\top\mathbf{Q}_{k+1,\tau}^{-1}\mathbf{F}_{k+1,\tau}\end{smallmatrix}\big]$，被边缘化的 $24\times24$ 两端块 $\boldsymbol{\Xi}$ 是 \cref{eq:steam-cov-assembly} 中 $\{t_k,t_{k+1}\}$ 的主子块（即两端边缘逆、两段新因子的对角贡献 $\mathrm{diag}(\mathbf{F}_{\tau,k}^\top\mathbf{Q}_{\tau,k}^{-1}\mathbf{F}_{\tau,k},\ \mathbf{E}_{k+1,\tau}^\top\mathbf{Q}_{k+1,\tau}^{-1}\mathbf{E}_{k+1,\tau})$、以及扣除项三者之和）。求逆 $\hat{\mathbf{P}}(\tau,\tau)=(\,\cdots)^{-1}$ 即得查询时刻协方差，单次 $O(1)$。\cref{eq:steam-cov-interp} 与已全式给出的\emph{外推}协方差 \cref{eq:steam-extrap-cov} 现在深度对称——插值是“两端夹 $\tau$”、外推是“单端悬臂 $\tau$”，同一套“先验因子装配 $+$ 边缘化”机理。

\paragraph{外推（预测未来态，$\tau>t_K$）。}
比插值更简单：以最后一态“悬臂”外推。均值直接用常速度先验
\begin{equation}
  \hat{\mathbf{T}}(\tau)=\exp\!\big((\tau-t_K)\,\hat{\boldsymbol{\varpi}}(t_K)^\wedge\big)\hat{\mathbf{T}}(t_K),\qquad
  \hat{\boldsymbol{\varpi}}(\tau)=\hat{\boldsymbol{\varpi}}(t_K).
  \label{eq:steam-extrap}
\end{equation}
协方差用仅含 $\{t_K,\tau\}$ 的两状态小问题，边缘化 $t_K$、再用 SMW（\cref{lem:smw}）化简，得
\begin{equation}
  \hat{\mathbf{P}}(\tau,\tau)=\mathbf{E}_{\tau,K}^{-1}\big(\mathbf{F}_{\tau,K}\,\hat{\mathbf{P}}(t_K,t_K)\,\mathbf{F}_{\tau,K}^\top+\mathbf{Q}_{\tau,K}\big)\mathbf{E}_{\tau,K}^{-\top},
  \label{eq:steam-extrap-cov}
\end{equation}
$\mathbf{F}_{\tau,K},\mathbf{E}_{\tau,K}$ 同 \cref{eq:steam-FE}（时间取 $\tau,K$）、$\mathbf{Q}_{\tau,K}=\mathbf{Q}(\tau-t_K)$。\cref{eq:steam-extrap-cov} 形如 EKF 的协方差\emph{预测步} $\check{\boldsymbol{\Sigma}}=\mathbf{F}\hat{\boldsymbol{\Sigma}}\mathbf{F}^\top+\boldsymbol{\Sigma}_{\mathbf{w}}$（\cref{thm:kf-from-map}），但适配到 $\SE(3)\times\mathbb{R}^6$ 状态——连续时间外推与滤波预测在此漂亮地打通。

\begin{insight}[三套时间戳可解耦：测量时刻 $\ne$ 估计时刻 $\ne$ 查询时刻]
既然能插值，主解里的状态数就\emph{不必}等于测量数。可让三套时间戳各司其职：(i) \emph{测量}时刻（数据到来，通常很多）、(ii) \emph{估计}时刻（主 MAP 实际求解的状态，可少）、(iii) \emph{查询}时刻（事后要取的位姿，可任意多）。测量时刻的因子通过插值“挂”到相邻估计状态上，从而\emph{大幅压缩主解规模}。这正是连续时间方法应对高频/异步传感器的威力之源，也与 \cref{ch:vio} 里“用连续轨迹融合 IMU 与相机”一脉相承。
\end{insight}

\begin{pitfall}[把约束式插值误当成 GP 先验插值]
\textbf{错误想法}：“连续时间无非就是在两关键帧之间插个值，跟 \cref{ch:vio} 里‘强制匀速’的几何插值一回事。”\textbf{现象}：把先验当成\emph{硬约束}（强制段内恒定广义速度），于是先验既不可调权重、也无协方差可言，查询协方差无从谈起。\textbf{根本原因}：二者范式不同。约束式插值\emph{规定}轨迹形状（一条确定曲线）；STEAM 的 GP 先验\emph{定义一整个轨迹分布}，用代价项\emph{鼓励}（而非强制）解贴近常速度，并与观测\emph{权衡}——故先验有 PSD $\mathbf{Q}_C$ 可调、查询有协方差可算。\textbf{正确做法}：把运动先验当\emph{带权重的因子}（\cref{eq:steam-error} 的 $\mathbf{Q}_k^{-1}$ 即权），PSD 小则先验强（轨迹更平滑、更不信观测），PSD 大则先验弱。\textbf{自检}：你的“插值”能不能输出协方差？能，才是 GP 先验插值。
\end{pitfall}

\begin{insight}[同一框架的高阶与跨域推广]
本节用 WNOA（白噪声加速度）讲透了机理，但这\emph{不是}唯一选择，框架可直接换更高阶或物理化的 SDE。\textbf{(1) WNOJ（白噪声加加速度，white-noise-on-jerk）}\cite{tang2019wnoj}：把白噪声注入\emph{三阶}导（$\dddot{\boldsymbol{\xi}}=\mathbf{w}$），先验鼓励\emph{常加速度}，状态扩为 $[\boldsymbol{\xi};\boldsymbol{\psi};\dot{\boldsymbol{\psi}}]$、$\mathbf{A}$ 带宽加一，仍块带状、仍 $O(K)$。\textbf{注意}：WNOA 因 $\mathbf{A}^2=\mathbf{0}$ 幂零才可手算出 \cref{eq:steam-Phi-Q} 的闭式 $\boldsymbol{\Phi},\mathbf{Q}$；WNOJ/Singer 的 $\mathbf{A}$ \emph{不再幂零}，此时 $\boldsymbol{\Phi},\mathbf{Q}$ 应由 \cref{thm:matderiv-vanloan}（Van Loan 矩阵指数法）统一机械算出（该附录定理正是为此类连续$\to$离散化而设，并显式核对与本节幂零手算在 WNOA 上一致）。\textbf{(2) Singer 模型}\cite{wong2020singer}：在加速度上加一阶马尔可夫衰减（介于常速度与常加速度之间，可由数据拟合衰减时间常数），更贴合真实运动谱。\textbf{(3) 连续体机器人}\cite{lilge2022continuum}：把“时间”换成“弧长”，本节几乎原样搬到连续体机器人的状态估计——其简化的 Cosserat 杆先验与 WNOA 轨迹先验在数学上高度相似。这些都印证了连续时间 GP 框架的统一性：换 SDE 只换 $\mathbf{A},\mathbf{L},\boldsymbol{\Phi},\mathbf{Q}$，整套稀疏求解与查询机理不变。
\end{insight}

\begin{practice}
\begin{enumerate}
  \item （推导）验证 \cref{eq:steam-Phi-Q}：由 $\mathbf{A}^2=\mathbf{0}$ 证 $\boldsymbol{\Phi}(t,s)=\mathbf{I}+\mathbf{A}(t-s)$，并对 $\mathbf{Q}(\Delta t)$ 的积分逐块算出 $\tfrac13\Delta t^3,\tfrac12\Delta t^2,\Delta t$ 三个系数。
  \item （概念）解释 \cref{eq:steam-error} 第二行为何要乘 $\boldsymbol{\mathcal{J}}_r^{-1}$：从 \cref{thm:se3-kinematics} 的 $\boldsymbol{\psi}=\boldsymbol{\mathcal{J}}_r^{-1}\boldsymbol{\varpi}$ 出发，说明局部坐标率与机体系广义速度之别。
  \item 由 \cref{prop:steam-sparse} 说明：当路标数 $\gg$ 位姿数时为何 Schur 消元路标更优；当位姿/速度数 $\gg$ 路标数时为何反而该直接对 $\mathbf{A}_{11}$ 做稀疏 Cholesky（提示：$\mathbf{A}_{11}^{-1}$ 稠密、$\mathbf{A}_{22}$ 块对角）。
  \item （连接 \cref{ch:lie}）证明常速度（WNOA）先验下 \cref{eq:steam-interp} 退化为三次 Hermite 插值：取 $\boldsymbol{\Phi},\mathbf{Q}$ 的标量化形式，验证位置关于 $\tau$ 是三次、速度是二次。
  \item （跨章综合）把 \cref{eq:steam-extrap-cov} 与 \cref{thm:kf-from-map} 的 EKF 协方差预测 $\check{\boldsymbol{\Sigma}}=\mathbf{F}\hat{\boldsymbol{\Sigma}}\mathbf{F}^\top+\boldsymbol{\Sigma}_{\mathbf{w}}$ 对照，指出 $\mathbf{E}_{\tau,K},\mathbf{F}_{\tau,K},\mathbf{Q}_{\tau,K}$ 各自对应 KF 里的什么量。
  \item （概念，连接 \cref{sec:nlopt-cov}）为什么连续时间下“查询协方差”离不开\emph{选择性稀疏逆}（Takahashi）？若每次查询都显式求 $\boldsymbol{\Lambda}^{-1}$ 全阵，复杂度会怎样退化？
\end{enumerate}
\end{practice}

\subsection{另一条路：参数样条与 GP 的取舍}
\label{sec:est-steam-spline}

本节开头说连续时间有\emph{两条}路；上面把 GP（非参数）一支讲到了底。这里补全\emph{参数}一支——样条，并点明二者的本质区别。两者都把轨迹写成连续函数、都用局部支撑换稀疏，差别只在“正则从哪来”。

\paragraph{参数样条：位姿 $=$ 已知基函数的加权和。}
把轨迹写成有限个\emph{控制点} $\mathbf{c}_k$ 与\emph{已知}时间基函数 $\psi_k(t)$ 的加权和\cite{carlone2026handbook}（向量空间先讲，李群版见下）
\begin{equation}
  \mathbf{p}(t)=\sum_{k=1}^{K}\psi_k(t)\,\mathbf{c}_k,\qquad
  \dot{\mathbf{p}}(t)=\sum_{k=1}^{K}\dot{\psi}_k(t)\,\mathbf{c}_k .
  \label{eq:steam-spline-basis}
\end{equation}
待估的不再是“每个时刻的状态”，而是这\emph{一组系数}。基函数取 B 样条等\emph{局部支撑}函数：在某时刻 $t$ 只有少数几个 $\psi_k(t)\ne0$，于是一条观测 $\mathbf{z}_i=h_i(\mathbf{p}(t_i),\boldsymbol{\ell})+\mathbf{n}_i$ 只触及那几个邻近控制点（再加它看到的路标），化为一个\emph{小定义域因子}——因子图照样稀疏。速度、加速度由\emph{对基函数求导} \cref{eq:steam-spline-basis} 直接得到，故依赖速度的传感器（如 IMU）也能用同一组系数表达。线性化、白化、稀疏分解与离散 SLAM\emph{逐字相同}，只是变量从“位姿”换成“系数”。

\paragraph{累积形式与李群样条。}
直接用 \cref{eq:steam-spline-basis} 把李群元素“加权求和”没有意义（$\SE(3)$ 不是向量空间）。出路是\emph{累积}（cumulative）形式：先锚在首控制点，再叠加相邻控制点之\emph{差}
\begin{equation}
  \mathbf{p}(t)=\psi_1^c(t)\,\mathbf{c}_1+\sum_{k=2}^{K}\psi_k^c(t)\,(\mathbf{c}_k-\mathbf{c}_{k-1}),
  \qquad \psi_k^c(t)=\sum_{\ell=k}^{K}\psi_\ell(t),
  \label{eq:steam-spline-cumul}
\end{equation}
其中累积基 $\psi_k^c$ 是普通基的\emph{尾和}。把“求和”换成李群乘法、把“差”换成\emph{相对位姿} $\mathbf{T}_k\mathbf{T}_{k-1}^{-1}$、用 $\exp/\ln$ 进出代数，即得 $\SE(3)$ 上的累积样条。以一阶（线性）样条为例，当 $\tau\in[t_{k-1},t_k)$、记归一时间 $\alpha=\tfrac{\tau-t_{k-1}}{t_k-t_{k-1}}\in[0,1)$、$\psi_k^c(\tau)=\alpha$：
\begin{equation}
  \mathbf{T}(\tau)=\exp\!\Big(\big(\alpha\,\ln(\mathbf{T}_k\mathbf{T}_{k-1}^{-1})^\vee\big)^{\wedge}\Big)\,\mathbf{T}_{k-1}
  =\big(\mathbf{T}_k\mathbf{T}_{k-1}^{-1}\big)^{\alpha}\,\mathbf{T}_{k-1},
  \label{eq:steam-spline-lie}
\end{equation}
即“沿相对位姿走比例 $\alpha$”。高阶样条只是激活更多控制点、累积基更光滑，结构不变。

\begin{note}[扰动取在哪一侧]
\cref{eq:steam-spline-lie} 把插值增量\emph{左乘}在 $\mathbf{T}_{k-1}$ 上。部分文献（含 Handbook\cite{carlone2026handbook}、Barfoot\cite{barfoot2024state}）默认在\emph{传感器系}做插值、故取\emph{左}扰动；本书统一\emph{右}扰动主线（\cref{ch:lie}），二者经伴随 $\boldsymbol{\mathcal{T}}=\mathrm{Ad}(\mathbf{T})$ 精确互换，不影响结论。沿用本书右扰动 $\mathbf{T}_k=\mathbf{T}_{\mathrm{op},k}\,\mathrm{Exp}(\boldsymbol{\epsilon}_k)$ 时，把控制点扰动 $\boldsymbol{\epsilon}_{k-1},\boldsymbol{\epsilon}_k$ 传到插值位姿扰动 $\boldsymbol{\epsilon}(\tau)$ 的一阶关系形如
\begin{equation}
  \boldsymbol{\epsilon}(\tau)\approx\big(\mathbf{I}-\boldsymbol{\mathcal{A}}(\alpha)\big)\boldsymbol{\epsilon}_{k-1}+\boldsymbol{\mathcal{A}}(\alpha)\,\boldsymbol{\epsilon}_k,
  \qquad \boldsymbol{\mathcal{A}}(\alpha)=\alpha\,\boldsymbol{\mathcal{J}}\!\big(\alpha\,\boldsymbol{\phi}_{k,k-1}\big)\,\boldsymbol{\mathcal{J}}\!\big(\boldsymbol{\phi}_{k,k-1}\big)^{-1},
  \label{eq:steam-spline-jac}
\end{equation}
$\boldsymbol{\phi}_{k,k-1}=\ln(\mathbf{T}_{\mathrm{op},k}\mathbf{T}_{\mathrm{op},k-1}^{-1})^\vee$、$\boldsymbol{\mathcal{J}}$ 为 $\SE(3)$ 雅可比（\cref{ch:lie}）。把 \cref{eq:steam-spline-jac} 代进观测残差的链式法则，就得到样条因子对两端控制点的雅可比——与 STEAM 把误差线性化到相邻状态（\cref{eq:steam-lin}）的做法同构。
\end{note}

\begin{insight}[样条 vs GP：唯一的本质区别是“有没有运动先验因子”]
\cref{eq:steam-spline-lie} 与 GP 在常速度先验下的查询式 \cref{eq:steam-interp}/\cref{eq:steam-Tquery} \emph{形式相同}——都是“在两个 $\SE(3)$ 控制点之间做（李群）线性插值”。区别不在插值，而在\emph{解出控制点的目标里有什么}：
\begin{itemize}
  \item \textbf{样条}：控制点只被\emph{观测}约束。光滑性来自基函数的选择（阶数、节点密度）。控制点太多易\emph{过拟合}、太少则\emph{过平滑}，且常需额外加一个正则项 $\|\mathbf{c}\|^2$ 压描述长度\cite{carlone2026handbook}。
  \item \textbf{GP（STEAM）}：除观测外，相邻估计时刻间还有一族\emph{运动先验因子}（\cref{eq:steam-error}，权重为 $\mathbf{Q}_k^{-1}$）。这族因子\emph{就是}正则项，由 SDE（如 WNOA）自动给出、可调 PSD，并\emph{附带协方差}——故 GP 能算查询不确定度而纯样条不能（\cref{sec:est-steam-query} 的协方差插值无对应物）。
\end{itemize}
一句话：\emph{GP $=$ 样条 $+$ 一族自洽的运动先验因子}\cite{carlone2026handbook}。这解释了本书为何以 GP 为主线（运动先验与观测在同一 MAP 里权衡、查询有协方差）；而 \cref{eq:steam-spline-jac} 表明，若已用样条，也可借同样的“控制点插值”把多个测量时刻\emph{挂}到少数控制点上，与 GP 用插值减少估计时刻（inducing points，\cref{sec:est-steam-query} 的三套时间戳）异曲同工。需要说明的是，二者并非互斥——给样条补上运动先验因子，也能获得类似的正则\cite{carlone2026handbook}。
\end{insight}

% ════════════════════════════════════════════════════════════════
\section{实践：把最小二乘交给求解器}
\label{sec:est-practice}

理论落地时，你几乎不手写正规方程，而是\emph{描述因子、交给库}。三类主流工具：稀疏线性代数（Eigen/CHOLMOD 直接解 \cref{eq:normal-eq}）、通用非线性最小二乘（Ceres）、因子图专用（GTSAM、g2o）。本节给最小可跑的骨架，求解器内部（GN/LM、稀疏分解、变量排序）留 \cref{ch:nlopt}。

\paragraph{信息矩阵与正规方程（Eigen）。}
线性情形可直接组装 $\boldsymbol{\Lambda}=\mathbf{H}^\top\mathbf{W}^{-1}\mathbf{H}$ 并用 Cholesky 解。

\begin{codebox}[C++]{Eigen:组装并求解正规方程 $\boldsymbol{\Lambda}\hat{\mathbf{x}}=\boldsymbol{\eta}$}
#include <Eigen/Dense>
using namespace Eigen;
// H: m x n 测量雅可比;Winv: m x m 信息(逆协方差，块对角);z: m 观测
MatrixXd Lambda = H.transpose() * Winv * H;   // 信息矩阵 Lambda = H^T W^-1 H
VectorXd eta    = H.transpose() * Winv * z;    // 信息向量 eta = H^T W^-1 z
VectorXd x_hat  = Lambda.ldlt().solve(eta);    // Cholesky(LDL^T) 解，利用对称正定
MatrixXd Sigma  = Lambda.inverse();            // 解的协方差 = Lambda^-1（大规模时勿显式求逆）
\end{codebox}

\noindent 工程提醒：大规模问题\emph{绝不}显式求 $\boldsymbol{\Lambda}^{-1}$（稠密、$O(n^3)$）；用稀疏 Cholesky 解线性系统，按需回代特定块的边际协方差（\cref{ch:nlopt} 的稀疏求解）。数值稳健时常用\emph{平方根}信息形式（$\sqrt{}$SAM，QR/Cholesky 因子，见 \cref{sec:est-fg-linear}）。

\paragraph{因子图（GTSAM 风格伪代码）。}
现代写法是声明变量与因子，让库去线性化与求解：

\begin{codebox}[C++]{GTSAM 风格:把 \cref{fig:factor-graph} 建成因子图}
NonlinearFactorGraph graph;
// 先验因子:锚定首位姿，消除 gauge 自由度（对应 tab:factors 先验项）
graph.addPrior(X(1), prior_pose, priorNoise);
// 里程计因子:连接相邻位姿（BetweenFactor 对应里程计/回环）
graph.emplace_shared<BetweenFactor<Pose>>(X(1), X(2), odo_12, odoNoise);
graph.emplace_shared<BetweenFactor<Pose>>(X(2), X(3), odo_23, odoNoise);
// 观测因子:位姿看到路标（对应重投影/方位残差）
graph.emplace_shared<BearingRangeFactor<Pose,Point>>(X(1), L(1), meas_11, obsNoise);
// ... 给初值，交给优化器（GN/LM 见 nlopt 章）
Values result = LevenbergMarquardtOptimizer(graph, initial).optimize();
Marginals marginals(graph, result);   // 取边际协方差 = Lambda^-1 的对应块
\end{codebox}

\noindent 注意 \texttt{addPrior} 对应 \cref{tab:factors} 的先验因子、\texttt{BetweenFactor} 对应里程计/回环、观测因子对应重投影/方位——与本章理论一一对应。\texttt{noise} 模型就是各因子的 $\boldsymbol{\Sigma}_i$。库内部正是在做本章的“因子图 $\to$ 线性化 $\to$ 变量消元/稀疏分解 $\to$ 回代”（\cref{thm:elimination}）。

\begin{practice}
\begin{enumerate}
  \item 用 Eigen 实现一维小车例（\cref{eq:car-solution}），组装 $\mathbf{H},\boldsymbol{\Sigma}$ 并解出 $\hat{\mathbf{x}}$ 与 $\hat{\boldsymbol{\Sigma}}$；验证去掉先验因子后 $\boldsymbol{\Lambda}$ 奇异。
  \item 在 GTSAM/g2o 任一库里复现 \cref{fig:factor-graph}（含一条回环因子），对比有/无回环的轨迹误差，并用 \texttt{Marginals} 取出某位姿的边际协方差。
  \item （跨章综合）对同一因子图，分别用“信息矩阵 Cholesky”（\cref{sec:est-fg-linear}）与“变量消元”（\cref{thm:elimination}）求解，验证两者给出相同 $\hat{\mathbf{x}}$，并比较自然排序与 COLAMD 排序下因子的非零元数。
\end{enumerate}
\end{practice}

% ════════════════════════════════════════════════════════════════
\section*{常见误解汇总}
\addcontentsline{toc}{section}{常见误解汇总}
\begin{table}[H]\centering\small
\caption{本章常见误解与正确理解}
\begin{tabular}{p{0.42\linewidth}p{0.5\linewidth}}
\toprule
误解 & 正确理解 \\
\midrule
“SLAM 是解方程组。” & 观测多于未知且带噪，精确解一般不存在；是\emph{概率估计}（\cref{sec:est-why}）。 \\
“MLE 和 MAP 一样。” & 仅当先验均匀时同解；一般 MAP $=$ MLE $+$ 先验（\cref{sec:est-mle-map}）。 \\
“最小二乘只是工程凑合。” & 高斯下它\emph{精确}等于 MAP（\cref{thm:map-ls}）；非高斯下它是 BLUE（高斯–马尔可夫）。 \\
“协方差大 $=$ 没信息。” & 看信息矩阵 $\boldsymbol{\Lambda}=\boldsymbol{\Sigma}^{-1}$ 的特征值才准（\cref{sec:est-crlb}）。 \\
“不相关就独立。” & 仅高斯下等价（\cref{thm:gauss-indep}）；一般只有“独立$\Rightarrow$不相关”。 \\
“条件化和边缘化都用同一种运算。” & 协方差形式：边缘易（取块）、条件难（舒尔补）；信息形式相反（\cref{tab:cov-info}）。 \\
“滤波比平滑省所以更优。” & 滤波边缘化历史使信息矩阵\emph{稠密}；平滑保持稀疏，常更准更可扩展（\cref{sec:est-batch-recursive,sec:est-smooth-filter}）。 \\
“卡尔曼滤波是另一种独立的方法。” & 它是批量 MAP 的\emph{前向递归特例}，线性高斯下无近似（\cref{thm:kf-from-map}）。 \\
“因子图节点包含测量。” & 节点只有待估变量，测量藏在因子里（\cref{def:factor-graph}）。 \\
“变量排序会改变结果。” & 任意排序给出\emph{相同}解，只改变 fill-in/速度（COLAMD，\cref{sec:est-fg-linear}）。 \\
“消息传递（GBP）只是另一种解法，结果和消元一样。” & 无环图上等价；有环图上 GBP 是\emph{近似}——收敛时均值精确但协方差偏乐观（\cref{sec:est-gbp}）。 \\
“连续时间就是在关键帧间插值，样条和 GP 没区别。” & 插值式相同，但 GP 比样条多一族\emph{运动先验因子}（自带正则、可算查询协方差）（\cref{sec:est-steam-spline}）。 \\
\bottomrule
\end{tabular}
\end{table}

% ════════════════════════════════════════════════════════════════
\section*{本章小结}
\addcontentsline{toc}{section}{本章小结}

\begin{table}[H]\centering\small
\caption{符号表}
\begin{tabular}{lll}
\toprule
符号 & 含义 & 首次出现 / 备注 \\
\midrule
$\mathbf{x},\mathbf{y},\mathbf{z},\mathbf{u}$ & 状态 / 路标 / 观测 / 控制 & \cref{sec:est-why}；批量时 $\mathbf{x}$ 含全部位姿与路标 \\
$\mathbf{w},\mathbf{v}$ & 过程 / 观测噪声 & $\sim\mathcal{N}(\mathbf{0},\boldsymbol{\Sigma}_{\mathbf{w}}),\mathcal{N}(\mathbf{0},\boldsymbol{\Sigma}_{\mathbf{v}})$ \\
$\boldsymbol{\Sigma},\boldsymbol{\Lambda}$ & 协方差 / 信息矩阵 & $\boldsymbol{\Lambda}=\boldsymbol{\Sigma}^{-1}$（\cref{eq:info-form}） \\
$\boldsymbol{\eta}$ & 信息向量 & $\boldsymbol{\eta}=\boldsymbol{\Sigma}^{-1}\boldsymbol{\mu}$ \\
$\check{\mathbf{x}},\hat{\mathbf{x}}$ & 先验/预测 / 后验/校正估计 & Barfoot check/hat 约定 \\
$\|\mathbf{e}\|_{\boldsymbol{\Sigma}}^2$ & 平方马氏距离 & $=\mathbf{e}^\top\boldsymbol{\Sigma}^{-1}\mathbf{e}$（\cref{def:gaussian}） \\
$H(\mathbf{x}),I(\mathbf{x},\mathbf{y})$ & 香农熵 / 互信息 & \cref{eq:gauss-entropy,eq:gauss-mutual-info}；主动感知用 \\
$\mathbf{V}$ & 协方差平方根 & $\boldsymbol{\Sigma}=\mathbf{V}\mathbf{V}^\top$，采样/白化（\cref{eq:gauss-sample}） \\
$\mathbf{S}$ & 舒尔补 & \cref{thm:schur-cond} 条件协方差 \\
$\mathbf{H},\mathbf{W},\mathbf{A}_i$ & 堆叠雅可比 / 权 / 白化雅可比 & 正规方程 \cref{eq:normal-eq}；$\mathbf{A}_i=\boldsymbol{\Sigma}_i^{-1/2}\mathbf{H}_i$ \\
$\mathcal{O}$ & 可观测性矩阵 & \cref{eq:observability} \\
$\boldsymbol{\mathcal{I}}_{\boldsymbol{\theta}}$ & Fisher 信息矩阵 & \cref{eq:crlb}；线性高斯下 $=\boldsymbol{\Lambda}$ \\
$\mathbf{K}_k$ & 卡尔曼增益 & \cref{thm:kf-from-map} \\
$\phi_i,\mathbf{s}_j$ & 因子 / 分隔集 & $\phi_i\propto\exp(-\tfrac12\|\mathbf{e}_i\|_{\boldsymbol{\Sigma}_i}^2)$；\cref{def:factor-graph,thm:elimination} \\
$\mathbf{R},\mathbf{L}$ & 平方根（Cholesky）因子 & $\boldsymbol{\Lambda}=\mathbf{R}^\top\mathbf{R}=\mathbf{L}\mathbf{L}^\top$（\cref{sec:est-fg-linear}） \\
$\mathbf{A}_{\mathcal{G}},\boldsymbol{\mathcal{L}},\boldsymbol{\mathcal{L}}_w$ & 图关联矩阵 / （加权）拉普拉斯 & $\boldsymbol{\mathcal{L}}=\mathbf{A}_{\mathcal{G}}^\top\mathbf{A}_{\mathcal{G}}$（\cref{eq:fim-laplacian-def}）；FIM $=\boldsymbol{\mathcal{L}}_w\otimes\mathbf{I}$ \\
$b(x_i),\,\boldsymbol{\eta}_{s\to i}$ & GBP 置信 / 因子到变量消息 & 信息形式相加$+$舒尔补（\cref{eq:gbp-belief,eq:gbp-factor-msg}） \\
\multicolumn{3}{l}{\emph{连续时间 / STEAM（\cref{sec:est-steam}）}} \\
$\psi_k(t),\psi_k^c(t),\mathbf{c}_k$ & 样条基 / 累积基 / 控制点 & 参数法（\cref{eq:steam-spline-basis,eq:steam-spline-cumul}） \\
$\boldsymbol{\varpi},\boldsymbol{\gamma}_k$ & 机体系广义速度 / 局部连续态 $[\boldsymbol{\xi}_k;\boldsymbol{\psi}_k]$ & \cref{eq:steam-local-sde,eq:steam-lti} \\
$\mathbf{Q}_C$ & 白噪声功率谱密度（PSD） & \cref{eq:steam-local-sde}；离散累积噪声记 $\mathbf{Q}(\Delta t)$ \\
$\boldsymbol{\Phi}(t,s),\mathbf{A}$ & GP 转移函数 / 连续系统矩阵 & \cref{eq:steam-Phi-Q}（WNOA 下 $\mathbf{A}^2=\mathbf{0}$） \\
$\boldsymbol{\Lambda}(\tau),\boldsymbol{\Psi}(\tau)$ & GP 插值权（左/右段） & \cref{eq:steam-interp,eq:steam-LambdaPsi} \\
$\check{\mathbf{P}},\hat{\mathbf{P}}$ & 轨迹先验/后验协方差（滤波惯例，$=\check{\boldsymbol{\Sigma}}/\hat{\boldsymbol{\Sigma}}$） & \cref{eq:steam-cov-bracket} \\
\bottomrule
\end{tabular}
\end{table}

\begin{table}[H]\centering\small
\caption{定理速查表（功能描述：这条让你能做什么）}
\begin{tabularx}{\linewidth}{lLl}
\toprule
名称 & 一句话功能 & 出处 \\
\midrule
贝叶斯 & 由先验与似然算后验 $\propto$ 似然 $\times$ 先验 & \cref{eq:bayes} \\
舒尔补条件化 & 联合高斯测一半、条件出另一半，协方差收缩 & \cref{thm:schur-cond} \\
信息形式边缘化 & 在信息矩阵上消掉一个变量（滑窗/边缘化） & \cref{eq:info-marginal} \\
独立$=$不相关（高斯） & 高斯下协方差为零即可判独立 & \cref{thm:gauss-indep} \\
高斯过非线性 & 线性化把高斯推过非线性，得 EKF 预测协方差 & \cref{prop:gauss-nonlinear} \\
信息相加 & 多条独立观测融合 $=$ 信息矩阵/向量相加 & \cref{prop:info-add} \\
SMW 引理 & 信息形式 $\leftrightarrow$ 协方差形式 $+$ 增益互转 & \cref{lem:smw} \\
CRLB & 给出无偏估计方差的理论下界 & \cref{thm:crlb} \\
互信息（高斯） & 量化一次观测能消去多少不确定度（主动感知） & \cref{eq:gauss-mutual-info} \\
高斯采样 & 由协方差平方根 $\mathbf{V}$ 从标准正态造出任意高斯 & \cref{eq:gauss-sample} \\
MAP $=$ 加权最小二乘 & 高斯下 $\arg\max p=\arg\min\sum\|\cdot\|^2$ & \cref{thm:map-ls} \\
正规方程 & 一步解线性高斯 MAP，$\hat{\boldsymbol{\Sigma}}=\boldsymbol{\Lambda}^{-1}=$ CRLB & \cref{thm:normal-eq} \\
MAP $=$ 贝叶斯均值 & 线性高斯下三条路径同解（众数$=$均值） & \cref{thm:map-bayes} \\
可观测性 & 判定解是否存在唯一（$\boldsymbol{\Lambda}$ 可逆） & \cref{eq:observability} \\
KF $=$ 批量 MAP 递归 & 把批量解递归化、边缘化历史得卡尔曼滤波 & \cref{thm:kf-from-map} \\
因子图 MAP & $\arg\max\prod\phi_i=\arg\min\sum\|\cdot\|^2$ & \cref{thm:fg-map} \\
消元 $=$ 分解 & 因子图变量消元就是稀疏 Cholesky/QR & \cref{thm:elimination} \\
GBP（消息传递） & 纯局部/分布式近似推断；有环图若收敛则均值精确、协方差偏乐观 & \cref{eq:gbp-factor-msg} \\
FIM $=$ 图拉普拉斯 & 已知朝向位姿图的 Fisher 信息 $=\boldsymbol{\mathcal{L}}_w\otimes\mathbf{I}$，精度由图连通度定 & \cref{thm:fim-laplacian} \\
STEAM 不破稀疏 & 连续时间 GP 先验下信息阵仍块三对角箭头形，求解照旧高效 & \cref{prop:steam-sparse} \\
样条 vs GP & 两者都是控制点间插值；GP $=$ 样条 $+$ 运动先验因子（自带正则与协方差） & \cref{sec:est-steam-spline} \\
GP 插值/外推查询 & 解一次批量后以 $O(1)$ 查询任意时刻位姿/速度\emph{及协方差} & \cref{eq:steam-interp,eq:steam-extrap} \\
\bottomrule
\end{tabularx}
\end{table}

\begin{table}[H]\centering\small
\caption{知识点总表（本章 $\to$ 后续）}
\begin{tabular}{p{0.27\linewidth}llp{0.27\linewidth}}
\toprule
知识点 & 难度 & 脉络层 & 在哪里展开 \\
\midrule
概率/高斯/信息形式/Schur/SMW & $\bigstar\bigstar$ & 骨干 & \cref{ch:eskf}（滤波）反复用 \\
MLE/MAP/最小二乘 & $\bigstar$ & 骨干 & 所有后端章 \\
正规方程/信息矩阵/CRLB & $\bigstar\bigstar$ & 骨干 & \cref{ch:nlopt} 迭代解之 \\
可观测性/gauge & $\bigstar\bigstar\bigstar$ & 次优 & 标定、VIO 初始化 \\
批量 vs 递归（等价/稀疏） & $\bigstar\bigstar\bigstar$ & 次优 & \cref{ch:eskf} 完整 KF/RTS \\
因子图/平滑 vs 滤波 & $\bigstar\bigstar$ & 骨干 & \cref{ch:vo,ch:vio,ch:lidar_slam} 后端 \\
变量消元/Bayes 树/iSAM & $\bigstar\bigstar\bigstar\bigstar$ & 前沿 & \cref{ch:vio} 增量后端 \\
高斯置信传播（GBP） & $\bigstar\bigstar\bigstar$ & 进阶（选读） & \cref{sec:est-gbp}；分布式/图处理器 \\
FIM $=$ 图拉普拉斯/最优设计 & $\bigstar\bigstar\bigstar$ & 进阶 & \cref{sec:est-fim-laplacian}；主动 SLAM/剪枝 \\
连续时间/GP 回归（STEAM） & $\bigstar\bigstar\bigstar\bigstar$ & 进阶 & \cref{sec:est-steam} \\
参数样条 vs GP & $\bigstar\bigstar\bigstar$ & 进阶 & \cref{sec:est-steam-spline} \\
GN/LM/稀疏求解 & $\bigstar\bigstar\bigstar$ & 次优（推迟） & \cref{ch:nlopt} \\
\bottomrule
\end{tabular}
\end{table}

\paragraph{延伸阅读。}
\begin{itemize}
  \item Barfoot《State Estimation for Robotics》\cite{barfoot2024state} 第 1--3 章：概率论基元、线性高斯批量/递归估计、可观测性与 CRLB 的严谨处理（本章理论骨架与几乎全部推导的来源）。
  \item SLAM Handbook\cite{carlone2026handbook} 第 1 章 \& Dellaert/Kaess 因子图综述\cite{dellaert2017factor}：因子图、MAP、线性化/白化、稀疏 Cholesky/QR（$\sqrt{}$SAM）、变量消元、平滑 vs 滤波、Bayes 树/iSAM2（本章因子图视角的来源）。
  \item SLAM Handbook\cite{carlone2026handbook} 第 2 章：连续时间状态表示——参数样条（B 样条/累积样条/李群累积样条）与高斯过程两条路的并列处理（\cref{sec:est-steam-spline} 样条一支的来源）。
  \item SLAM Handbook\cite{carlone2026handbook} 第 6 章 §6.2：Fisher 信息阵与位姿图拉普拉斯的对应、D/A/E-最优设计与生成树（\cref{sec:est-fim-laplacian} 的来源）；其 §6.1 的可证最优求解器（半定松弛/SE-Sync）见 \cref{ch:nlopt}。
  \item GBP 教程\cite{ortiz2021gbp} 与 Davison–Ortiz\cite{davison2019futuremapping2}：高斯置信传播的消息推导、收敛性与面向空间智能的硬件视角（\cref{sec:est-gbp} 的来源）。
  \item 视觉 SLAM 十四讲\cite{gaoxiang2019slam14} 第 6 章 §6.1：从概率到最小二乘的直觉叙事与一维小车例。
\end{itemize}

\paragraph{与后续章节的关系。}
本章把 SLAM 钉成“因子图上的 MAP $=$ 最小二乘”，并证明它如何被批量稀疏分解或递归滤波解出。两条主线由此展开：\emph{递归}地解它 $\to$ 卡尔曼/误差状态滤波（\cref{ch:eskf}，本章 \cref{thm:kf-from-map} 已证其为批量 MAP 的前向特例）；\emph{迭代}地解非线性版 $\to$ 高斯牛顿/LM 与稀疏求解、Schur 消元（\cref{ch:nlopt}，本章正规方程 \cref{eq:normal-eq} 是其每步内核）。其上层即视觉里程计（\cref{ch:vo}）、激光 SLAM（\cref{ch:lidar_slam}）、VIO（\cref{ch:vio}）的后端。

\begin{table}[H]\centering\small
\caption{本章与后续章节的关系}
\begin{tabular}{lll}
\toprule
后续章节 & 与本章的关系 & 本章哪个知识点为其铺垫 \\
\midrule
\cref{ch:eskf} 卡尔曼/ESKF & 本章 MAP 的递归特例 & \cref{thm:kf-from-map}、信息/协方差形式 \\
\cref{ch:nlopt} 非线性优化 & 解本章正规方程的迭代法 & \cref{eq:normal-eq}、稀疏/Schur/白化 \\
\cref{ch:vo} 视觉里程计 & 在因子图上做 MAP & \cref{thm:fg-map}、重投影因子 \\
\cref{ch:vio} VIO & 增量因子图后端 & 因子图、Bayes 树/iSAM、边缘化 \\
\cref{ch:lidar_slam} 激光 SLAM & 位姿图/PGO 是因子图特例 & \cref{tab:est-slam-variants}、回环因子 \\
\bottomrule
\end{tabular}
\end{table}

\begin{table}[H]\centering\small
\caption{故障排查}
\begin{tabular}{p{0.28\linewidth}p{0.32\linewidth}p{0.32\linewidth}}
\toprule
症状 & 可能原因 & 对策 \\
\midrule
解不唯一 / 数值漂移 & $\boldsymbol{\Lambda}$ 奇异（gauge/不可观测） & 加先验因子锚定；检查 $\mathbf{H}$ 零空间 / $\mathcal{O}$ 满秩（\cref{sec:est-observability}） \\
优化发散 / 陷局部极小 & 初值差、$h$ 强非线性 & 好初值；鲁棒核；见 \cref{ch:nlopt} \\
协方差越来越稠密、变慢 & 误用滤波边缘化历史 & 改平滑/因子图保稀疏（\cref{sec:est-smooth-filter}） \\
某方向误差异常大 & 该方向信息弱（$\boldsymbol{\Lambda}$ 小特征值） & 加观测/约束；查可观测性 \\
权重失衡、某项主导 & $\boldsymbol{\Sigma}_i$ 设定不当 & 按传感器真实噪声标定 $\boldsymbol{\Sigma}_i$；卡方一致性检验（\cref{prop:chi2}） \\
求解很慢但结果对 & 变量排序差、fill-in 多 & 用 COLAMD 重排（\cref{sec:est-fg-linear}） \\
\bottomrule
\end{tabular}
\end{table}

\paragraph{研究实践建议。}
给新手：先用 GTSAM/g2o 把一个小因子图跑通，对照本章的因子表（\cref{tab:factors}）逐项确认残差与噪声模型，再回头读正规方程——“先会用、再懂内核”比反过来快。给有经验者：在自己的后端里显式监控信息矩阵的稀疏度与 fill-in，用卡方一致性检验（\cref{prop:chi2}）核查噪声协方差是否标定正确；做滑窗时务必同步更新边缘化后的信息向量（\cref{eq:info-marginal}），这是最常见的隐蔽 bug。

% ════════════════════════════════════════════════════════════════
\section*{附录：本章推导补遗}
\label{sec:est-appendix}
\addcontentsline{toc}{section}{附录：本章推导补遗}

\begin{derivation}[舒尔补条件化 \cref{eq:schur-cond} 的二次型配方细节]
对联合协方差作块 $\mathbf{LDU}$ 分解：
\[
  \begin{bmatrix}\boldsymbol{\Sigma}_{xx}&\boldsymbol{\Sigma}_{xy}\\\boldsymbol{\Sigma}_{yx}&\boldsymbol{\Sigma}_{yy}\end{bmatrix}
  =\begin{bmatrix}\mathbf{I}&\boldsymbol{\Sigma}_{xy}\boldsymbol{\Sigma}_{yy}^{-1}\\\mathbf{0}&\mathbf{I}\end{bmatrix}
   \begin{bmatrix}\mathbf{S}&\mathbf{0}\\\mathbf{0}&\boldsymbol{\Sigma}_{yy}\end{bmatrix}
   \begin{bmatrix}\mathbf{I}&\mathbf{0}\\\boldsymbol{\Sigma}_{yy}^{-1}\boldsymbol{\Sigma}_{yx}&\mathbf{I}\end{bmatrix},
  \quad \mathbf{S}=\boldsymbol{\Sigma}_{xx}-\boldsymbol{\Sigma}_{xy}\boldsymbol{\Sigma}_{yy}^{-1}\boldsymbol{\Sigma}_{yx}.
\]
求逆得联合\emph{信息}矩阵的相应分块，代入高斯指数 $-\tfrac12[\,\cdot\,]^\top\boldsymbol{\Lambda}[\,\cdot\,]$ 并按 $\mathbf{x}$ 配方，交叉项把均值平移到 $\boldsymbol{\mu}_x+\boldsymbol{\Sigma}_{xy}\boldsymbol{\Sigma}_{yy}^{-1}(\mathbf{y}-\boldsymbol{\mu}_y)$，剩余二次型的系数矩阵为 $\mathbf{S}^{-1}$，即条件协方差 $=\mathbf{S}$。于是 \cref{eq:schur-cond} 成立。$\mathbf{S}$ 即 $\boldsymbol{\Sigma}_{yy}$ 的舒尔补。\hfill$\square$
\end{derivation}

\begin{derivation}[高斯负对数 $\to$ 马氏二次型 \cref{eq:neglog,eq:single-ls}]
由 \cref{eq:gaussian}，$p(\mathbf{z}\mid\mathbf{x})=\big((2\pi)^M\det\boldsymbol{\Sigma}_{\mathbf{v}}\big)^{-1/2}\exp(-\tfrac12\|\mathbf{z}-h(\mathbf{x})\|_{\boldsymbol{\Sigma}_{\mathbf{v}}}^2)$。取 $-\ln(\cdot)$：归一化常数项 $\tfrac12\ln((2\pi)^M\det\boldsymbol{\Sigma}_{\mathbf{v}})$ 不含 $\mathbf{x}$；剩 $\tfrac12\|\mathbf{z}-h(\mathbf{x})\|_{\boldsymbol{\Sigma}_{\mathbf{v}}}^2$。$\arg\max p=\arg\min(-\ln p)=\arg\min\|\mathbf{z}-h(\mathbf{x})\|_{\boldsymbol{\Sigma}_{\mathbf{v}}}^2$（正系数 $\tfrac12$ 不影响 $\arg\min$）。批量情形对 \cref{eq:factorize} 的乘积取负对数即化乘为和，得 \cref{eq:batch-ls}。\hfill$\square$
\end{derivation}

\begin{derivation}[信息形式下的观测融合 \cref{prop:info-add}]
两独立高斯 $\mathcal{N}(\boldsymbol{\mu}_1,\boldsymbol{\Sigma}_1),\mathcal{N}(\boldsymbol{\mu}_2,\boldsymbol{\Sigma}_2)$ 之积的指数为 $-\tfrac12\big[(\mathbf{x}-\boldsymbol{\mu}_1)^\top\boldsymbol{\Sigma}_1^{-1}(\mathbf{x}-\boldsymbol{\mu}_1)+(\mathbf{x}-\boldsymbol{\mu}_2)^\top\boldsymbol{\Sigma}_2^{-1}(\mathbf{x}-\boldsymbol{\mu}_2)\big]$。按 $\mathbf{x}$ 收集二次项系数得 $\boldsymbol{\Lambda}=\boldsymbol{\Sigma}_1^{-1}+\boldsymbol{\Sigma}_2^{-1}$，一次项系数得 $\boldsymbol{\eta}=\boldsymbol{\Sigma}_1^{-1}\boldsymbol{\mu}_1+\boldsymbol{\Sigma}_2^{-1}\boldsymbol{\mu}_2$，配方后仍是高斯，均值 $\boldsymbol{\Lambda}^{-1}\boldsymbol{\eta}$。归纳到 $K$ 项即 \cref{eq:info-sum}；带映射 $\mathbf{G}_k$ 时二次项系数为 $\sum_k\mathbf{G}_k^\top\boldsymbol{\Sigma}_k^{-1}\mathbf{G}_k$，得 \cref{eq:info-sum-G}。这正是正规方程 \cref{eq:normal-eq} 两端皆为累加的来由。\hfill$\square$
\end{derivation}

\begin{derivation}[SMW 引理 \cref{eq:smw} 的分块求逆证明]
把 $\mathbf{M}=\big[\begin{smallmatrix}\mathbf{A}^{-1}&-\mathbf{B}\\\mathbf{C}&\mathbf{D}\end{smallmatrix}\big]$ 分别作 LDU 与 UDL 分解：
\[
  \mathbf{M}=\begin{bmatrix}\mathbf{I}&\mathbf{0}\\\mathbf{C}\mathbf{A}&\mathbf{I}\end{bmatrix}
  \begin{bmatrix}\mathbf{A}^{-1}&\mathbf{0}\\\mathbf{0}&\mathbf{D}+\mathbf{C}\mathbf{A}\mathbf{B}\end{bmatrix}
  \begin{bmatrix}\mathbf{I}&-\mathbf{A}\mathbf{B}\\\mathbf{0}&\mathbf{I}\end{bmatrix}
  =\begin{bmatrix}\mathbf{I}&-\mathbf{B}\mathbf{D}^{-1}\\\mathbf{0}&\mathbf{I}\end{bmatrix}
  \begin{bmatrix}\mathbf{A}^{-1}+\mathbf{B}\mathbf{D}^{-1}\mathbf{C}&\mathbf{0}\\\mathbf{0}&\mathbf{D}\end{bmatrix}
  \begin{bmatrix}\mathbf{I}&\mathbf{0}\\\mathbf{D}^{-1}\mathbf{C}&\mathbf{I}\end{bmatrix}.
\]
两式各自求逆。LDU 逆的左上块为 $\mathbf{A}-\mathbf{A}\mathbf{B}(\mathbf{D}+\mathbf{C}\mathbf{A}\mathbf{B})^{-1}\mathbf{C}\mathbf{A}$，UDL 逆的左上块为 $(\mathbf{A}^{-1}+\mathbf{B}\mathbf{D}^{-1}\mathbf{C})^{-1}$；二者相等即 \cref{eq:smw}。比较右下块得 $(\mathbf{D}+\mathbf{C}\mathbf{A}\mathbf{B})^{-1}=\mathbf{D}^{-1}-\mathbf{D}^{-1}\mathbf{C}(\mathbf{A}^{-1}+\mathbf{B}\mathbf{D}^{-1}\mathbf{C})^{-1}\mathbf{B}\mathbf{D}^{-1}$，比较非对角块得另两式。取 $\mathbf{A}=\boldsymbol{\Sigma}$、$\mathbf{B}=\mathbf{C}^\top$、$\mathbf{C}=\mathbf{C}$、$\mathbf{D}=\mathbf{R}$ 即得卡尔曼增益两写法的等价。\hfill$\square$
\end{derivation}

\begin{derivation}[KF 预测步协方差 \cref{eq:kf-info-form} 的期望推导]
预测均值：$\check{\mathbf{x}}_k=\mathbb{E}[\mathbf{A}_{k-1}\mathbf{x}_{k-1}+\mathbf{v}_k+\mathbf{w}_k]=\mathbf{A}_{k-1}\hat{\mathbf{x}}_{k-1}+\mathbf{v}_k$（$\mathbb{E}[\mathbf{w}_k]=\mathbf{0}$）。预测协方差：$\check{\boldsymbol{\Sigma}}_k=\mathbb{E}[(\mathbf{x}_k-\check{\mathbf{x}}_k)(\cdots)^\top]=\mathbf{A}_{k-1}\underbrace{\mathbb{E}[(\mathbf{x}_{k-1}-\hat{\mathbf{x}}_{k-1})(\cdots)^\top]}_{\hat{\boldsymbol{\Sigma}}_{k-1}}\mathbf{A}_{k-1}^\top+\underbrace{\mathbb{E}[\mathbf{w}_k\mathbf{w}_k^\top]}_{\boldsymbol{\Sigma}_{\mathbf{w},k}}=\mathbf{A}_{k-1}\hat{\boldsymbol{\Sigma}}_{k-1}\mathbf{A}_{k-1}^\top+\boldsymbol{\Sigma}_{\mathbf{w},k}$（用 $\mathbf{x}_{k-1}$ 与 $\mathbf{w}_k$ 不相关）。校正步把 $(\check{\mathbf{x}}_k,\check{\boldsymbol{\Sigma}}_k)$ 与观测 $\mathbf{y}_k=\mathbf{C}_k\mathbf{x}_k+\mathbf{n}_k$ 作联合高斯并条件化（\cref{eq:schur-cond}），即得 \cref{eq:kf-info-form}；完整递推与各等价形式见 \cref{ch:eskf}。\hfill$\square$
\end{derivation}

\begin{derivation}[高斯过非线性 \cref{eq:gauss-nonlinear} 的逐行矩阵代数（配方 $\to$ 积分 $\to$ SMW）]\label{der:gauss-nonlinear-full}
把线性化 $g(\mathbf{x})\approx\boldsymbol{\mu}_y+\mathbf{G}(\mathbf{x}-\boldsymbol{\mu}_x)$（$\boldsymbol{\mu}_y=g(\boldsymbol{\mu}_x)$，$\mathbf{G}=\partial g/\partial\mathbf{x}|_{\boldsymbol{\mu}_x}$）代入 $p(\mathbf{y})=\int p(\mathbf{y}\mid\mathbf{x})p(\mathbf{x})\,\mathrm{d}\mathbf{x}$，合并两个高斯指数并按 $\mathbf{x}$ 整理（$\boldsymbol{\Sigma}_{\mathbf{v}}$ 为观测噪声协方差）：
\[
  p(\mathbf{y})=\eta\,\mathrm{e}^{-\frac12(\mathbf{y}-\boldsymbol{\mu}_y)^\top\boldsymbol{\Sigma}_{\mathbf{v}}^{-1}(\mathbf{y}-\boldsymbol{\mu}_y)}
  \!\int\!\mathrm{e}^{-\frac12(\mathbf{x}-\boldsymbol{\mu}_x)^\top(\boldsymbol{\Sigma}_{xx}^{-1}+\mathbf{G}^\top\boldsymbol{\Sigma}_{\mathbf{v}}^{-1}\mathbf{G})(\mathbf{x}-\boldsymbol{\mu}_x)}
  \,\mathrm{e}^{(\mathbf{y}-\boldsymbol{\mu}_y)^\top\boldsymbol{\Sigma}_{\mathbf{v}}^{-1}\mathbf{G}(\mathbf{x}-\boldsymbol{\mu}_x)}\,\mathrm{d}\mathbf{x}.
\]
\textbf{(1) 配方}：定义 $\mathbf{F}$ 满足 $\mathbf{F}^\top(\mathbf{G}^\top\boldsymbol{\Sigma}_{\mathbf{v}}^{-1}\mathbf{G}+\boldsymbol{\Sigma}_{xx}^{-1})=\boldsymbol{\Sigma}_{\mathbf{v}}^{-1}\mathbf{G}$，则积分内两因子凑成
\[
  \mathrm{e}^{-\frac12((\mathbf{x}-\boldsymbol{\mu}_x)-\mathbf{F}(\mathbf{y}-\boldsymbol{\mu}_y))^\top(\mathbf{G}^\top\boldsymbol{\Sigma}_{\mathbf{v}}^{-1}\mathbf{G}+\boldsymbol{\Sigma}_{xx}^{-1})((\mathbf{x}-\boldsymbol{\mu}_x)-\mathbf{F}(\mathbf{y}-\boldsymbol{\mu}_y))}\,
  \mathrm{e}^{\frac12(\mathbf{y}-\boldsymbol{\mu}_y)^\top\mathbf{F}^\top(\mathbf{G}^\top\boldsymbol{\Sigma}_{\mathbf{v}}^{-1}\mathbf{G}+\boldsymbol{\Sigma}_{xx}^{-1})\mathbf{F}(\mathbf{y}-\boldsymbol{\mu}_y)}.
\]
\textbf{(2) 积分吸常数}：第二个因子不含 $\mathbf{x}$，提到积分号外；剩余对 $\mathbf{x}$ 是纯高斯，积分得与 $\mathbf{y}$ 无关的常数，并入归一化 $\eta\to\rho$。于是 $p(\mathbf{y})=\rho\,\mathrm{e}^{-\frac12(\mathbf{y}-\boldsymbol{\mu}_y)^\top\mathbf{M}(\mathbf{y}-\boldsymbol{\mu}_y)}$，其中 $\mathbf{M}=\boldsymbol{\Sigma}_{\mathbf{v}}^{-1}-\mathbf{F}^\top(\mathbf{G}^\top\boldsymbol{\Sigma}_{\mathbf{v}}^{-1}\mathbf{G}+\boldsymbol{\Sigma}_{xx}^{-1})\mathbf{F}$。\textbf{(3) 代入 $\mathbf{F}$ 并用 SMW}：由定义式转置解出 $\mathbf{F}=(\mathbf{G}^\top\boldsymbol{\Sigma}_{\mathbf{v}}^{-1}\mathbf{G}+\boldsymbol{\Sigma}_{xx}^{-1})^{-1}\mathbf{G}^\top\boldsymbol{\Sigma}_{\mathbf{v}}^{-1}$，代回
\[
  \mathbf{M}=\boldsymbol{\Sigma}_{\mathbf{v}}^{-1}-\boldsymbol{\Sigma}_{\mathbf{v}}^{-1}\mathbf{G}(\mathbf{G}^\top\boldsymbol{\Sigma}_{\mathbf{v}}^{-1}\mathbf{G}+\boldsymbol{\Sigma}_{xx}^{-1})^{-1}\mathbf{G}^\top\boldsymbol{\Sigma}_{\mathbf{v}}^{-1}
  =(\boldsymbol{\Sigma}_{\mathbf{v}}+\mathbf{G}\boldsymbol{\Sigma}_{xx}\mathbf{G}^\top)^{-1},
\]
末步取 SMW \cref{eq:smw} 中 $\mathbf{A}^{-1}=\boldsymbol{\Sigma}_{\mathbf{v}}$、$\mathbf{B}=\mathbf{G}$、$\mathbf{D}=\boldsymbol{\Sigma}_{xx}^{-1}$、$\mathbf{C}=\mathbf{G}^\top$。故 $\mathbf{y}\sim\mathcal{N}\big(g(\boldsymbol{\mu}_x),\,\boldsymbol{\Sigma}_{\mathbf{v}}+\mathbf{G}\boldsymbol{\Sigma}_{xx}\mathbf{G}^\top\big)$，即 \cref{eq:gauss-nonlinear}。\hfill$\square$
\end{derivation}
